% concatenated from example_paper.tex
%%%%%%%% ICML 2025 EXAMPLE LATEX SUBMISSION FILE %%%%%%%%%%%%%%%%%

\documentclass{article}

% Recommended, but optional, packages for figures and better typesetting:
\usepackage{microtype}
\usepackage{graphicx}
\usepackage{subfigure}
\usepackage{diagbox}
\usepackage{url}
\usepackage{booktabs} % for professional tables
\usepackage{natbib}
% hyperref makes hyperlinks in the resulting PDF.
% If your build breaks (sometimes temporarily if a hyperlink spans a page)
% please comment out the following usepackage line and replace
% \usepackage{icml2025} with \usepackage[nohyperref]{icml2025} above.
\usepackage{hyperref}

\usepackage[a4paper,margin=3cm]{geometry}
% Attempt to make hyperref and algorithmic work together better:
\newcommand{\theHalgorithm}{\arabic{algorithm}}

% Use the following line for the initial blind version submitted for review:
% \usepackage{icml2025}

% If accepted, instead use the following line for the camera-ready submission:
% \usepackage[accepted]{icml2025}

% For theorems and such
\usepackage{amsmath}
\usepackage{amssymb}
\usepackage{mathtools}
\usepackage{amsthm}
\usepackage{circuitikz}


% Personal command

\newcommand{\jh}[1]{{\color{red} #1}}
\newcommand{\jfa}[1]{{\color{green} #1}}
\newcommand{\ar}[1]{{\color{purple} #1}}
\newcommand{\cd}[1]{{\color{blue} #1}}
\newcommand{\lh}[1]{{\color{orange} #1}}
%sets
\def\R{{\mathbb{R}}}
\def\C{{\mathbb{C}}}
\def\N{{\mathbb{N}}}
\def\cK{\mathbf{K}}
\def\cL{\mathbf{L}}
\def\sS{\mathbf{S}}
\def\sphere{\mathbb{S}}
\def\1{\mathbb{1}}
\def\Abar{\overline{\A}}
\def\xcont{(x_t)_{t\ge 0}}
\def\xz{(x_t,z_t)_{t\ge 0}}
\def\xdisc{\{\tilde x_k \}_{k\ge 0}}
\def\xgen{\{x_t \}_{t\in \mathcal{K}}}
\def\mart{(M_t)_{t\ge 0}}
\def\xdisctilde{\{\tilde{x}_k \}_{k\in \N}}
\def\pl{{\text{PL}^\mu}}
\def\rsi{{\text{RSI}^\nu}}
%Operators
\newcommand{\A}{\mathcal{A}}
\renewcommand{\P}{\mathbb{P}}
\newcommand{\bpar}[1]{\left(#1\right)}
\newcommand{\prox}[2]{\text{prox}_{#2g}\left(#1\right)}
\newcommand{\Lv}{\mathcal{L}_v}

\newcommand{\E}[1]{\mathbb{E}\left[#1 \right]}
\newcommand{\condE}[1]{\mathbb{E}_{k}\left[#1 \right]}
\newcommand{\norm}[1]{\left\lVert#1\right\rVert}
\newcommand{\dotprod}[1]{\left< #1\right>}
\newcommand{\grad}[1]{\nabla f(#1)}
\newcommand{\gradsto}[1]{\nabla f_{v_n} (#1)}
\newcommand{\sqc}{\text{SQC}_{\tau}^{\mu}}
\newcommand{\qc}{\text{QC}_{\tau}}
\newcommand{\ls}{\text{LS}^{L}}
\newcommand{\lsf}[1]{\text{LS}^{#1}}
\newcommand{\strongconv}{\text{SC}^{\mu}}
\newcommand{\strongconvf}[1]{\text{SC}^{#1}}
\newcommand{\ac}{\text{AC}^{a}}
\newcommand{\acf}[1]{\text{AC}^{#1}}
\newcommand{\qg}{\text{QG}_{+}^{L_0}}
\newcommand{\qginf}{\text{QG}_{-}^{\mu_0}}
\newcommand{\qgf}[1]{\text{QG}_{+}^{#1}}
\newcommand{\qginff}[1]{\text{QG}_{-}^{#1}}
\newcommand{\sqcf}[2]{\text{SQC}_{#1}^{#2}}
\newcommand{\plf}[1]{\text{PL}^{#1}}
\newcommand{\rsif}[1]{\text{RSI}^{#1}}
\newcommand{\gradstobatch}[2]{\tilde \nabla_{#2} (#1)}
\newcommand{\footremember}[2]{%
    \footnote{#2}
    \newcounter{#1}
    \setcounter{#1}{\value{footnote}}%
}
\newcommand{\footrecall}[1]{%
    \footnotemark[\value{#1}]%
} 
\newenvironment{sproof}{%
  \renewcommand{\proofname}{Sketch of Proof}\proof}{\endproof}
\newcommand{\off}[1]{}
\newcommand{\Prb}[1]{\mathbb{P}\left[#1\right] }
\newcommand{\ind}{\perp\!\!\!\!\perp}
\newcommand{\bigO}{\mathcal{O}}
\newcommand{\eqdef}{\overset{\text{def}}{=}}
\def\restriction#1#2{\mathchoice
              {\setbox1\hbox{${\displaystyle #1}_{\scriptstyle #2}$}
              \restrictionaux{#1}{#2}}
              {\setbox1\hbox{${\textstyle #1}_{\scriptstyle #2}$}
              \restrictionaux{#1}{#2}}
              {\setbox1\hbox{${\scriptstyle #1}_{\scriptscriptstyle #2}$}
              \restrictionaux{#1}{#2}}
              {\setbox1\hbox{${\scriptscriptstyle #1}_{\scriptscriptstyle #2}$}
              \restrictionaux{#1}{#2}}}
\def\restrictionaux#1#2{{#1\,\smash{\vrule height .8\ht1 depth .85\dp1}}_{\,#2}} 
\def\argmin{\textup{argmin}\,}
\def\arginf{\textup{arginf}\,}



% if you use cleveref..
\usepackage[capitalize,noabbrev]{cleveref}

%%%%%%%%%%%%%%%%%%%%%%%%%%%%%%%%
% THEOREMS
%%%%%%%%%%%%%%%%%%%%%%%%%%%%%%%%
\theoremstyle{plain}
\newtheorem{theorem}{Theorem}[section]
\newtheorem{proposition}[theorem]{Proposition}
\newtheorem{lemma}[theorem]{Lemma}
\newtheorem{corollary}[theorem]{Corollary}
\theoremstyle{definition}
\newtheorem{definition}[theorem]{Definition}
\newtheorem{assumption}[theorem]{Assumption}
\theoremstyle{remark}
\newtheorem{remark}[theorem]{Remark}
\newtheorem*{remark*}{Remark}

% Todonotes is useful during development; simply uncomment the next line
%    and comment out the line below the next line to turn off comments
%\usepackage[disable,textsize=tiny]{todonotes}
\usepackage[textsize=tiny]{todonotes}
\usepackage{enumitem}

% The \icmltitle you define below is probably too long as a header.
% Therefore, a short form for the running title is supplied here:
% \icmltitlerunning{Nesterov Acceleration with Aiming Condition}\title{Your Title Here}
\author{Julien Hermant\\
Univ. Bordeaux, IMB}
\date{}

\title{Acceleration for Polyak–\L{}ojasiewicz Functions with a Gradient Aiming Condition}
\author{Julien Hermant\\
Univ. Bordeaux, IMB}
\date{}

\begin{document}

% \twocolumn[
% \icmltitle{Acceleration for Polyak–\L{}ojasiewicz Functions with a Gradient Aiming Condition}

% It is OKAY to include author information, even for blind
% submissions: the style file will automatically remove it for you
% unless you've provided the [accepted] option to the icml2025
% package.

% List of affiliations: The first argument should be a (short)
% identifier you will use later to specify author affiliations
% Academic affiliations should list Department, University, City, Region, Country
% Industry affiliations should list Company, City, Region, Country

% You can specify symbols, otherwise they are numbered in order.
% Ideally, you should not use this facility. Affiliations will be numbered
% in order of appearance and this is the preferred way.
% \icmlsetsymbol{equal}{*}

\maketitle

\vskip 0.3in


% this must go after the closing bracket ] following \twocolumn[ ...

% This command actually creates the footnote in the first column
% listing the affiliations and the copyright notice.
% The command takes one argument, which is text to display at the start of the footnote.
% The \icmlEqualContribution command is standard text for equal contribution.
% Remove it (just {}) if you do not need this facility.

%\printAffiliationsAndNotice{}  % leave blank if no need to mention equal contribution
% \printAffiliationsAndNotice{\icmlEqualContribution} % otherwise use the standard text.

\begin{abstract}
% It is known that smooth Polyak-\L{}ojasiewicz (PL) functions are such that minimization with momentum algorithms cannot significantly improve the convergence bound of gradient descent, contrasting with the acceleration phenomenon occurring in the strongly convex case. According to the literature, the intermediate class of non-convex functions called \textit{strongly quasar-convex} seems to fill this gap.  In this work we show that momentum actually does not necessarily yields better theoretical bounds for these functions. 
% %are flaws in our current understanding of acceleration with these functions. 
% We propose an alternative perspective, identifying that a key ingredient to accelerate with momentum for PL functions is an \textit{aiming condition}, which quantifies how well the descent direction aims toward a minimizer.
It is known that when minimizing smooth Polyak-\L{}ojasiewicz (PL) functions, Nesterov's momentum algorithm cannot significantly improve the convergence bound of gradient descent, contrasting with the acceleration phenomenon occurring in the strongly convex case. To bridge this gap, the literature has proposed \textit{strongly quasar-convex} functions as an intermediate non-convex class, for which the accelerated bounds of the strongly convex case has been suggested to persist. We show that strong quasar-convexity actually does not ensure a similar acceleration phenomenon. As an alternative, we study PL functions under an \textit{aiming condition} that measures how well the descent direction points toward a minimizer. This viewpoint sheds new light on the possibility of theoretical acceleration with Nesterov's momentum when minimizing PL functions.
% We propose an alternative perspective based on an \textit{aiming condition} that measures how well the descent direction points toward a minimizer. This perspective clarifies the geometric ingredient enabling provable benefit of momentum when minimizing PL functions.
\end{abstract}

\section{Introduction}

%Momentum \cite{polyak1964some,nesterov1983method} is among the most important mechanisms of modern machine learning solvers, because of its empirical effectiveness \cite{sutskever2013importance,he2016deep}.
Momentum \cite{polyak1964some,nesterov1983method} is an important mechanism of several modern optimizers. In machine learning, it is observed to substantially improve the optimization of large-scale models in practice \cite{sutskever2013importance,he2016deep}.
Although the theoretical benefit of incorporating momentum in vanilla gradient descent is well established in the convex and strongly convex setting \cite{nemirovskij1983problem,nesterovbook}, non-convexity is, however, an inherent property of many problems \cite{dauphin2014identifying,ge2016matrix,ge2017no,bhojanapalli2016global,li2018visualizing}. Importantly, for some non-convex settings it is not possible to demonstrate acceleration using momentum \cite{lowerboundI,PLlowerbound}. 
Focusing on neural networks, some works establish acceleration properties \cite{wang2021modular, liu2022provable,liu2022convergence, liao2024provable}, but they rely on simplifying assumptions--such as a very high degree of over-parameterization--that makes the learning problem similar to a strongly convex quadratic minimisation problem.
%, under which the objective is similar to a quadratic function.

%but their over-parameterization is so important that they behave like quadratic functions.
Another line of work lies in considering classes of non-convex functions that relax strong convexity.
In particular, the class of Polyak-\L{}ojasiewicz (PL) functions \cite{Polyak1963,lojasiewicz1963topological} has emerged as an important framework and has stimulated extensive interest \cite{karimi2016linear,fazel2018global,oymak2019overparameterized,altschuler2021averaging,merigot2021non,apidopoulos2022convergence,apidopoulos2025heavy,gess2023exponential}. A continuously differentiable function $f:\R^d \to \R $ belongs to this class if there exists $\mu>0$ such that the following inequality holds for all $x\in\R^d$
\[ \norm{\nabla f(x)}^2 \ge 2\mu\Big(f(x)-\min_{u\in \R^d}f(u)\Big).\]
%has asserted itself as a key class and has stimulated extensive interest
%When it comes to showing theoretically an improvement using momentum when minimizing non-convex functions, the results are less clear-cut. In particular, it is theoretically established that some classes of functions are such that the vanilla gradient descent already achieves optimal convergence speed, namely for functions with Lipschitz gradient \cite{lowerboundI}, or for Polyak-\L{}ojasiewicz (PL) functions \cite{PLlowerbound}. 
Among the remarkable properties of PL functions, we note that (i) it is a necessary and sufficient condition to achieve linear convergence using gradient descent \cite{abbaszadehpeivasti2023conditions}, and (ii) it can effectively model over-parameterized neural networks \cite{ liu2022loss,barboni2022global,chen2023overparameterized, xulocal}. Yet, PL functions are such that gradient descent's convergence bound cannot be significantly improved using momentum \cite{PLlowerbound}. This seems to contradict the often empirically observed benefit of momentum, and thus motivates the search for more structured relaxations of strong convexity.
An intermediate non-convex class that is expected to capture acceleration by momentum is \textit{strong quasar-convexity} \cite{hinder2020near}, sometimes referred to as \textit{weak-quasi-strong convexity }\cite{necoara2019linear,bu2020note}. A continuously differentiable function $f:\R^d \to \R $ belongs to this class if there exists a minimizer $x^\ast\in \R^d$ and parameters $(\tau,\mu) \in (0,1]\times \R_+^{\ast}$ such that for all $x\in \R^d$
\[\min_{u\in \R^d}f(u) \ge f(x) + \frac{1}{\tau}\dotprod{\nabla f(x),x^\ast-x} + \frac{\mu}{2}\norm{x-x^\ast}^2.\]
There has been growing interest in strongly quasar-convex functions in recent years \cite{gower2021sgd,wang2023continuized,alimisis2024characterization,alimisis2024geodesic,pun2024online,hermant2025continuized,chenefficient2025,lara2025delayed,de2025extending,farzin2025minimisation}.
%, whose connection to deep learning was justified in \cite{guptanesterov}. 
In the $L$-smooth setting, minimizing these functions using momentum is often referred to as \textit{accelerated optimization} \cite{hinder2020near,fu2023accelerated,hermant2024study}, suggesting better convergence bounds than gradient descent. This expectation partly stems from the fact that its complexity bounds take a form similar to those obtained in the smooth, strongly convex case. In this work, we show that this intuition is misleading.

% If a function $f : \R^d \to \R$ verifies this inequality and is $L$-smooth, gradient descent converges to the minimum value at speed $\bigO((1-\tau \mu/L)^n)$, while a suitable momentum method achieves $\bigO((1-\tau \sqrt{\mu/L})^n)$, with $n$ being the iteration number. In the case $\tau = 1$, we rates writes similarly as for $L$-smooth, $\mu$-strongly convex functions. For these functions, the condition number $L/\mu$ is typically large, so that $\mu/L \ll\sqrt{\mu/L}$, and the momentum bound is therefore significantly sharper.
% This improvement, corresponding to a squared dependence on \(\mu/L\), is commonly referred to as \emph{acceleration}.
% %It is common to mention this gain of a squared dependence on $\mu/L$ as \textit{acceleration}. 
% As it also occurs for strongly quasar-convex functions, we are tempted to conclude that a similar acceleration phenomenon occurs for these functions. In this work, we show that this intuition is misleading.




\subsection{Contributions}
%Our work addresses the question of what is missing to PL functions to ensure better convergence bounds over gradient descent when using momentum.
Our work addresses what additional structure beyond the PL condition is required to obtain improved theoretical convergence bounds over gradient descent when using momentum. We summarize our contributions as follows:

% \textbf{(1.)} We show in Section~\ref{sec:sqc_pitfalls} that, contrasting with the literature, minimizing strongly quasar-convex functions with momentum does not necessarily ensure better convergence bounds compared to using gradient descent. This result from the fact that the bounds for PL functions can be sharper in some cases.
\noindent \textbf{(1.)} 
We exhibit in Section~\ref{sec:sqc_pitfalls} flaws of strong quasar-convexity, and show that in contrast with existing interpretations, minimizing these functions with momentum does not necessarily yield better theoretical convergence bounds than gradient descent. We demonstrate that an acceleration phenomenon similar to strong convexity does not occur in general for this class.
% \noindent \textbf{(1.)} We show that the acceleration property of strongly quasar-convex functions has been stated in a manner that is overly simplistic, and optimistic. This partly results from an intrinsic flaw in its parameterization, see Section~\ref{sec:sqc_pitfalls}. \jh{est ce qu'il faut plutot dire que la comparaison sous hypothèse différentes pose des problèmes ?}


\noindent \textbf{(2.)}  We propose in Section~\ref{sec:aiming_cond} a different viewpoint by considering an \textit{aiming condition}
\begin{equation}\label{eq:intro_ac}
    \dotprod{\nabla f(x),x-x^\ast} \ge a \norm{\nabla f(x)}\norm{x-x^\ast},
\end{equation}
for $1 \ge a>0$, which quantifies how well the descent direction aims toward a minimizer $x^\ast$. We show how \eqref{eq:intro_ac} relates to strong-quasar convexity, and show that it helps to define a specific reparameterization of this latter class. In Section~\ref{sec:cv_results}, we study the convergence of gradient descent, gradient flow, the Nesterov momentum algorithm and the Nesterov ODE for functions that are both PL and satisfy the aiming condition. Using a stochastic parameterization of Nesterov momentum's algorithm, we show that we obtain accelerated bounds when the alignment constant $a$ is sufficiently large (\cref{sec:mom_cv}). 

% \noindent \textbf{(2.)} We study directly PL functions in Section~\ref{sec:cv_results}, under an \textit{aiming condition}
% \begin{equation}\label{eq:intro_ac}
%     \dotprod{\nabla f(x),x-x^\ast} \ge a \norm{\nabla f(x)}\norm{x-x^\ast},
% \end{equation}
% for $1 \ge a>0$, which quantifies how much the descent direction aims toward a minimizer $x^\ast$.  We show that it gives a stronger weight to the geometry around the minimizer (Section~\ref{sec:local_weight}), and that acceleration using momentum will hold conditionally to a large enough constant $a$ (Section~\ref{sec:mom_cv}). 

\noindent \textbf{(3.)} 
%We relax the aiming condition, requiring it only to hold in average, along the trajectory (Section~\ref{sec:avg}). Despite this relaxation, we show that we can achieve similar rate of convergence. This formalizes that we can achieve acceleration as long as the aiming condition is verified along a large enough portion of the optimization path, matching more accurately the behavior of optimizers in practice. 
We relax the aiming condition by requiring it to hold only on average along the optimization trajectory, recovering similar convergence rates (Section~\ref{sec:avg}). 
 It shows that acceleration can be achieved as long as \eqref{eq:intro_ac} holds over a sufficiently large fraction of the optimization path,
 better reflecting the behavior of practical optimization algorithms.
%
% Despite this weaker assumption, we recover essentially the same convergence rates.
% This shows that acceleration can be achieved as long as the aiming condition holds over a sufficiently large fraction of the optimization path,
% which better reflects the behavior of practical optimization algorithms.

\noindent \textbf{(4.)} We provide numerical experiment to support our theoretical insights. We exhibit a two-dimensional PL function with a unique minimizer that fails to
satisfy~\eqref{eq:intro_ac}, resulting in gradient descent to outperform momentum methods during the early stages of optimization (Section~\ref{sec:numerical_1}). We then provide an heuristic validation of our theoretical convergence bounds on the hard function for PL functions from \cite{PLlowerbound}, see Section~\ref{app:heuristic}.

% When minimizing this function, we observe that momentum methods are slower than gradient
% descent in the early stages of the optimization process, even under optimal parameter tuning.

% We design a $2$-dimensional function PL function with unique minimizer such that \eqref{eq:intro_ac} is not verified. It results in momentum algorithms minimizing this function slower than standard gradient descent for the $1000$th first iterations, even if tuning the momentum parameter to optimize the decrease at this horizon.  

% \jh{Paradigme worst-case global est très rigide}

% \jh{Idée motivation (pour montrer limite framework actuel): 2 fonctions avec même mauvaise cst de quasar-convexité, mais une avec mauvaise constante très localisée. Ces deux fonctions ont théoriquement la même vitesse de convergence, est ce que ça a du sens ?

% Autre idée motivation: la géométrie le long de la trajectoire est plus significative que la géométrie globale. Exemple de fonction avec des pathologies très ciblées qui flinguent sa régularitée, mais qui va être fortement convexe presque sûrement le long de sa trajectoire.}

% \paragraph{Tempor related works}
% Acceleration using momentum is established for neural networks under non-realistic setting. Such as very large networks \cite{wang2021modular,liu2022convergence, liao2024provable}, such that we recover a quadratic-like behavior.
%  mean field-limit \cite{wu2022mean}.


% -Intro générale 
% -importance non convexity
% -acceleration momentum: well established in (strongly-convex), more complicated in non-convex: PL mort, cas où on est presque fortement convex (neural net). Arrive strongly quasar-convex.

% Polyak-\L{}ojasiewicz (PL) functions \cite{Polyak1963,lojasiewicz1963topological}  

% \jh{Local condition are sometimes preferred/the only available \cite{xulocal}}

% There has been increasing interest in strongly quasar-convex functions (carry the hope of acceleration) in recent years \cite{hinder2020near,gower2021sgd,wang2023continuized,fu2023accelerated, hermant2024study,ablaev2025gradient,de2025extending,martinez2025smooth} + \jh{us}, whose connection to deep learning was justified in \cite{guptanesterov}.





% \jh{aiming condition gives a stronger weight to the geoemtry around minimizers}

% \jh{Relaxed assumption pour comprendre efficacité des algo d'optim}


% \jh{Plan possible: limit de SQC, alternative, moyenisation, XP. Voir pour inclure des section continues exclusives.}
\subsection{Related Works on Aiming Conditions}\label{sec:related_aiming}
Studying SGD,
\cite{liu2023aiming,guptanesterov,alimisis2025we} use conditions close to quasar-convexity they name \textit{aiming condition}. While they consider the interpretation of pointing toward a minimizer, this condition actually carries more structure. Studying gradient descent, \cite{guille2024no,singh2025directionality} consider \eqref{eq:intro_ac} from distinct perspectives: the first studies it empirically along the optimization path, the second uses it to prove a stability result, in the idea of better understanding the Edge of Stability phenomenon \cite{cohengradient}.  \cite{hermant2024study} establishes that a PL function with a unique minimizer is also strongly quasar-convex for some parameters, as long as \eqref{eq:intro_ac} is verified, but the derived parameters are such that momentum does not yield accelerated bounds, see Appendix~\ref{app:related:uaac}.
\section{Background}
In this work, $f:\R^d \to \R$ is continuously differentiable. We note $f^\ast := \min_{x \in \R^d} f(x)$, $x^\ast \in {\arg\min}_{x\in \R^d} f(x)$. 
%Noting $\mathcal{K} := \N$ or $\R_+$, $\xgen$ can indifferently be a process indexed by continuous or discrete time. 
$k \in \N$ denotes a discrete index and $t \in \R_+$ a continuous one. $\mathcal{E}(1)$ denotes the exponential law of parameter $1$. 

\subsection{Algorithms and Equations}
We introduce the algorithms and dynamical systems that we study throughout this work.
%, with their continuous counterparts.

\paragraph{Gradient Descent and Gradient Flow \:}
Among the most classical first-order algorithms is gradient descent \eqref{eq:gd}
\begin{equation}\label{eq:gd}\tag{GD}
    \tilde x_{k+1} =\tilde x_k - \gamma \nabla f(\tilde x_k).
\end{equation}
It is the building block of many important algorithms used in modern applications \cite{hinton2012rmsprop,ghadimi2013stochastic,kingma2014adam,romano2017little}, and is still an important object of study \cite{malitsky2020adaptive,cohengradient}. 
 Writing gradient descent as $ \frac{x_{k+1} - x_k}{\gamma } = -\nabla f(x_k)$, setting the \textit{Ansatz} $\tilde x_k \approx x_{k\gamma }$ for some smooth curve $\xcont$, we consider the ODE limit of \eqref{eq:gd} as the stepsize $\gamma $ vanishes to zero, named gradient flow \eqref{eq:gf} 

\begin{equation}\label{eq:gf}\tag{GF}
    \dot{x}_t = - \nabla f(x_t). 
\end{equation}
With words, \eqref{eq:gf} is a continuous version of \eqref{eq:gd}. This dynamical system viewpoint is longstanding for gaining insights into the corresponding algorithms \cite{ambrosio2005gradient, benaim2006dynamics}. %\jh{parallel limit passage mean field ?}

\paragraph{Nesterov's Momentum: Ode and Algorithm \:}
If \eqref{eq:gd} benefits from a rather simple formulation, there are many situations where modifications of this algorithm accelerate its convergence speed. Among these mechanisms, an important one is \textit{momentum}. The first momentum algorithm was Polyak's Heavy Ball \cite{polyak1964some}, which writes
\begin{equation}\label{eq:hb}\tag{HB}
     \tilde x_{k+1} = \tilde x_k + \underbrace{\alpha(\tilde x_k-\tilde x_{k-1})}_{\text{momentum}} - \underbrace{\eta \nabla f(\tilde x_k)}_{\text{gradient step}}.
\end{equation}
In this work, we focus on the Nesterov momentum \eqref{eq:nm} version \cite{nesterovbook}. 
\begin{align}\label{eq:nm}\tag{NM}\left\{
    \begin{array}{ll}
        \tilde{y}_k &= \tilde{x}_k+ \alpha_k(\tilde{z}_{k} - \tilde{x}_k),  \\
        \tilde{x}_{k+1} &=  \tilde{y}_k  - \gamma \nabla f(\tilde{y}_k),\\
        \tilde{z}_{k+1} &=   \tilde{z}_{k} + \beta_k(\tilde{y}_k - \tilde{z}_{k}) - \gamma' \nabla f(\tilde{y}_k).
    \end{array}
\right.
\end{align}
Although its formulation is slightly more involved than Polyak’s Heavy Ball,
the latter exhibits worse convergence guarantees in some settings \cite{ghadimi2015global,goujaud2025provable}.
% If it's formulation is slightly more complicated than Polyak's Heavy Ball, 
% %it is more stable and 
% the latter performs worse in some settings \cite{ghadimi2015global,goujaud2025provable}.
As for gradient descent, taking the limit of \eqref{eq:nm} as stepsize $\gamma$ vanishes yields a limit Nesterov momentum ODE \eqref{eq:nmo}, \cite[Appendix A.1]{pmlr-v202-kim23y}
\begin{align}\label{eq:nmo}\tag{NMO}\left\{
    \begin{array}{ll}
        \dot{x}_t &= \eta_t(z_t-x_t) - \gamma \nabla f(x_t), \\
        \dot{z}_t &= \eta'_t(x_t-z_t) - \gamma' \nabla f(x_t). 
    \end{array}
\right.
\end{align}
%This equation is closely related to the well known Heavy Ball equation (\jh{ref us, ou copier preuve}). 
Since the strong connection between these ODEs and the corresponding algorithm has been established \cite{suboydcandes}, their study has gained a rising interest both in the convex \hypertarget{cont_param}{{}} \cite{siegel2021accelerated,shi2018understanding,attouch2020firstorder, aujdossrondPL,li2024linear,guo2024speed} and non-convex case \cite{okamura2024primitive,hermant2024study,guptanesterov,renaud2025provably}. 
%\jh{écrire quelque part que ça s'étudiait quand même avant} 

\noindent
\textbf{Continuized Parameterization \:}
We use a specific parameterization of \eqref{eq:nm} with stochastic parameters that arises from the \textit{continuized} method \cite{even2021continuized}. Precisely, for some $\eta,\eta' \in \R$ such that $\eta+\eta'>0$, and some $\{T_k \}_{k \in \N}$ such that $T_{k+1}-T_k \overset{i.i.d}{\sim} \mathcal{E}(1)$, we fix in \eqref{eq:nm}:
\begin{align*}
    &\alpha_k = \frac{\eta}{\eta+\eta'}\bpar{1 -e^{-(\eta + \eta')(T_{k+1}-T_k)} },\\
    &\beta_k = \eta' \frac{\bpar{1- e^{-(\eta+\eta')(T_{k+1}-T_k)}}}{\eta' + \eta e^{-(\eta+\eta')(T_{k+1}-T_k)}}.
\end{align*}
With this parameterization, the algorithm is inherently stochastic, yielding convergence guarantees holding with high probability. 
%The motivation for this viewpoint is that it can be more suited to non-convex optimization 
In the case of minimizing smooth strongly-quasar convex functions, the best known bounds are achieved using this parameterization \cite{wang2023continuized}. In contrast, more classical deterministic parameterization achieves it only if assuming more structure \cite{hermant2024study}, or when using search-procedures to compute parameters \cite{hinder2020near,nesterov2021primal,guminov2023accelerated}, see details in Appendix~\ref{app:related_work_tab}.
%This viewpoint can be particularly well suited to non-convex optimization,
%where deterministic parameterization may lead to weaker results \cite{wang2023continuized,hermant2025continuized}. 
The analysis of \eqref{eq:nm} using the continuized parameterization involves a Poisson-driven SDE; we refer to \cite{hermant2025continuized} for a methodological introduction to this continuized viewpoint.
% A version of \eqref{eq:nm} with stochastic parameter appeared to circumvent theoretical limitations in non-convex setting. It arises from a \textit{continuized} point of view \cite{even2021continuized}. We use it in this work, which induces our result with \eqref{eq:nm} holds with high probability.
% \begin{assumption}[Error Bound]\label{ass:eb}
% $F: \mathbb{R}^d \mapsto \mathbb{R}$ is $\nu$-Error Bound ($\nu$-$EB$) if
% \begin{align*}
%    \forall x \in \mathbb{R}^d, \quad \lVert \nabla F(x) \rVert \geqslant \nu \lVert x- x^\ast \rVert
% \end{align*}
% \end
\subsection{Classical setting for acceleration}\label{sec:strongly_convex}
\begin{table*}[]
\caption{Comparison of convergence rates up to a numerical constant, for \eqref{eq:gf}, \eqref{eq:nmo}, \eqref{eq:gd} and \eqref{eq:nm} properly parameterized, for functions belonging to $\strongconv$, $\sqc$ and $\pl$, intersected with $\ls$ for algorithms. Most can be found in the literature, see Appendix~\ref{app:related_work_tab}.}
\label{table:sc_sqc_rates}
\vspace{0.3em}
\centering
\begin{tabular}{|l|l|l|l|l|}
\hline
\diagbox{Function class}{Equation/Algorithm}     & GF           & NMO                   & GD                     & NM                     \\ \hline
$\mu$-strongly convex ($\strongconv$)       & $ \mu$       & $ \sqrt{\mu}$         & $\mu/L$        & $\sqrt{\mu/L}$        \\ \hline
$(\tau,\mu)$-strongly quasar-convex ($\sqc$)& $\tau \mu$ & $ \tau \sqrt{\mu} $ & $\tau \mu/L$ & $\tau  \sqrt{\mu/L}$ \\ \hline
$\mu$-Polyak-\L{}ojasiewicz ($\pl$)       & $ \mu$       & $(-)$         & $\mu/L$        & $\mu/L$        \\ \hline
\end{tabular}
\end{table*}


We first introduce a geometrical setting in which momentum-based acceleration is long established.
\begin{definition}
    For $L>0$, we call $\ls$ the set of $L$-smooth functions, namely $f$ such that $\forall x,y\in \R^d,
 ~ f(x)-f(y)-\dotprod{\nabla f(y),x-y}  \le \frac{L}{2}\norm{x-y}^2$.
\end{definition}
\begin{definition}
    For $\mu>0$, we call $\strongconv$ the set of $\mu$-strongly convex functions, namely $f$ such that $\forall x,y \in \R^d,~ f(x)+\dotprod{\nabla f(x),y-x} + \frac{\mu}{2}\norm{x-y}^2 \le f(y).$
\end{definition}
Intuitively, $L$-smoothness and $\mu$-strong convexity provide global upper and lower bounds on the curvature of $f$.
% Intuitively, $L$-smoothness ($\mu$-strong convexity) defines an upper-bound (lower bound) on the curvature of $f$.
For a $C^2$ function, $f \in\strongconv \cap \ls $ implies that any eigenvalue $\lambda(x)$ of $\nabla^2 f(x)$ is such that $\mu \le \lambda(x) \le L$, for any $x\in\R^d$. The ratio $L/\mu \ge 1 $ is called \textit{condition number}, and characterizes how ill-conditioned the minimization problem is—a higher condition number typically reflecting a more difficult problem. $\strongconv \cap \ls$ is a favorable setting for minimization: there exists a unique critical point $x^\ast$, satisfying $f(x^\ast) = f^\ast$,\hypertarget{cr}{{}} and linear convergence of function values is ensured by first order methods, namely
%\textit{i.e.} for $\xgen$ 
\[\log (f(x_t)-f^\ast) \le -\Lambda t + \log(K(x_0)), \]
with $\Lambda > 0$ being the \textbf{convergence rate}, $K(x_0) > 0$ depends on the initialization. Linear convergence also holds in the discrete case, \textit{i.e.}, replacing $\xcont$ by $\xdisc$ and $t$ by $k$.
%and $\mathcal{K}$ being $\N$ or $\R_+$. 
In this work, we specifically focus on the convergence rate value $\Lambda$.
In the $\strongconv \cap \ls$ regime, the improvement of \eqref{eq:nm} over \eqref{eq:gd} is well established \cite{nesterov2004introductory}, see %a summary of 
the convergence rates in Table~\ref{table:sc_sqc_rates}. Precisely, as typically $L/\mu \gg 1$, improving from $\mu/L$ to $\sqrt{\mu/L}$ yields a substantial gain, often called \textit{acceleration}. However, many problems of interest do not fall within $\strongconv \cap \ls$. This motivates to identify more general settings where momentum-based acceleration remains possible.
 \subsection{Relaxed setting}\label{sec:relax_setting}
A first relaxation of $\strongconv$ and $\ls$ consists in lower and upper
global quadratic bounds of the function.
% Straightforward relaxations of $\strongconv$ and $\ls$ are global quadratic bounds on the function.
        \begin{definition}\label{ass:upp_qg}
                For $\mu_0,L_0 > 0$, we call: $\qginf$ the set of functions $f$ that verify
  $  f(x) - f^\ast \ge \frac{\mu_0}{2}\lVert x - x^\ast \rVert^2$, $\qg$ the set of functions $f$ that verify 
  $  f(x) - f^\ast \le \frac{L_0}{2}\lVert x - x^\ast \rVert^2$.
         \end{definition}
         $\qginf \cap \qg$ includes functions with pathological behavior, \textit{e.g.} having an infinite amount of critical points in any vicinity of $x^\ast$ \cite[Proposition 9]{hermant2024study}. This may discard the possibility of targeting a global minimum $f^\ast$.
          A key relaxation of strong convexity that does not preclude this objective is the class of Polyak-\L{}ojasiewicz (PL) functions \cite{Polyak1963,lojasiewicz1963topological}.   
\begin{definition}
    For $\mu>0$, we call $\pl$ the set of functions $f$ that satisfies $\forall x \in \R^d,~ \norm{\nabla f(x)}^2 \ge 2\mu(f(x)-f^\ast).$
\end{definition}
It immediately follows that for $f \in \pl$, any critical point is a global minimizer of $f$. From an optimization perspective, two key results hold for the class $\pl \cap \ls$: (i) gradient descent achieves a linear convergence rate \cite{Polyak1963}, and (ii) this rate is optimal among first-order algorithms \cite{PLlowerbound}. As a consequence, \eqref{eq:nm} cannot yield any significant theoretical improvement over this class, see Table~\ref{table:sc_sqc_rates}. This motivates studying more structured non-convex relaxations of $\strongconv \cap \ls$. The class of strongly quasar-convex functions (SQC) is expected to fulfill the demand.
\begin{definition}
    For $(\tau,\mu)\in (0,1]\times \mathbb{R}^\ast_+$, we call $\sqcf{\tau}{\mu}$ the set of functions $f$ that verify $\forall x\in \mathbb{R}^d,~ f^\ast \ge f(x) + \frac{1}{\tau}\dotprod{\nabla f(x),x^\ast-x} + \frac{\mu}{2}\norm{x-x^\ast}^2$.
\end{definition}
% We have that $\strongconv \subset \sqc \subset  \plf{\mu \tau^2}$ for any $\tau \in (0,1]$ \cite{hermant2024study}, making it a relaxation of $\strongconv$, while being more structured than $\pl$. 
$\sqc$ is a relaxation of $\strongconv$, while being more structured than $\pl$, see \eqref{eq:inclusions}. One possible application of these functions is the training of generalized linear models \cite[Section 3]{wang2023continuized}. Compared to $\pl$, an important supplementary structure of $\sqc$ is to keep the property of $\strongconv$ to having a unique minimizer, although it can be relaxed \cite{guptanesterov,hermant2025continuized}. Thse functions can exhibit can exhibit strong oscillations, see Figure~\ref{fig:figure_sqc}. Unlike PL functions, this class admits distinct convergence rates for \eqref{eq:gd} and \eqref{eq:nm} when intersected with $\ls$, see \cref{table:sc_sqc_rates}.

% $\sqc$ is a relaxation of $\strongconv$, while being more structured than $\pl$, see \eqref{eq:inclusions}. One possible application of these functions is the training of generalized linear models \cite[Section 3]{wang2023continuized}. Functions in $\sqc$ keep the property of $\strongconv$ to having a unique minimizer, but 
%A notable property of this class is that, unlike with PL functions, using momentum allows different convergence rates for \eqref{eq:gd} and \eqref{eq:nm}, see Table~\ref{table:sc_sqc_rates}. The rates seems to mirror those of the strongly convex case up to a $\tau$ factor. It should be kept in mind that this apparent similarity partly stems from a notational convention, using the parameter $\mu$ for both class $\strongconv$ and $\sqc$. We further discuss this point in Section~\ref{sec:sqc_pitfalls}.
% Unlike PL functions, this class admits distinct convergence rates for \eqref{eq:gd} and \eqref{eq:nm}, resembling the strongly convex case up to a factor $\tau$, see Table~\ref{table:sc_sqc_rates}.
% This apparent similarity should be treated with care, as it partly results from a notational convention that reuses the parameter 
% $\mu$ for both the classes $\strongconv$ and $\sqc$, as discussed in Section~\ref{sec:specific_pitfall}.
%Using the parameter $\mu$ for both class $\strongconv$ and $\sqc$.%that mirror those of the strongly convex case--%A notable feature of this class is that using momentum, contrary to PL-functions, one can deduce convergence rates that write similarly as in the strongly convex case 
% up to a $\tau$ factor--, see Table~\ref{table:sc_sqc_rates}. However, this apparent similarity stems from a notational convention, using the parameter $\mu$ for both class $\strongconv$ and $\sqc$. We further discuss this point in Section~\ref{sec:sqc_pitfalls}. \jh{***}

\noindent
\textbf{Relation between classes \:}
The following inclusions hold \cite{karimi2016linear,hermant2024study}
\begin{equation}\label{eq:inclusions}
    \strongconv \subset \sqcf{1}{\mu} \subset{\sqc} \subset \plf{\mu\tau^2} \subset \qginff{\mu \tau^2},\quad \ls \subset \qgf{L}.
\end{equation}
    As observed by \cite{guille2021study}, the constants inherited from these inclusions may be severely suboptimal. For instance, oscillations of the functions outside a vicinity of minimizers may affect significantly the $\ls$ constant but not the $\qg$ constant, such that we could have that $f \in \ls \cap \qg$ with $L_0\ll L$. Similarly, a sudden drop of the gradient norm may affect the $\pl$ constant while letting the $\qginf$ almost untouched, such that $f \in \pl \cap \qginf$ while $\mu_0 \gg \mu$. This induces that comparing convergences rates across different classes  must be handled with care when using same constants for different classes, $\mu$ or $L$ classically. 

       % As observed by \cite{guille2021study}, the constants inherited from these inclusions may be severely suboptimal. 
       % %This means that we should handle comparison between convergence rates across several set of functions with care. 
       % \jh{As a simple example, consider a one-dimensional example $f$ that is quadratic on most of the definition space, but such that the curvature drops on a small interval, outside of a vicinity of the minimizer. Precisely, define $f$ such that $\min f = f(0)= 0$, $f''(t) = 1$ for all $t \in \R \backslash A$, where $A$ is such that $0 \notin A$. Let us fix $A = [1,1.2]$. For $0<\varepsilon<1$, on $[1,1.1]$, we define $f''(t) = 10(\varepsilon-1)t + 11-10 \varepsilon$, on $(1.1,1.2]$ we define $f''(t) =10(1-\varepsilon)t +  12\varepsilon-11$. This definition ensure the curvature $f''$ drops on $A$, reaching a minimal value $\varepsilon > 0$, aribitrarily small. By the second-order characterization of $L$-smoothness and strong convexity, we have that $f$ thus defined belongs to $\strongconvf{\varepsilon} \cap \lsf{1}$.
       % According to the inclusions above, $f$ belongs to $\qginff{\varepsilon} \cap \lsf{1}$. However, it is direct that $f$ belongs at least to $\qginff{1/2}$, independently of $\varepsilon$. For such examples, characterizing $f$ as being a strongly convex function leads to highly suboptimal convergence bounds. $f \in \strongconvf{\varepsilon} \cap \lsf{1}$ leads to the convergence rate $\varepsilon$ for \eqref{eq:gd} and $\sqrt{\varepsilon}$ for \eqref{eq:nm}, see Table~\ref{table:sc_sqc_rates}. In contrast, using that $f \in \qginff{1/2} \cap \lsf{1} \cap \strongconvf{0}$, existing convergence rates (ignoring numerical constants) are $1/2$ for \eqref{eq:gd} \cite{guille2021study} and $1/\sqrt{2}$ for \eqref{eq:nm} \cite{aujdossrondPL}.
       % In our case, $f$ belongs to $\qginff{1/2}$, which implies rates $2$ and $\sqrt{2}$, independently of $\varepsilon$. 
       %Therefore, for the given function $f$, characterizing its convergence using strong convexity can be significantly sub-optimal.  %This strongly affects the strong convexity constant, not the quadratic growth constant.
       %As shown by the result under quadratic growth, such curvature drop does not significantly affect the convergence bound. This shows that for a given function, characterizing it with the "stronger set" possible is not necessarily relevant from an optimization point of view. ALTERNATIVE: 
       % In this example, the strong convexity constant is heavily impacted by a phenomenon--the curvature drop outside a vicinity of the minimizer-- that actually does not affect the behavior of the algorithm. In particular, using the "strongest assumption" into a function fits in to characterize its convergence is not necessarily the best solution.   }
       %we could have that $f \in \ls \cap \qg$ while $L_0\ll L$, and similarly $f \in \pl \cap \qginf$ while $\mu_0 \gg \mu$,
      % In particular, it is possible to have $f \in \ls \cap \qg$ with $L_0 \ll L$, and similarly
%$f \in \pl \cap \qginf$ with $\mu_0 \gg \mu$, see Figure~\ref{fig:quad_bounds}
%in the appendix.

% More generally, this shows that comparing convergences rates across different classes must be handle with the highest care. In particular, using the same constants for different classes--$\mu$ or $L$ classically--may misleadingly suggest that it quantify the same information. As we show in the next section, such confusion has been made between $\strongconv$ and $\sqc$.
  %     which means that the constants derived from these inclusions can be highly suboptimal. This can be particularly critical for non-convex functions, see Figure~\ref{fig:quad_bounds} in Appendix. 
       %\jh{donner deuxieme example ? Ca fera plus solide, par exemple example convex de guillestudy}. 
     %  It should be kept in mind that using the same notation for the constants of these set could lead to misleading across-classes comparisons of the convergence rates. \jh{probablement pas clair}
       %Using different constants to parameterize these sets is thus meaningful, despite these inclusions. 
       % Importantly, it means that comparing convergence rates for different classes of functions is not straightforward, as they could belong to these class for different constant $\mu$.
       


% \begin{remark}[Worst case]\label{rem:worst_case}
%     For $f \in \strongconv \cap \ls$, the "worst" situation for an algorithm is to approach the minimizer following the lowest curvature direction, as low gradient values induce slow convergence. This is why the dynamical systems \eqref{eq:gf} and \eqref{eq:nmo} have their convergence rate depending on $\mu$, see Table~\ref{table:sc_sqc_rates}. This straightforward interpretation is possible because $\strongconv$ "only" carries this interpretation. For $f \in \pl \cap \ls$, the same dependence of $\mu$ occurs, but it carries more possible interpretation. \jh{figure avec quadratique 2d sc, et une autre avec direction plateau   Our further result in Section machin will allow to separate both informations}
% \end{remark}

\section{Acceleration under Strong Quasar-Convexity: Hidden Pitfalls}\label{sec:sqc_pitfalls}
% Motivated by an analogy with the complexity bounds obtained for functions in $\strongconv \cap \ls$, see Table~\ref{table:sc_sqc_rates}, using momentum to optimize functions in $\sqc \cap \ls$ is often described as accelerated optimization \cite{hinder2020near,fu2023accelerated}. Following this line of reasoning, one might expect that \textit{for a function belonging to $\sqc \cap\ls$, the complexity bound of \eqref{eq:nm} improves upon that of \eqref{eq:gd}}. This intuition is misleading.
Using momentum to optimize functions in $\sqc \cap \ls$ has been described as accelerated optimization \cite{hinder2020near,fu2023accelerated,hermant2024study}. This would place this class as an interesting intermediate class between $\strongconv \cap \ls$ and $\pl \cap \ls$: retaining the acceleration property of the first, while allowing non-convex behavior as the second does. 
The interpretation of retaining the acceleration property of  $\strongconv \cap \ls$ is
motivated by an analogy between complexity bounds. We show that performing such analogy is misleading. 
%Then, we give evidences that the claim of acceleration is flawed in general. 

\paragraph{A hard PL instance is strongly quasar-convex}
% \jh{The previous discussion raises the question wether it is merely a matter of suboptimal characterization of the convergence speed, as for our synthetic example at the end of Section~\ref{sec:relax_setting}, or if this is more serious. The fact that the PL characterization can be sharper is of importance, as for this class there exists no accelerated rates with momentum.}
\cite{PLlowerbound} introduces a \textit{hard function} belonging to $\pl \cap \ls$ such that that the number of needed gradient calls needed to obtain a point $\tilde x$ ensuring $f(\tilde x)-f^\ast \le \varepsilon$ scales at best as $\bigO(\frac{L}{\mu}\log(\frac{1}{\varepsilon}))$ for any first-order algorithms. This bound corresponds to a $\mu/L$ convergence rate, inducing that, up to numerical constants, the convergence rate of \eqref{eq:gd} is optimal, see Table~\ref{table:sc_sqc_rates}. Perhaps surprisingly, this hard function satisfies some strong regularity, for instance it has a unique global minimizer. More interesting in our case, it is actually strongly-quasar convex.
\begin{proposition}\label{prop:hard_fonc_sqc}
   The hard function belonging to $\pl \cap \ls$ from \cite{PLlowerbound}, for which the complexity bound of \eqref{eq:gd} can be improved at most by a numerical factor, also belongs to $\sqcf{\tau}{\mu'} \cap \ls$ for some parameters $(\tau,\mu')\in (0,1]\times \R_+^\ast$.
\end{proposition}
The proof is essentially a collection of known result, see Appendix~\ref{app:hard_fonc_sqc}. This proposition definitely shows that an acceleration of the same nature than for $\strongconv \cap \ls$ is not possible in general when considering $\sqc \cap \ls$. Indeed, in the case of $\strongconv \cap \ls$, the ratio of the convergence rates between \eqref{eq:nm} and \eqref{eq:gd} is $\sqrt{\mu/L}/(\mu/L) =  \sqrt{L/\mu}$, which goes to infinity with the conditioning $L/\mu$ going to infinity. In particular, if we fix $L$, this shows that for this class, when taking $\mu \to 0$ the bound of \eqref{eq:nm} is arbitrarily better compared with \eqref{eq:gd}. Such behavior is not possible for the function from \cite{PLlowerbound}.
%, which seems to indicate that a similar gain of squared root conditioning occurs. %Because of the impossibility result for functions in $\pl \cap \ls$, 
%Precisely, following this line of reasoning, one might expect that \textit{for a function belonging to $\sqc \cap\ls$, \eqref{eq:nm} improves upon that of \eqref{eq:gd}}.%the complexity bound of \eqref{eq:nm} improves upon that of \eqref{eq:gd}}. 
%This does not hold in general. 
% We provide two observations that shows this conclusion is over-simplistic.
%\noindent \textbf{The hard-PL function from \cite{PLlowerbound} is strongly-quasar convex}


\paragraph{Weaker Assumption, Sharper Rate}
%While $\sqc \cap \ls$ defines a more structured class than $\pl\cap \ls$, convergence bounds given for the latter could actually be sharper in some cases.
%While $\sqc \cap \ls \subset \pl\cap \ls$, namely $\sqc \cap \ls$ defines a more structured class than $\pl\cap \ls$, convergence bounds given for the latter could actually be sharper in some cases. %\jh{This results from the fact that in this case, $f$ could belong to $\pl$ with a way larger $\mu$ that what can offer parameters from $\sqc$.}
%In order to compare the convergence rates of \eqref{eq:gd} and \eqref{eq:nm} for a function in $\sqc \cap \ls$, it is natural to compare the convergence rates of these algorithms known for the class $\sqc \cap \ls$, namely $\tau \mu/L$ and $\tau \sqrt{\mu/L}$, see Table~\ref{table:sc_sqc_rates}.
To compare the convergence rates of \eqref{eq:gd} and \eqref{eq:nm} when minimizing functions in
$\sqc \cap \ls$, it is natural to rely on the rates established in the literature for this
class, namely $\tau\mu/L$ for \eqref{eq:gd} and $\tau\sqrt{\mu/L}$ for \eqref{eq:nm},
see Table~\ref{table:sc_sqc_rates}.
However, any function belonging to $\sqc \cap \ls$ also belongs to $\plf{\mu'} \cap \ls$, for some parameter $\mu'$. %This mean that we could also consider the convergence rate of $\eqref{eq:gd}$ for the class $\plf{\mu'}\cap \ls$, namely $\mu'/L$. 
As a result, one
may alternatively consider the convergence rate $\mu'/L$ of gradient descent derived for functions in
$\plf{\mu'} \cap \ls$.
Strikingly, in some cases, this latter rate
\textbf{can be sharper} than those given for the class $\sqc \cap \ls$.
In Figure~\ref{fig:pl_sqc_low_corr}, we plot two examples of functions that belong to $\sqc \cap \ls$, and so also belong to $\plf{\mu'}  \cap \ls$, for some parameters. For a given pair of algorithm and function class, we numerically compute the parameters of this class such that the associated convergence rate is maximized. 
%such that the associated convergence rate is maximized.
%For each algorithm,  we numerically compute the parameters of the class of functions such that its convergence rate is maximized.
For the first example, we find that the rate of \eqref{eq:nm} characterized by $\sqc$ is $\approx 10^{2}$ times larger than the one of \eqref{eq:gd} characterized by $\plf{\mu'}$.
%(we give  an intuitive explanation in Appendix~\jh{??} of why such a difference).
However, for the second example, we find that the rate of \eqref{eq:gd} characterized by $\plf{\mu'}$ is $\approx 10^2$ times better than the rate of \eqref{eq:nm} characterized by $\sqc$. Thus, the worst-case bounds of \eqref{eq:gd} and \eqref{eq:nm} under $\sqc \cap \ls$ may be improved if considering the worst-case bound of \eqref{eq:gd} under $\pl \cap \ls$. 




These findings indicate that our understanding of acceleration with strong quasar convexity is lacking. In fact, this class suffers from some problems, as we discuss in the next section.
\subsection{Problems of Strong Quasar-Convexity}\label{sec:specific_pitfall}
The intuition that a similar acceleration phenomenon occurs for $\sqc \cap \ls$ and $\strongconv \cap \ls$ can results from an appearing improvement from $\mu/L$ with \eqref{eq:gd} to $\sqrt{\mu/L}$ with \eqref{eq:nm} in both cases, see Table~\ref{table:sc_sqc_rates}, and from carrying the interpretation of $L/\mu$ as a condition number. This is somewhat deceptive, as we detail below. 

\paragraph{Parameter Interpretation}
The constant $\mu$ in $\strongconv$ parametrizes a global lower bound on the curvature. Reusing the notation $\mu$ in a similar way in $\sqc$ may therefore suggest a similar interpretation. In particular, for $f \in \sqc \cap \ls$, one may be tempted to view $L/\mu$ as a conditioning measurement, analogously to the setting $\strongconv \cap \ls$. Under this interpretation, one would expect $L/\mu \ge 1$, see Section~\ref{sec:strongly_convex}. 
However, since $f \in \sqcf{\tau}{\mu} \Rightarrow f \in \sqcf{\theta \tau}{\mu/\theta}$, for any $\theta \in (0,1]$ \cite[Observation 5]{hinder2020near}, it is possible to have $f \in \sqcf{\tau}{\mu} \cap \ls$ for some parameters satisfying $\mu > L$. This observation challenges both the condition-number interpretation of $L/\mu$ and calls into question the interpretation of the convergence rates in this setting.
% g

% Importantly, it means we can have $f \in \sqcf{\tau}{\mu} \cap \ls$ for some parameters satisfying $\mu > L$, questioning both the condition-number interpretation and the convergence rate interpretation.
%Moreover in this case,  transforming the convergence rate into $\sqrt{\mu/L}$ is not an improvement. This questions the interpretation of the associated convergence rates.
%Note also that if the rates of \eqref{eq:gf} and \eqref{eq:gd} are uniform with respect to the aforementioned reparameterization, it is not the case for \eqref{eq:nmo} and \eqref{eq:nm}, which worsen by a factor $\sqrt{\theta}$.

\paragraph{The Optimal Choice of Parameters}
For many function classes, there exists a natural notion of optimal parameters.
For instance, if $f \in \pl$
for a range of constants $\mu>0$, the tighter class to describe $f$ is $\plf{\hat{\mu}}$ with $\hat \mu$
the largest admissible value.
Accordingly, the rates reported in Table~\ref{table:sc_sqc_rates} are maximized by choosing $\mu = \hat \mu$.
In contrast, because of the double parameterization, there is no
canonical way to choose among all admissible pairs $(\tau,\mu)$ an optimal one for a function in $\sqc$. Choosing the largest admissible $\tau$ and then the largest admissible $\mu$ leads to a different parameter pair than optimizing in the reverse order, see Appendix~\ref{app:choice_param_sqc}.
%that optimize the convergence
%bounds of Table~\ref{table:sc_sqc_rates}.
%In particular, optimizing the rate with respect to $\tau$ first and then $\mu$
%leads to a different parameter choice than optimizing with respect to $\mu$ first
%and then $\tau$ (see %Appendix~\ref{app:choice_param_sqc}).
More critically, for a given function $f$, among the admissible pair $(\tau,\mu)$ ensuring $f \in \sqc$, the one 
that maximizes a given convergence rate depends on the algorithm or
dynamic under consideration, as illustrated in \cref{fig:sqc_flaws}.
%, comparison of convergence rates across algorithms should be regarded with cares. 
As a consequence, statements of the form
\emph{“the rates of \eqref{eq:gf} and \eqref{eq:nmo} scale as
$\tau\mu$ and $\tau\sqrt{\mu}$”} without clarifying how the parameters are chosen may result in misleading comparisons.
\begin{figure}
    \centering
    \includegraphics[width=0.3\linewidth]{figures/f_1d_graph.pdf}
    \hspace{-7pt}
    \includegraphics[width=0.3\linewidth]{figures/surface_fx_y.png}
    \begin{tabular}{|l|l|l|}
\hline
Algorithm (class) & $\Lambda$ (left plot)   & $\Lambda$ (right plot)       \\ \hline
GD ($\pl$)        & $3.2 \cdot 10^{-4}$  & $1.6 \cdot 10^{-2}$ \\ \hline
GD ($\sqc$)       & $1.3 \cdot 10^{-2}$   & $7.7 \cdot 10^{-5}$\\ \hline
NM ($\sqc$)       & $1.8 \cdot 10^{-2}$ & $1.8 \cdot 10^{-4}$  \\ \hline
\end{tabular}
    \caption{ Left: graph of $f(t) = 5(t+0.19\sin(5t))^2$, minimizer is $0$. Right: graph of $f(x,y) = 0.5(0.5x^2 -y)^2 + 0.05x^2$, minimizer is $(0,0)$. Bottom: table giving for each function the numerical value of the theoretical convergence rate for \eqref{eq:gd} or \eqref{eq:nm}, precising what class of function is used to characterize the bound (see Table~\ref{table:sc_sqc_rates}), where $\Lambda$ is the convergence rates. See implementation details in Appendix~\ref{app:param_details}. It shows that depending on the functions, $\pl$-based bound may be sharper than  $\sqc$-based bounds, and conversely.}
    \label{fig:pl_sqc_low_corr}
\end{figure}
\begin{figure}[t]
    \centering
    %\includegraphics[width=0.8\linewidth]{figures/sqc_bound_cont.pdf}
    \includegraphics[width=0.85\linewidth]{figures/sqc_bound.pdf}


     \caption{For a range of parameters $\tau$, we numerically compute the highest admissible $\mu$ ensuring $f \in \sqc $, and plot the numerical values of the associated theoretical convergence rates for \eqref{eq:gd} and \eqref{eq:nm} under $\sqc \cap \ls$, namely $\tau \mu/L$ and $\tau \sqrt{\mu/L}$, with $f(t) = 5(t+0.19\sin(5t))^2$.  See implementation details in Appendix~\ref{app:param_details}.
    %For each $\tau$, we empirically find the highest $\mu$ such that the function belongs to $\sqc$. We only consider value of $\tau$ such that the associated $\mu$ is positive. Note that the pair 
Note that the pair $(\tau, \mu)$ that maximizes the convergence rate of \eqref{eq:gd} differs from the pair that maximizes the one of \eqref{eq:nm}.} 
    \label{fig:sqc_flaws}
\end{figure}

% These first observations are induced by the double parameterization in $\sqc$. We note that these issues are conceptually solved by replacing these assumptions by belonging to $\sqcf{\tau}{0} \cap \qginf$, where $\sqcf{\tau}{0}$ is the set of $\tau$-quasar convex functions. First, the interpretation is clearly that $\mu_0$ parameterize a lower bound on the local curvature around the minimizer, and $\tau$ is \jh{complete}

% \subsection{Strong Quasar-Convexity and Acceleration}\label{sec:sqc_non_accelerate}
% In the previous section, we discussed the problem of the interpretation and choice of the parameters for the class $\sqc$, which raise questions about how to interpret the existing convergence rates. Beyond this, we show now that i) belonging to $\sqc \cap \ls$ does not ensure the convergence bound of \eqref{eq:nm} to be better than the one of \eqref{eq:gd}, and ii) one can deduce from the literature about PL functions that there exists functions in $\sqc \cap \ls$ for which gradient's descent convergence speed cannot be significantly improved.



% \section{Acceleration under Strong Quasar-Convexity: Hidden Pitfalls}\label{sec:sqc_pitfalls}

% As improving over gradient descent's convergence rate using momentum is not possible with $\pl \cap \ls$, this motivates to study more structured non-convex relaxations of $\strongconv \cap \ls$. The following class seems to fulfill this demand.
% \begin{definition}
%     For $(\tau,\mu)\in (0,1]\times \mathbb{R}^\ast_+$, we call $\sqcf{\tau}{\mu}$ the set of $(\tau,\mu)$-strongly quasar-convex functions, namely $f$ that verify $x\in \mathbb{R}^d,~ f^\ast \ge f(x) + \frac{1}{\tau}\dotprod{\nabla f(x),x^\ast-x} + \frac{\mu}{2}\norm{x-x^\ast}^2$.
% \end{definition}
% We have that $\strongconv \subset \sqc \subset  \plf{\mu \tau^2}$ for any $\tau \in (0,1]$ \cite{hermant2024study}. One possible application of these functions is the training of generalized linear models \cite[Section 3]{wang2023continuized}. Functions in $\sqc$ keeps the property of $\strongconv$ to having a unique minimizer, but can exhibit strong oscillations, see Figure~\ref{fig:figure_sqc}.
% \begin{assumption}\label{ass:sqc}
%   $f$ verifies
%     $f^\ast \ge f(x) + \frac{1}{\tau}\dotprod{\nabla f(x),x^\ast-x} + \frac{\mu}{2}\norm{x-x^\ast}^2$ for all $x\in \mathbb{R}^d$, and some $(\tau,\mu)\in (0,1]\times \mathbb{R}^\ast_+$.
% \end{assumption}
% An important feature of this class is that using momentum, contrary to PL-functions, one can deduce convergence rates that write similarly as in the strongly convex case, for which acceleration occurs. We summarize the rates in 

%The previous reasoning also hold in other settings, see Appendix \jh{app}. 
% The previous discussion shows that $\sqc \cap \ls$ does not guarantee improved convergence rates using \eqref{eq:nm} in general.




% This shows that the pathological examples designed to prove optimality of gradient descent satisfies a strongly-quasar convex structure.
% In our numerical experiments to minimize this function, we tune optimally \eqref{eq:gd} and \eqref{eq:nm}, and were unable to achieve any acceleration with \eqref{eq:nm}. This indicates that the numerical constants are artificial. This indicates that the acceleration property \jh{avant ou apres ?} 
\subsection{Discussion}
The previous discussions shows that the potential benefit provided by using \eqref{eq:nm} instead of \eqref{eq:gd} in strong quasar-convex optimization is not yet fully understood. We note that the problems raised in \cref{sec:specific_pitfall} can be solve if splitting $\sqc$ into two assumptions, \textit{i.e.} assuming $f \in \qginf$ and that it is quasar-convex, namely $\exists \tau \in(0,1]$ such that $f^\ast \ge f(x) + \frac{1}{\tau}\dotprod{\nabla f(x),x^\ast-x}$. However, the other points remain. The problem stems in the fact that convergence rates deduced under $\pl$ are not directly comparable to those deduce for classes such as $\sqc$, as the underlying geometrical constants encompasses different information. Our idea in the remaining of the paper is then to keep working with the $\pl$ class, and pairing it with an additional and intuitive geometrical information.



\section{Aiming condition as an Alternative Characterization of Convergence}\label{sec:aiming_cond}

We introduce a class of function that we will further intersect with $\pl$.
We assume that there exists a unique global minimizer $x^\ast$.
% In order to understand why $\sqc$ enables to deduce different convergence rates for \eqref{eq:gd} and \eqref{eq:nm}, which is not possible with $\pl$, we can try to understand what is missing to functions in $\pl$ to belong to $\sqc$. A first straightforward difference is that $\sqc$ forces uniqueness of minimizer, which is not the case with $\pl$. An earlier work identified the last ingredient.
% $\pl$ does not allow acceleration, while $\sqc$ can. To better understand why, one can try to understand how $\sqc$ is more structured than $\pl$. A first straightforward difference is that only $\sqc$ forces uniqueness of minimizer.
% An earlier work identified the last ingredient.
% \begin{theorem}[\cite{hermant2024study}]\label{thm:uaac}
%     Let $f \in \pl$, with a unique minimizer $x^\ast$. Then, there exists $(\tau,\mu') \in (0,1] \times \mathbb{R}^\ast_+$ such that $f\in \sqcf{\tau}{\mu'}$ if and only if
%           \begin{equation}\label{eq:uaac}
%                           \exists a > 0,  \forall x \in \mathbb{R}^d, \quad 1 \geqslant \frac{\langle \nabla f(x),x - x^\ast \rangle}{\lVert  \nabla f(x) \rVert \lVert x - x^\ast \rVert} \geqslant a> 0.
%           \end{equation}
% \end{theorem}
% Intuitively, the condition \eqref{eq:uaac} indicates that the descent direction must form an acute angle with the direction pointing towards the minimizers.
We initially assume the existence
of a unique minimizer $x^\ast$. We note that as discussed in \cref{sec:sqc_pitfalls}, functions that serve as lower bound for $\pl \cap \ls$ are designed in this setting, such that it remains an interesting setting to discuss acceleration. Anyway this assumption can be relaxed, as discussed in
Section~\ref{sec:avg} 
\begin{definition}\label{def:ac}
    For $a\in(0,1]$, we call $\ac$ the set of functions satisfying the \textit{aiming condition},\textit{ i.e.} $\forall x\in \R^d,~ \dotprod{\nabla f(x),x-x^\ast} \ge a \norm{\nabla f(x)}\norm{x-x^\ast}$.
\end{definition}


The parameter $a$ measures the alignment between the descent direction $-\nabla f$ and the
direction toward the minimizer. Values of $a$ close to $1$ indicate a strong alignment. In one dimension, any function in $\ac$ has perfect alignment, yielding $a=1$.
Such strong alignment  is particularly favorable to momentum methods, as it ensures that
the accumulated momentum reinforces progress toward $x^\ast$.
In contrast, if $a$ is small, momentum may amplify steps taken in poorly aligned directions.
For example, the rightmost function in Figure~\ref{fig:pl_sqc_low_corr} exhibits
a small aiming condition value, $a \approx 3\cdot10^{-4}$, which helps to explain why the convergence rate of
\eqref{eq:gd} is sharper than that of \eqref{eq:nm} in this case.  We discuss the important connection between $\ac$ and $\sqc$ in Appendix~\ref{app:link}. Before, we provide a toy construction that belongs to $\ac$.




% Although the aiming condition benefits of a direct intuition, it has seldom be used as is as an assumption. We first give a toy example belonging to this set. Then, we discuss its link with other assumption, and then a limit of this characterization.

\paragraph{A toy example}
 Let $f:\R_+ \to \R_+$ such that $f(0) = 0$ and $f'(t)> 0$ for all $t> 0$. Let $g : \R^d \to \R$, such that $g\ge 1$. Define $h :\R^d \to \R$
  \begin{equation}\label{eq:h}
    h(x) = f(\lVert x \rVert)g\left(\frac{x}{\lVert x \rVert} \right).
\end{equation}
$f$ describes the radial movement, while $g$ describes the angular oscillations. Such constructions were introduced as quasar-convex examples \cite{lee2016optimizingstarconvex}. In particular, $f \in \sqc$ implies $h \in \sqc$ \cite[Proposition 8]{hermant2024study}. 
 The more $g$ oscillates, the lower is $a$. This is formalized with the following result.
\begin{proposition}\label{lem:toy_example_ac}
Let $h$ defined as \eqref{eq:h}. Then, $h \in \ac$ with
    \[a= \inf_{r>0,u \in \mathbb{S}^{d-1}} \frac{1}{\sqrt{1+ \bpar{\frac{f(r)}{rf'(r)}}^2 \bpar{\frac{\norm{\nabla_{\sphere} g(u)}}{g(u)}}^2}},  \]
    where $\nabla_{\sphere} g(u)$ is the spherical gradient of $g$ restricted to $ \mathbb{S}^{d-1}$.
\end{proposition}
The spherical gradient $\nabla_{\sphere} g(u)$ encodes the angular variation. If $g$ is constant or has no angular variation, \textit{e.g.} $g(x) = \norm{x}^2$, then $\nabla_{\sphere} g(u)= 0$. In this case $h$ is a radial function and it follows that $h \in \acf{1}$. Supplementary structure on $f$--such as belonging to $\pl$ or $\qginf$--will transfer to $h$. However it appears clearly in our result that it will also impact the value of $a$ through the ratio $f(r)/(rf'(r))$.
\subsection{Link between aiming condition and strong quasar convexity}\label{app:link}
In this section, we investigate the link between $\ac \cap \pl$ and relaxations of strong convexity. 
\begin{lemma}\label{lem:scalar_prod_bound}
Let $f \in \pl \cap \ac \cap \qg $. Let $
\mu_0 := \sup\{\mu' \ge \mu : f \in \qginff{\mu'}\}$.
Then, for all $x\in\R^d$,
\[ f^\ast \ge f(x) + \frac{1}{a}\sqrt{\frac{L_0}{\mu}}
\dotprod{\nabla f(x), x^\ast - x}
 + \frac{\sqrt{\mu \mu_0}}{2}\norm{x-x^\ast}^2.
\]
That is, $f \in \sqcf{a\sqrt{\mu/L_0}}{~ \sqrt{\mu \mu_0}}$.
\end{lemma}
\begin{proof}
If $x=x^\ast$, the inequality holds trivially. Assume $x\neq x^\ast$.

\medskip
\noindent\textbf{Step 1 (AC + PL).}
Since $f\in\ac$, we have
\[
\dotprod{\nabla f(x),x-x^\ast} \ge a \norm{\nabla f(x)}\norm{x-x^\ast},
\quad\text{hence}\quad
\dotprod{\nabla f(x),x^\ast-x} \le -a \norm{\nabla f(x)}\norm{x-x^\ast}.
\]
Since $f\in\pl$, we also have $\norm{\nabla f(x)} \ge \sqrt{2\mu}\sqrt{f(x)-f^\ast}$, so
\begin{equation}\label{eq:pl_ac_core}
\dotprod{\nabla f(x),x^\ast-x}
\le -a\sqrt{2\mu}\,\sqrt{f(x)-f^\ast}\,\norm{x-x^\ast}.
\end{equation}

\medskip
\noindent\textbf{Step 2 (use $\qg$).}
Since $f\in\qg$,
\[
f(x)-f^\ast \le \frac{L_0}{2}\norm{x-x^\ast}^2
\quad\Rightarrow\quad
\norm{x-x^\ast}\ge \sqrt{\frac{2}{L_0}}\sqrt{f(x)-f^\ast}.
\]
Plugging this into \eqref{eq:pl_ac_core} yields
\begin{equation}\label{eq:pl_ac_0_corr}
\dotprod{\nabla f(x),x^\ast-x}
\le -a\sqrt{2\mu}\sqrt{\frac{2}{L_0}}\,(f(x)-f^\ast)
= -2a\sqrt{\frac{\mu}{L_0}}\,(f(x)-f^\ast).
\end{equation}

\medskip
\noindent\textbf{Step 3 (use $\qginff{\mu_0}$).}
By definition of $\mu_0$, we have
\[
f(x)-f^\ast \ge \frac{\mu_0}{2}\norm{x-x^\ast}^2
\quad\Rightarrow\quad
\sqrt{f(x)-f^\ast}\ge \sqrt{\frac{\mu_0}{2}}\norm{x-x^\ast}.
\]
Plugging this into \eqref{eq:pl_ac_core} yields
\begin{equation}\label{eq:pl_ac_1_corr}
\dotprod{\nabla f(x),x^\ast-x}
\le -a\sqrt{2\mu}\sqrt{\frac{\mu_0}{2}}\norm{x-x^\ast}^2
= -a\sqrt{\mu \mu_0}\,\norm{x-x^\ast}^2.
\end{equation}


\medskip
\noindent\textbf{Step 4 (combine the two bounds).}
From \eqref{eq:pl_ac_0_corr} and \eqref{eq:pl_ac_1_corr}, we have
\[
\dotprod{\nabla f(x),x^\ast-x}
\le -\max\Bigl\{2a\sqrt{\frac{\mu}{L_0}}\,(f(x)-f^\ast),\; a\sqrt{\mu\mu_0}\,\norm{x-x^\ast}^2\Bigr\}.
\]
Since $\max\{u,v\}\ge \frac{u+v}{2}$ for $u,v\ge 0$, we obtain
\[
\dotprod{\nabla f(x),x^\ast-x}
\le
-\frac12\Bigl(
2a\sqrt{\frac{\mu}{L_0}}\,(f(x)-f^\ast) + a\sqrt{\mu\mu_0}\,\norm{x-x^\ast}^2
\Bigr),
\]
which is the claim.
\end{proof}
This lemma shows that $\pl \cap \ac \cap \qg $ induces a specific parameterization of strong-quasar convexity.  We note that this parameterization does not suffer from the problem raised in Section~\ref{sec:specific_pitfall}. Importantly, it also follows that similar convergence analysis ideas can be used for $\pl \cap \ac \cap \qg $ and $\sqc$.  
Conversely, if assuming some smoothness regularity, strong-quasar convexity implies the aiming condition.
\begin{proposition}\label{prop:converse_link}
    Let $f \in \sqc$. If $f$ also satisfies 
    \begin{equation}
         \forall x \in \R^d,~ \norm{\nabla f(x)}\le L\norm{x-x^\ast},
    \end{equation}
    then $f \in \acf{\frac{2\mu}{(2-\tau)L}}$.
\end{proposition}
See the proof in Appendix~\ref{app:prop:converse}.





%In particular it indicates that stating "\textit{under $\sqc$, compared to $\eqref{eq:gd}$, better convergence bounds for \eqref{eq:nm} hold}" is misleading. 
% \textbf{The optimal choice of parameters \:}
%  For most classes of functions, there exists a canonical way to chose the "optimal" parameters for this class, at least theoretically. By instance in the case of $\pl$, for a given $f$, 
%  %µ\jh{(}we have that $\mu > \mu'$ implies $f \in \pl \Rightarrow f \in \plf{\mu'}$\jh{)}, and each convergence rate for $\pl$
%   each convergence rate in Table~\ref{table:sc_sqc_rates} is optimized with the choice $\mu^\star := {\arg \sup}_{\mu > 0} \{\mu, s.t. f \in \pl \}$. In the case of $\sqc$ there is no canonical way to chose the parameters $(\tau,\mu)$ that optimize the convergence bounds. In particular, $\arg \sup_\tau \arg \sup_\mu \{ f \in \sqc \}$ and $\arg \sup_\mu \arg \sup_\tau \{ f \in \sqc \}$ lead to different choice of parameters $(\tau,\mu)$, see Appendix~\ref{app:choice_param_sqc}. Worse, as shown in Figure~\ref{fig:sqc_flaws}, the pair of parameters that maximizes a given convergence rate vary cross algorithms/equations. This means that stating "\textit{the rates of \eqref{eq:gf} vs \eqref{eq:nmo} compares as $\tau \mu$ vs $\tau \sqrt{\mu}$}" is misleading.

% \textbf{Strongly quasar-convex vs PL \:}



% Contrarily to $\pl$, $\sqc$ enables to say that there is a possibility that for such functions, the rate in the bound of \eqref{eq:nm} is better than \eqref{eq:gd}. One may question what makes this occur. Restricting to $f\in \pl$ with a unique minimizer, a former result establish what is needed to also belong to $\sqc$ for some $(\tau,\mu)$.

% \begin{theorem}[\cite{hermant2024study}]\label{thm:uaac}
%     Let $f \in \pl$, with a unique minimizer. Then, there exists $(\tau,\mu') \in (0,1] \times \mathbb{R}^\ast_+$ such that $f\in \sqcf{\tau}{\mu'}$ if and only if there exists some $a > 0$ such that
%            \begin{equation}\label{UAAC CONDITION}
%              \forall x \in \mathbb{R}^d, \quad 1 \geqslant \frac{\langle \nabla f(x),x - x^\ast \rangle}{\lVert  \nabla f(x) \rVert \lVert x - x^\ast \rVert} \geqslant a> 0.
%          \end{equation}
% \end{theorem}
% This result indicates that the descent direction must form an acute angle with the direction pointing toward the minimizer. This intuitively hint why it enables a theoretical benefit of momentum, as it ensures that the momentum pushes the trajectory in the right direction, namely toward the minimizer. This questions rather a constant $a$ with low value significantly degrades the rate of $\sqc$. In particular, one may expect $\pl$ to be less sensitive to this quantity. We perform empirical computation of $L$, the best $\mu$ in $\pl$, and the best pair $(\tau,\mu)$ in $\sqc$ to maximize the rate of \eqref{eq:nm} (see Table~\ref{table:sc_sqc_rates}), for the two functions displayed in Figure\jh{??}. For the 1d function (impliying $a = 1$) displayed on the left plot, we find that the rate of \eqref{eq:nm} characterized by $\sqc$ is $\approx 10^{12}$ times larger than the one of \eqref{eq:gd} characterized by $\pl$ (we give  an intuitive explanation in Appendix~\jh{??} of why such a difference). The function displayed on the right side is built such that $a$ is very low on the displayed domain. We find that the rate of \eqref{eq:gd} characterized by $\pl$ is $\approx 10^2$ times better than the rate of \eqref{eq:nm} characterized by $\sqc$. Importantly, it indicates that we can have $f \in \sqc$, while having that rate of convergence with momentum is not better than \eqref{gd}.








\section{Acceleration for PL Functions with an Aiming Condition}\label{sec:cv_results}

% It is intuitive that belonging to $\ac$ with $a$ close to $1$ will benefit to momentum methods, as it ensures that the momentum pushes the trajectory in the right direction, namely toward the minimizer. In particular, the right-most function displayed in Figure~\ref{fig:pl_sqc_low_corr} is such that $a$ is very small, which hints why the bound of \eqref{eq:gd} characterized with $\pl$ is better.

In Section~\ref{sec:local_weight}, we observe the impact of the aiming condition when studying standard gradient methods \eqref{eq:gf} and \eqref{eq:gd}. We consider in Section~\ref{sec:mom_cv} the convergence of 
%momentum methods 
\eqref{eq:nmo} and \eqref{eq:nm}, and the question of acceleration.

% To quantify the impact of the aiming condition, we focus on the class
% $\pl \cap \ac$, allowing to draw new insights on convergence behavior.
% Notably, this class addresses two of the issues identified in
% Section~\ref{sec:sqc_pitfalls} for $\sqc$: it yields a clearer interpretation of
% the parameters and admits a canonical choice of optimal constants, namely the largest admissible values. \jh{peut être faire miniplan: section 4.1, 4.2}


% The parameter $a$ measures the alignment between the gradient direction and the
% direction toward the minimizer.
% In one dimension, any function in $\ac$ necessarily has $a=1$, while larger values
% of $a$ more generally indicate better alignment.
% This property is favorable to momentum methods, whereas small values of $a$ may
% lead to poorly aligned momentum directions, as illustrated in
% Figure~\ref{fig:pl_sqc_low_corr}.
% To quantify the impact of the aiming condition, we focus on the class $\pl \cap \ac$,
% which admits a clearer parameter interpretation and a canonical choice of optimal
% constants, namely the largest admissible values.

% In order to push further this intuition and incorporate it in our analysis, we introduce the following class of functions.
% \begin{definition}\label{def:ac}
%     For some $1\ge a >0$, we call $\ac$ the set of functions that verify the \textit{aiming condition}, namely $\forall x\in \R,~ \dotprod{\nabla f(x),x-x^\ast} \ge a \norm{\nabla f(x)}\norm{x-x^\ast}$.
% \end{definition}
% As a special case, $1$-dimensional functions belong to $\acf{1}$. 


\subsection{$\ac$ Gives Weight to the Local Information at $x^\ast$}\label{sec:local_weight}
We first compare convergence bounds obtained for \eqref{eq:gf}, depending on $f$ belongs to $\ac$ or not.
\begin{theorem}\label{thm:gf_pl_&_ac}
    Let $\xcont \sim \eqref{eq:gf}$. 

    \noindent
    (i) \cite[Theorem 3]{polyak2017lyapunov}  If $f \in \pl$,  then
    $$ f(x_t)-f^\ast \le (f(x_0)-f^\ast)e^{-2 \mu t}. $$
    (ii) If $f \in \pl \cap \ac \cap \qg$, then
    $$ f(x_t)-f^\ast \le \bpar{1+\sqrt\frac{L_0}{\mu_0}}(f(x_0)-f^\ast)e^{-a\sqrt{\mu \mu_0}t}, $$
    where $\mu_0 := \sup\{\mu' \ge\mu : f \in \qginff{\mu'}\} $.
\end{theorem}
See the proof in Appendix~\ref{app:gf_pl_&_ac}. The assumption $f \in  \qg$ for bound (ii) is fairly mild. When considering algorithms, we will assume $f \in \ls$, such that they will naturally belong to $\qg$. We note that under the aiming condition, there is a worse constant factor $1+\sqrt{L_0/\mu_0}$, which mainly affects the bound for early iterations. In this work, we specifically focus on \hyperlink{cr}{convergence rates}.

\paragraph{Comparison of Convergence Rates} In addition to the PL constant $\mu$, the \hyperlink{cr}{convergence rates} in (ii) depends on a quadratic lower bound parameter $\mu_0$ and the aiming condition parameter $a$. As a result, rate (ii) improves upon rate (i) whenever $a \ge 2\sqrt{\mu/\mu_0}$. So, for functions with bad alignment, characterizing the convergence rate solely through the PL condition may lead to sharper rates. It is interesting to observe that $\mu_0$ appears in (ii), but not in (i). In the presence of flat plateaus, we may have $\mu \ll \mu_0$. This is due to $\mu_0$ being less affected by global irregularities of the function landscape, as we discuss next.

 \noindent\textbf{The Aiming Condition Exploits the Local Information at $x^\ast$ \:} For functions in $\strongconv \cap \ls$, the slowest direction toward the minimizer is the one with the smallest curvature~$\mu$, since it corresponds to the flattest direction (left plot of Figure~\ref{fig:worst_case} in the appendix).  
For functions in $\pl \cap \ls$, the situation is less clear: the trajectory may approach the minimizer either through low-curvature directions near $x^\ast$ or through plateau-like regions (right plot of Figure~\ref{fig:worst_case}). Thus, it is ambiguous whether the relevant $\mu$ is determined by local curvature at $x^\ast$ or by the plateaus, as $\pl$ handles both information.
Bound~(i) of Theorem~\ref{thm:gd_pl_&_ac} does not separate these effects, but bound~(ii) does: it reveals a trade-off between global plateaus and local curvature~$\mu_0$. In the special case $a=1$, if $\mu=\mu_0$, flatness near $x^\ast$ dominates and we recover the rate $\mathcal{O}(e^{-\mu t})$. If $\mu<\mu_0$, the plateaus matter, yielding a rate $e^{-\sqrt{\mu\mu_0}}$, which shows that plateaus do not fully determine the convergence.

\paragraph{Convergence of \eqref{eq:gd}}  Assuming the function belongs to $\ls$, we deduce similar convergence rates for \eqref{eq:gd}, with the difference that the $L$-smoothness constant intervenes.
\begin{theorem}\label{thm:gd_pl_&_ac}
    Let $f \in \ls$, $\xdisctilde \sim \eqref{eq:gd}$, $\gamma = \frac{1}{L}$.

    \noindent
    (i) \cite[Theorem 1]{karimi2016linear} If we also have $f \in \pl$,  then
    $$ f(\tilde x_k)-f^\ast \le \bpar{1-\frac{\mu}{L}}^k(f(\tilde x_0)-f^\ast). $$
    (ii) If we also have $f \in \pl \cap \ac $, then
    $$ f(\tilde x_k)-f^\ast \le  K_0\bpar{1-a\frac{\sqrt{\mu \mu_0}}{L}}^k(f(\tilde x_0)-f^\ast). $$
    with $K_0 := \bpar{1+\sqrt\frac{L_0}{\mu_0}}$, $\mu_0 := \sup\{\mu' \ge \mu: f \in \qginff{\mu'}\} $ and $L_0 :=\inf\{L' \le  L:f \in \qgf{L'}\}$.
\end{theorem}
Broadly speaking, the convergence rates of Theorem~\ref{thm:gd_pl_&_ac} are the same as those of Theorem~\ref{thm:gf_pl_&_ac}, divided by $L$, such that we can draw similar insights, see the proof in Appendix~\ref{app:gd_pl_ac}.

\subsection{Acceleration using Momentum with Large Enough Aiming Condition Constant}\label{sec:mom_cv}
We state the convergence rates of \eqref{eq:nmo} and \eqref{eq:nm}, under the aiming condition.
\begin{theorem}\label{thm:pl_accel}
Let $f\in \pl \cap \ac \cap \ls$, $L_0 := \inf \{L'\le L:  f \in \qgf{L'} \}$, $\mu_0 := \sup\{\mu' \ge \mu : f \in \qginff{\mu'}\}$.

(i)  Let $\xcont \sim \eqref{eq:nmo}$ and constant parameters $\gamma\ge 0$,  $\gamma' = (\mu_0 L_0)^{-1/4}$, $\eta = (\mu_0 L_0)^{1/4}$, $\eta' = a\bpar{\mu_0/L_0}^{1/4}\sqrt{\mu}$. Then,
     \begin{equation*}
     f(x_t)-f^\ast \le K_0 e^{-a \bpar{\frac{\mu_0}{L_0}}^{\frac{1}{4}}\sqrt{\mu}t},
 \end{equation*}
 where $K_0 :=  \bpar{1+\sqrt\frac{L_0}{\mu_0}}(f(x_0)-f^\ast)$ 
 
(ii) Let $\xdisctilde$ verify \eqref{eq:nm} with the \hyperlink{cont_param}{continuized parameterization}, with 
$\eta = (\mu_0 L_0)^{1/4}/\sqrt{L},\quad \gamma = 1/L,\quad \gamma' = (( \mu_0 L_0)^{1/4}\sqrt{ L})^{-1}, $ and $\eta' = a\bpar{\mu_0/L_0}^{1/4}\sqrt{\mu/L}$. Then,  
    with probability  $1-\frac{1}{c_0} - \exp(-(c_1-1-\log(c_1))k)$, for some $c_0 > 1$ and $c_1 \in (0,1)$ we have
$$ f(\tilde x_k)-f^\ast \le K_1e^{-a\bpar{\frac{\mu_0}{L_0}}^{1/4}\sqrt{\frac{\mu}{L}}(1-c_1)k}, $$
where $K_1 := c_0 \bpar{1+\sqrt\frac{L_0}{\mu_0}}(f(x_0)-f^\ast)$.
\end{theorem}

See the proof in Appendix~\ref{app:pl_accel}. The result of Theorem~\ref{thm:pl_accel} (ii) holds under high probability, due to the use of the \textit{continuized} version of Nesterov. %It is not clear whether this result can be obtained with a more classical parameterization of \eqref{eq:nm}. 
In the following, we compare convergence rates despite this difference.

\noindent
\textbf{Comparison of the Convergence Rates: Discrete Case \:} The \hyperlink{cr}{convergence rates}  in Theorem~\ref{thm:gd_pl_&_ac} and Theorem~\ref{thm:pl_accel} (ii) are
\[
\frac{\mu}{L},\qquad 
a\frac{\sqrt{\mu \mu_0}}{L},\qquad 
a\!\left(\frac{\mu_0}{L_0}\right)^{1/4}\!\sqrt{\frac{\mu}{L}}(1-c_1).
\]
Ignoring the factor $(1-c_1)$, the latter always dominates $a\frac{\sqrt{\mu\mu_0}}{L}$, such that we compare $\mu/L$ and $a(\mu_0/L_0)^{1/4}\sqrt{\mu/L}$. 
% Thus, the rate of \eqref{eq:gd} for functions in $\pl\cap\ls$ is worse than the rate of \eqref{eq:nm} for functions in $\pl\cap\ac\cap\ls$ whenever
Thus, for functions in $\pl\cap\ac\cap\ls$, the rate of  \eqref{eq:nm} is better than the rate of \eqref{eq:gd} whenever
\[
a \;\ge\; \left(\frac{L_0}{\mu_0}\right)^{1/4}\sqrt{\frac{\mu}{L}}.
\]
This bound suggests that, in the presence of oscillations and flat plateaus, leading to $\mu/L \ll 1$, there is increased tolerance for poor alignment while still allowing acceleration.
%This bound indicates that in the presence of flat plateaus, inducing $\mu \ll 1$, there is more room for aiming condition to be poor and having acceleration. 
We test this bound numerically on the PL lower-bound function of \cite{PLlowerbound}, and find that the two quantities scale similarly, in a sense described in Appendix~\ref{app:heuristic}, providing a heuristic support for the relevance of our bound.
% If $\frac{\mu}{L}$ is negligible relative to $\frac{\mu_0}{L_0}$, a rough guideline to have a theoretical benefit of momentum is that $a$ should exceed the square root of the condition number $\mu/L$. However in the case where $\mu_0/L_0 \approx \mu/L$, we should have $a \ge (\mu/L)^{1/4}$.
% This comparison indicates that poorly conditioned problem 
% This quantify the role of the aiming condition on the benefit of momentum. \jh{nul}



\paragraph{Acceleration for Dynamical Systems: Flatness Matters}
Comparing the rates of Theorem~\ref{thm:gf_pl_&_ac} (i) and Theorem~\ref{thm:pl_accel} (i) yields an improvement of momentum if $a \ge (L_0/\mu_0)^{1/4}\sqrt{\mu}$, namely a similar bound as a in the discrete case, up to the $\sqrt{L}^{-1}$ factor. Contrasting with the discrete case, the rate of Theorem~\ref{thm:pl_accel} (i) does not necessarily improve upon Theorem~\ref{thm:gf_pl_&_ac} (ii). Improvement occurs whenever
\[
\mu_0 L_0 < 1.
\]
Since $L_0 \ge \mu_0$, this implies that if $\mu_0>1$ (the function is sharp near the minimizer), then \eqref{eq:gf} is preferable, whereas if $L_0<1$ (flatness near minimizer), then \eqref{eq:nmo} performs better.
In the intermediate regime $\mu_0 < 1 < L_0$, the outcome depends on the balance between the smallest and largest local curvatures: if $\mu_0$ is much smaller than~$1$ while $L_0$ is only moderately above~$1$, flatness dominates and \eqref{eq:nmo} wins; conversely, in the opposite case.
The advantage of \eqref{eq:nmo} on flat landscapes is consistent with previous observations for the Heavy Ball dynamics when $f\in\pl$ \cite{apidopoulos2022convergence, gess2023exponential}, whose dynamics are close to \eqref{eq:nmo}.
%We verify this behavior empirically. \jh{fig} Note that this intuitive property does not appear when considering rates associated to $\sqc$.


\subsection{Decoupling information: more interpretability, but may create sub-optimality.} 
%From Lemma~\ref{lem:scalar_prod_bound}, it follows that we do not really introduce and study a new class of function. We rather introduce a new prism to study the strongly quasar convex-like functions, precisely by decoupling information: angle with $\ac$, and curvature with $\pl$. This gives more intuition/interpretation, precising the role of each, where the information is more blurrily blended with $\sqc$. This comes at a cost: this may create sub-optimality. This follows from the definition of the constant parameterizing the class. Define $a(x) = \dotprod{\nabla f(x),x-x^\ast}/({\norm{\nabla f(x)}\norm{x-x^\ast}})$, and assume the worst-case point $\hat{x} = \arginf_{x \in \R^d} a(x)$ is such that $a(\hat x) > 0$. Then, we have that $f \in \ac$ for $a = a(\hat x)$. For the PL condition, the $\mu$ constant is also determined by such a worst-case point. However, it is likely that it will not be the same as for the aiming condition. When an assumption blends the two information in one sole inequality, the associated worst case point balance between each condition. This limitation is rather artificial, because we use the condition at the same points. Precisely, the aiming condition will be used in our convergence proofs through Lemma~\ref{lem:scalar_prod_bound}. But its proof remain valid for $a$ defined as  \[a := \frac{\dotprod{\nabla f(\hat x),\hat x-x^\ast}}{\norm{\hat x-x^\ast}\norm{\nabla f(\hat x)}},\]  where
%  \[\hat{x} :=\argmin_{x \in \R^d }\frac{\dotprod{\nabla f(x),x-x^\ast}}{\norm{x-x^\ast}\norm{\nabla f(x)}} \cdot \frac{\norm{\nabla f(x)}}{\sqrt{f(x)-f^\ast}}\]
%  This gives a worst-case point for a compromise between $\ac$ and $\pl$.

%  % and
%  % \[\mu := \frac{1}{2}\frac{\norm{\nabla f(\hat x)}}{f(\hat x)-f^\ast}.\] 

%  As mentioned in Section~\ref{sec:related_aiming}, \cite[Theorem 7]{hermant2024study} derives a different parameterization. The parameters they obtain are such that the bounds of \eqref{eq:nmo} and \eqref{eq:nm} reduce to those of \eqref{eq:gf} and \eqref{eq:gd}, see Appendix~\ref{app:related:uaac}. 
From Lemma~\ref{lem:scalar_prod_bound}, it follows that we do not study in this Section~\ref{sec:cv_results} a class of function that is fundamentally different from $\sqc$. We rather introduce a new characterization of these kind of functions in the idea of having better interpretability. This is done by decoupling the information: angular information with $\ac$, and curvature with $\pl$. In contrast, all these information are blended with $\sqc$, and according to Section~\ref{sec:sqc_pitfalls}, in a fairly hidden way.

This separation of information comes at a cost: this may create sub-optimality. A clear example is that for a quadratic function in $\strongconv \cap \ls$, one may have $a=2\sqrt{\mu L}/(\mu+L)$. Plugged into our Table bounds, this creates significantly poorer bounds compared with the known bounds, see Section~\ref{table:sc_sqc_rates}. This sub-optimality follows from the definition of the constants parameterizing the class, that follow from \textit{worst-case points}. Define $a(x) = \dotprod{\nabla f(x),x-x^\ast}/({\norm{\nabla f(x)}\norm{x-x^\ast}})$, and assume the worst-case point $\hat{x} = \arginf_{x \in \R^d} a(x)$ is such that $a(\hat x) > 0$. Then, we have that $f \in \ac$ for $a = a(\hat x)$. For the PL condition, the $\mu$ constant is also determined by such a worst-case point. However, the worst case points might not be the same for the two conditions; this is exactly what happens for a strongly convex quadratic function, where the worst aiming condition point is reached where the curvature is high. When an assumption blends the two information in a unique inequality, the associated worst case point balance between each condition. For such rigid functions such as quadratic, the different information (angular, curvature) are strongly coupled, such that a unique inequality can encompasses efficiently all the information. In a non-convex case these information may be decoupled. This may be seen with our toy example  in Section~\ref{sec:aiming_cond}, where we can fix $a=1$ while playing with the curvature through $f$. This indicates that such a decoupling can bring interesting insight for non-convex functions, compared to the more traditional $\sqc$ based on a unique inequality.    




Anyway, we note that the aforementioned limitation is rather artificial. A first reason is that it can be corrected if allowing a more involved definition of the classes. Precisely, the aiming condition will be used in our convergence proofs through Lemma~\ref{lem:scalar_prod_bound}. Its proof remains valid for $a$ defined as  \[a := \frac{\dotprod{\nabla f(\hat x),\hat x-x^\ast}}{\norm{\hat x-x^\ast}\norm{\nabla f(\hat x)}},\]  where
 \[\hat{x} :=\argmin_{x \in \R^d }\frac{\dotprod{\nabla f(x),x-x^\ast}}{\norm{x-x^\ast}\norm{\nabla f(x)}} \cdot \frac{\norm{\nabla f(x)}}{\sqrt{f(x)-f^\ast}}.\]
 Namely, the aiming condition constant may be considered at a point that creates a balance with the PL condition. This allows to keep the information of the aiming condition, while not inducing an artificial sub-optimality. A second justification is provided in the next section.

% \subsection{Dynamical 
% Systems}\label{sec:dynamical_systems}
% In this section, we focus on the dynamical systems \eqref{eq:gf} and \eqref{eq:nmo}. We start with \eqref{eq:gf}, and we compare the results obtained depending on $f$ belongs to $\ac$ or not.
% \begin{theorem}\label{thm:gf_pl_&_ac}
%     Let $\xcont \sim \eqref{eq:gf}$.
    
%     (i) If $f \in \pl$,  then
%     $$ f(x_t)-f^\ast \le (f(x_0)-f^\ast)e^{-2 \mu t}. $$
%     (ii) if $f \in \pl \cap \ac \cap \qg$, then
%     $$ f(x_t)-f^\ast \le \frac{\sqrt{L_0}(\sqrt{L_0}+\sqrt{\mu_0})}{2}\norm{x_0-x^\ast}^2e^{-a\sqrt{\mu \mu_0}t}, $$
%     where $\mu_0 := \sup_{\mu' \ge \mu}\{f \in \qginff{\mu'}\} $.
% \end{theorem}
% See the proof in Appendix~\ref{app:gf_pl_&_ac}. For point (ii), we also assume $f \in  \qg$ for technical reasons. In the subsequent section we will consider algorithms minimizing functions belonging to $\ls$, such that they will naturally belong to $\qg$.
% The bound (ii) in Theorem~\ref{thm:gf_pl_&_ac} indicates that if belongs to $\ac$ with $a$ close to $1$, \textit{e.g.} for $1$-dimensional functions, the \hyperlink{cr}{convergence rates} becomes $\bigO(\sqrt{\mu \mu_0})$, recalling $\mu_0 \ge \mu$. In particular we may have $\mu_0 \gg \mu$, such that it can be significantly better than $\mu$, see Section~\ref{sec:relax_setting}. The fact that $\mu_0$ appears in the convergence rate of (ii) has a deeper interpretation.
 
%  \noindent\textbf{The aiming condition exploits the local information at $x^\ast$ \:} For functions in $\strongconv \cap \ls$, the slowest direction toward the minimizer is the one with smallest curvature~$\mu$, since it corresponds to the flattest direction (left plot of Figure~\ref{fig:worst_case}).  
% For functions in $\pl \cap \ls$, the situation is less clear: the trajectory may approach the minimizer either through low-curvature directions near $x^\ast$ or through plateau-like regions (right plot of Figure~\ref{fig:worst_case}). Thus, it is ambiguous whether the relevant $\mu$ is determined by local curvature at $x^\ast$ or by the plateaus, as $\pl$ handles both informations.
% Bound~(i) of Thm.~\ref{thm:gd_pl_&_ac} does not separate these effects, but bound~(ii) does: it reveals a trade-off between global plateaus and local curvature~$\mu_0$. If $\mu=\mu_0$, flatness near $x^\ast$ dominates and we recover the rate $\mathcal{O}(e^{-\mu t})$. If $\mu<\mu_0$, the plateaus matter, yielding a rate $e^{-\sqrt{\mu\mu_0}}$, which shows that plateaus do not fully determine the convergence.

%\jh{(}Technically, making $\mu_0$ appear under the sole assumption $f\in\pl$ seems not possible. A possible intuition is that without $\ac$, we lose control on when the trajectory enters the locally curved region around the minimizer.\jh{)}

%  For functions in $\strongconv \cap \ls$, the worst direction along which approaching a minimizer is the direction of lowest curvature, characterized by $\mu$, because it correspond to the direction where the landscape is the flattest, see left plot of Figure~\ref{fig:worst_case}.
% For functions belonging to $\pl \cap \ls$, the intuition is less straightforward. The trajectory could also approach the minimizer along a direction that incorporate plateaus, see the right plot of Figure~\ref{fig:worst_case}. This means that it is not clear if the optimal $\mu$ such that $f \in \pl$ is determined by these plateaus or by the direction with lowest curvature at $x^\ast$, because $\pl$ handles both information. 
% %So, for a given function, the optimal $\mu$ constant such that it belongs to $\pl$ could be determined by the flat plateaus, or by a flat quadratic lower bound around minimizers.
% While bound (i) of Theorem~\ref{thm:gd_pl_&_ac} does not separate the two scenarios, it is the case with bound (ii), whose bound can be interpreted as a trade-off between the flatness of the plateaus, and the flatness around the minimizer $x^\ast$. Precisely, if $\mu = \mu_0$, the most important flatness of the function is around the minimizer, and we recover the rate $\bigO(e^{-\mu t})$. However if $\mu < \mu_0$, then intuitively the flatness of the plateaus prevail. The fact that we get a rate $e^{-\sqrt{\mu \mu_0}}$ instead of $e^{-\mu t}$ indicates that these plateaus does not totally determine the convergence rate, which also relies on $\mu_0$. So, taking into account the information of $\ac$ allows to separate these information.
% From a technical point of view, making appear $\mu_0$ when we only have $f \in \pl$ seems not possible. Intuitively, it can be explained by the fact that $f \notin \ac$ means that we have less theoretical control over the time the trajectory enters the local flatness around the minimizer.





% We now state a convergence result for the momentum ODE. 
% \begin{theorem}\label{thm:nmo_placgq}
% Let $f\in \pl \cap \ac \cap \qg$. Let $\xcont \sim \eqref{eq:nmo}$ with  constant parameters $\eta_t = (\mu_0 L_0)^{1/4}$, $\eta'_t = a\bpar{\frac{\mu_0}{L_0}}^{\frac{1}{4}}\sqrt{\mu}$, $\gamma'_t = (\mu_0 L_0)^{-1/4}$ and any $\gamma_t\ge 0$. Then,
%      \begin{equation}
%      f(x_t)-f^\ast \le K_0 e^{-a \bpar{\frac{\mu_0}{L_0}}^{\frac{1}{4}}\sqrt{\mu}t},
%  \end{equation}
%  where $K_0 := \frac{\sqrt{L_0}(\sqrt{L_0}+\sqrt{\mu_0})}{2}\norm{x_0-x^\ast}^2$ and where $\mu_0 := \arg\sup_{\mu' \ge \mu}\{f \in \qginff{\mu'}\} $.
% \end{theorem}
% The rate of Theorem~\ref{thm:nmo_placgq} improves upon (ii) in Theorem~\ref{thm:gf_pl_&_ac} whenever
% \[
% \mu_0 L_0 < 1.
% \]
% Since $L_0 \ge \mu_0$, this implies that if $\mu_0>1$ (the function is sharp near the minimizer), then \eqref{eq:gf} is preferable, whereas if $L_0<1$ (flatness near minimizer), then \eqref{eq:nmo} performs better.
% In the intermediate regime $\mu_0 < 1 < L_0$, the outcome depends on the balance between the smallest and largest local curvatures: if $\mu_0$ is much smaller than~$1$ while $L_0$ is only moderately above~$1$, flatness dominates and \eqref{eq:nmo} wins; conversely in the opposite case.
% The advantage of \eqref{eq:nmo} on flat landscapes is consistent with previous observations for the Heavy Ball dynamics when $f\in\pl$ \cite{apidopoulos2022convergence, gess2023exponential}. We verify this behavior empirically. \jh{fig} Note that this intuitive property does not appear when considering rates associated to $\sqc$.


% The rate of Theorem~\ref{thm:nmo_placgq} is better than (ii) in Theorem~\ref{thm:gf_pl_&_ac} as long as
% $$ \mu_0 L_0  < 1. $$
% Recalling $L_0\ge \mu_0$, we observe that if $\mu_0 > 1$, or the function is sharp around its minimizer, then \eqref{eq:gf} is better. If $L_0 <1$, the function is flat around its minimizer, and \eqref{eq:nmo} is better. If we have $\mu_0 < 1 < L_0$, there is a trade-off between the lowest and the largest curvature. If the $\mu_0$ is significantly lower than $1$ than $L_0$ is higher than $1$, then the "flatness" wins and $\eqref{eq:nmo}$ is better. Conversely in the opposite case.  The better performance of \eqref{eq:nmo} in flat landscapes is consistent with former observations for the Heavy Ball equation when $f \in \pl$ \cite{apidopoulos2022convergence}. 






% \subsection{Algorithms}\label{sec:accel_mom}
% In this section, we derive results for the algorithms \eqref{eq:gd} and \eqref{eq:nm}, assuming the functions belong to $\ls$. Classically, we recover similar convergence rates to those of the dynamical-system setting (Section~\ref{sec:dynamical_systems}), with the difference that the $L$-smoothness constant intervenes, introducing the conditioning ratio $\mu/L$.
% \begin{theorem}\label{thm:gd_pl_&_ac}
%     Let $f \in \ls$, $\xdisctilde \sim \eqref{eq:gd}$, with $\gamma = \frac{1}{L}$.
    
%     (i) If we also have $f \in \pl$,  then
%     $$ f(\tilde x_k)-f^\ast \le \bpar{1-\frac{\mu}{L}}^k(f(\tilde x_0)-f^\ast). $$
%     (ii) If we also have $f \in \pl \cap \ac $, then
%     $$ f(\tilde x_k)-f^\ast \le  K_0\bpar{1-a\frac{\sqrt{\mu \mu_0}}{L}}^k\norm{\tilde x_0-x^\ast}^2. $$
%     with $K_0 := \frac{\sqrt{L_0}(\sqrt{L_0}+\sqrt{\mu_0})}{2}$, $\mu_0 := \sup\{\mu' \ge \mu: f \in \qginff{\mu'}\} $ and $L_0 :=\inf\{L' \le  L:f \in \qgf{L'}\}$.
% \end{theorem}
% Grossly, the convergence rate of Theorem~\ref{thm:gd_pl_&_ac} are the same than Theorem~\ref{thm:gf_pl_&_ac}, divided by $L$. Dividing by $L$ does not affect the comments under Theorem~\ref{thm:gf_pl_&_ac}. We will take a more careful look at the result obtained for \eqref{eq:nm}.

% \begin{theorem}\label{thm:pl_accel}
% Let $f\in \pl \cap \ac \cap \ls$.  Denote $L_0 := \inf \{L'\le L:  f \in \qgf{L'} \}$, $\mu_0 := \sup\{\mu' \ge \mu : f \in \qginff{\mu'}\}$.
% Let $\xdisctilde$ verify \eqref{eq:nm} with the \hyperlink{cont_param}{continuized parameterization}, with 
% $\eta = \frac{(\mu_0 L_0)^{1/4}}{\sqrt{L}},\quad \gamma = \frac{1}{L},\quad \gamma' = \frac{1}{( \mu_0 L_0)^{1/4}\sqrt{ L}}, $ and $\eta' = a\bpar{\frac{\mu_0}{L_0}}^{1/4}\sqrt{\frac{\mu}{L}}$. Then,  
%     with probability  $1-\frac{1}{c_0} - e^{-(c_1-1-\log(c_1))k}$, for some $c_0 > 1$ and $c_1 \in (0,1)$ we have
% $$ f(\tilde x_k)-f^\ast \le K_1e^{-a\bpar{\frac{\mu_0}{L_0}}^{1/4}\sqrt{\frac{\mu}{L}}(1-c_1)k}, $$
% where $K_1 := c_0\frac{\sqrt{L_0}(\sqrt{L_0}-\sqrt{\mu_0})}{2}\norm{\tilde x_0-x^ \ast}^2$.
% \end{theorem}

% See the proof in Appendix~\ref{app:proof:pl_accel}. The result of Theorem~\ref{thm:pl_accel} holds under high probability, due to the use of the \textit{continuized} version of Nesterov. It is not clear whether this result can be obtained with a more classical parameterization of \eqref{eq:nm}.

% \noindent
% \textbf{Comparison of the convergence rates} The \hyperlink{cr}{convergence rates}  in Theorem~\ref{thm:gd_pl_&_ac} and Theorem~\ref{thm:pl_accel} are
% \[
% \frac{\mu}{L},\qquad 
% a\frac{\sqrt{\mu \mu_0}}{L},\qquad 
% a\!\left(\frac{\mu_0}{L_0}\right)^{1/4}\!\sqrt{\frac{\mu}{L}}(1-c_1).
% \]
% Ignoring the factor $(1-c_1)$, the latter always dominates $a\frac{\sqrt{\mu\mu_0}}{L}$. 
% % Thus, the rate of \eqref{eq:gd} for functions in $\pl\cap\ls$ is worse than the rate of \eqref{eq:nm} for functions in $\pl\cap\ac\cap\ls$ whenever
% Thus, for functions in $\pl\cap\ac\cap\ls$, the rate of  \eqref{eq:nm} is better than the rate of \eqref{eq:gd} whenever
% \[
% a \;\ge\; \left(\frac{L_0}{\mu_0}\right)^{1/4}\sqrt{\frac{\mu}{L}}.
% \]
% If $\frac{\mu}{L}$ is negligible relative to $\frac{\mu_0}{L_0}$, a rough guideline to have a theoretical benefit of momentum is that $a$ should exceed the square root of the condition number $\mu/L$. However in the case where $\mu_0/L_0 \approx \mu/L$, we should have $a \ge (\mu/L)^{3/4}$.
% Finally, compared to the discussion in Section~\ref{sec:sqc_pitfalls}, this viewpoint gives a more clear intuition about what makes momentum improve the convergence rate.


%  The \hyperlink{cr}{convergence rates} of Theorem~\ref{thm:gd_pl_&_ac} (i), (ii) and Theorem~\ref{thm:pl_accel} are respectively
% $$ \frac{\mu}{L},\quad a\frac{\sqrt{\mu \mu_0}}{L}, \quad a \bpar{\frac{\mu_0}{L_0}}^{1/4}\sqrt{\frac{\mu}{L}}(1-c_1). $$
% Discarding the factor $(1-c_1)$ (can be taken as $\frac{1}{2}$), we always have that $a\bpar{\frac{\mu_0}{L_0}}^{1/4}\sqrt{\frac{\mu}{L}} \ge  a\frac{\sqrt{\mu \mu_0}}{L}$. So, we compare the rate of \eqref{eq:gd} for functions in $\pl \cap \ls$, to the rate of \eqref{eq:nm} for functions in $\pl \cap \ac \cap \ls$, and the latter is better as long as
% $$ a \ge \bpar{\frac{L_0}{\mu_0}}^{1/4}\sqrt{\frac{\mu}{L}}. $$
% In many settings, $\frac{\mu}{L}$ is negligible relative to $\frac{\mu_0}{L_0}$. Hence, as a rough guideline, $a$ must exceed the square root of the condition number $\mu/L$.
% Finally, we note that compared to the last issue addressed in Section~\ref{sec:sqc_pitfalls} when comparing $\sqc$ and $\pl$ rates, our viewpoint offer a a clearer intuition wether momentum can improve or not.



% \jh{ce qui suit : trop long et nul. Comparer directement les rate en PL only et nest. Exhiber condition sur correlation}
% We have $e^{-\alpha k} \approx (1-\alpha)^k$, for $\alpha<1$ and $k$ not very small. In this regard, the rate of Theorem~\ref{thm:pl_accel} is approximately
% $$ f(\tilde x_k)-f^\ast \le K_0\bpar{1- (1-c_1)\frac{a}{\rho}\bpar{\frac{\mu_0}{L_0}}^{1/4}\sqrt{\frac{\mu}{L}}}^k.$$
% We see that appears the $\sqrt{\frac{\mu}{L}}$ that characterize strongly convex acceleration. The other most important factor is $a$, the correlation \jh{nom}. We have the bound $\frac{\mu_0}{L_0} \ge \frac{\mu}{L}$, which would make appear $\bpar{\frac{\mu}{L}}^{-3/4}$, which is still better than $\frac{\mu}{L}$, but weaker than $\sqrt{\frac{\mu}{L}}$. However, in any non-convex scenarios, we have that the inequality is strict, and in many situations we have $\frac{\mu_0}{L_0} \gg \frac{\mu}{L}$. 

% \subsection{Link with $\sqc$}
% In previous sections, we introduced $ \ac$ in order to address concerns related to $\sqc$, see Section~\ref{sec:sqc_pitfalls}.
% The following result draws a connection between these classes. 



\section{Relaxed Setting: Aiming Condition on Average}\label{sec:avg}
There is a gap between the convergence observed in practice and the theoretical bounds established for globally defined classes such as PL functions. 
%\jh{exemple de fonction (ex fig 2) où on regarde cv avec constantes theoriueq, et plus agressives ?}
The reason is intuitive: if the global landscape of non-convex functions can be highly irregular, nevertheless, the optimization path may remain in a favorable region. Importantly, the performance of first-order algorithms such as \eqref{eq:gd} and \eqref{eq:nm} only depends on the geometry surrounding their optimization path.

\noindent
\textbf{Naive Relaxation \: }A first direct idea to take this intuition into account is to consider an algorithm $\xdisc$ and a minimizer $\hat x$ such that $\tilde x_k \to_{k \to +\infty} \hat x$, and assuming
\begin{equation}\label{eq:naive_relax}
  \forall k\in \N,\dotprod{\nabla f(\tilde x_k),\tilde x_k-\hat x} \ge a \norm{\nabla f(\tilde x_k)}\norm{\tilde x_k-\hat x}.
\end{equation}
With words, $f$ is assumed to satisfy the $a$-aiming condition along the trajectory $\xdisc$, with respect to its convergence point $\hat x$. For a function in $\pl \cap \ls$, if $\xdisc \sim \eqref{eq:gd}$ or \eqref{eq:nm} verify \eqref{eq:naive_relax}, we would deduce for this sequence the same result as for Section~\ref{sec:cv_results}. 
The relaxation is twofold: (i) we get rid of the uniqueness minimizer assumption, and (ii) the convergence rate is not affected by potential pathological regions of $f$, as long as the sequence does not cross it. However, it suffices for \eqref{eq:naive_relax} to not be verified for a single iteration to make this reasoning fail. We thus derive a relaxation that, in the case where pathological regions remain a negligible section of the optimization path, still allows to keep accelerated results. 
%\jh{dire quelquepart qu'on suppose que la traj convergence, mais qu'en réalité on tune les pas pour bien converger, donc ça a du sens de considérer ce cas de figure}

\noindent
\textbf{Averaging Relaxation \: } We relax the viewpoint \eqref{eq:naive_relax}. We address the continuous case; the discrete case follows similarly, see Appendix~\ref{app:adptative_condition}.
%Looking at the right plot of Figure~\ref{fig:pl_sqc_low_corr}, it is clear that even if there exists regions of the landscape that make $f$ belong to $\ac$ with $a$ very small, a trajectory will quickly fall in a zone such that the aiming condition is actually closer to one. More generally, non-convexity makes possible for function to have pathological regions that, while degrading the geometrical parameters, remain a negligible section of the optimization path. We can formalize this intuition with trajectory-depending convergence rates.
%It can be challenging to design functions that belongs to structured class such as $\pl \cap \ac \cap \qg$, such that all the optimization path lies in a pathological region, which could be related to worst-case functions of this class for a given algorithm. Actually, for "natural" functions arising in practical applications, it seems reasonable to think it will globally lie in a reasonable region with small pathologic \jh{pt cite "vizualisation loss" avec le convex autour minim}.

\begin{assumption}\label{ass:avg_corelation}
    For some minimizer $x^\ast \in \R^d$, $\theta > 0$, $a_{avg} > 0$, and $\xcont$, we have
    \begin{align*}
        \int_0^t e^{\theta s}\Delta_sds
        \ge 0,
    \end{align*}
   $\Delta_s :=\dotprod{\nabla f(x_s),x_s-x^\ast}-a_{avg}\norm{\nabla f(x_s)}\norm{x_s-x^\ast}$.
\end{assumption}
A function $f$ satisfies Assumption~\ref{ass:avg_corelation} as long as the aiming condition holds along a specific average of the trajectory.
If $f \in \ac$, Assumption~\ref{ass:avg_corelation} holds with $a_{avg} = a$ for any %$\xcont$ and
initialization. However, we could have $a \ll a_{avg}$, or having that the path crosses some regions where the aiming condition is not verified, while still having $a_{avg} > 0$. %\jh{ref fig example}. %The $e^{\theta s}$ factor indicate that the poorness of the aiming condition in the early stage of the process does not significantly affect the constant.
\begin{theorem}\label{thm:gf_pl_&_avg_cor}
    Let $f \in \pl \cap \qg$, $\mu_0 := \sup\{\mu' \ge \mu: f \in \qginff{\mu'}\} $ and $a_{avg}>0$. If:
    
    (i) $\xcont \sim \eqref{eq:gf}$ for a given initialization $x_0 \in \R^d$ is such that Assumption~\ref{ass:avg_corelation} holds with $\theta = a_{avg}\sqrt{\mu \mu_0}$. , then
    $$ f(x_t)-f^\ast \le K_0e^{-a_{avg}\sqrt{\mu \mu_0}t}(f(x_0)-f^\ast). $$
    
    (ii) $\xcont \sim \eqref{eq:nmo}$ for $\eta' = a_{avg}\bpar{\mu_0/L_0}^{1/4}\sqrt{\mu}, \gamma' =  (\mu_0 L_0)^{-1/4},~\eta = (\mu_0 L_0)^{1/4}, ~ \gamma \ge 0$ and such that 
Assumption~\ref{ass:avg_corelation} holds with $\theta =  a_{avg}\bpar{\mu_0/L_0}^{1/4}\sqrt{\mu}$,  then
    $$f(x_t)-f^\ast \le K_0e^{- a_{avg}\bpar{\frac{\mu_0}{L_0}}^{\frac{1}{4}}\sqrt{\mu}} (f(x_0)-f^\ast).$$
    In both statements, $K_0 :=  \bpar{1+\sqrt\frac{L_0}{\mu_0}}.$
\end{theorem}
The proof is in Appendix~\ref{app:gf_pl_&_avg_cor}. We obtain the same convergence rate as for Theorem~\ref{thm:gf_pl_&_ac} (ii) and Theorem~\ref{thm:pl_accel} (i), replacing $a$ with $a_{avg}$. However, the rates of Theorem~\ref{thm:gf_pl_&_avg_cor} could be either significantly sharper, or allow for obtaining such a rate while the algorithm crosses regions with negative aiming condition values. This reasoning can be extended to other geometrical conditions, such as $\pl$. This perspective might explain why, in practice, grid-search selection of parameters often leads to relatively large constants, compared to what would be allowed by geometrical constants that parameterize globally the associated landscape.
%This formalizes that the algorithms are not sensitive to negligible pathological regions.




  
\section{Numerical Experiment}\label{sec:numerical}
We provide numerical experiments to support our theoretical findings. We first design a simple $2$d example that belongs to $\pl \cap \ls$, has a unique minimizer, but does not belong to $\ac$. This will result in momentum hurting the convergence speed in the early stage of optimization, see Section~\ref{sec:numerical_1}. Then, in Section~\ref{app:heuristic}, we provide a heuristic validation of the bound to have acceleration with \eqref{eq:nm} provided in Section~\ref{sec:mom_cv}.
%\jh{Peut être en première "xp", ou alors à mettre vers fin section 4, on peut reprendre figure 2 et regarder a et l'autre borne.}

%\jh{there is a clear intuition that aiming condition will benefit momentum method. Can we design functions where no aiming condition actually hurt convergence ?}



% \begin{figure}
%     \centering
%     \includegraphics[width=0.51\linewidth]{figures/toy_example_fig.pdf}
%     \hspace{-7pt}
%     \includegraphics[width=0.5\linewidth]{figures/toy_example_traj.pdf}
%         \includegraphics[width=0.48\linewidth]{figures/toy_example_fonc.pdf}
%                 \includegraphics[width=0.48\linewidth]{figures/toy_example_corr.pdf}
%     \caption{ Upper left: heatmap of $f(x,y) = 0.5(y-\sin(x))^2 + 0.001\cdot0.5x^2$. Blue arrows are the descent direction $-\nabla f$, the red cross is the unique minimizer at $(0,0)$. Upper right: plot of the first $1000$ iterations of the trajectories of \eqref{eq:gd} and \eqref{eq:nm_prime} on $f$.  Bottom left: Corresponding decrease of $\log(f)$. Bottom right: value of the aiming condition along iterates. We observe that the descent directions does not necessarily point toward the minimizer, which cause momentum to accumulate in the wrong direction, causing \eqref{eq:gd} to be ahead in term of decrease speed in the first $300$ iterations. }
%     \label{fig:toy_prob}
% \end{figure}
%\jh{fleche to minimizer + graphes vraiment petits}
\subsection{An Example with Negative Aiming Condition}\label{sec:numerical_1}
In this section, we design an example that shows how negative aiming condition values can hurt the convergence of momentum algorithms, and then discuss the aiming condition in neural networks. For our experiment, we consider an alternative formulation of \eqref{eq:nm} that we call \eqref{eq:nm_prime}.
\begin{equation}\label{eq:nm_prime}\tag{NM'}
    \begin{array}{c}
\left\{
\begin{aligned}
&\tilde y_k = \tilde x_k + \alpha(\tilde x_k -\tilde x_{k-1})  \\
&\tilde x_{k+1} =\tilde y_k - \gamma\, \nabla f(\tilde y_k)
\end{aligned}
\right.
\end{array}
\end{equation}
While \eqref{eq:nm} with the continuized parameterization is convenient for theoretical analysis, \eqref{eq:nm_prime} is more amenable to practical implementation, as it involves only two easily tunable parameters with clear interpretations: the stepsize $s$ and the momentum parameter $\alpha$. See \cite{defazio2019curved} for a discussion on the different forms of Nesterov momentum.

Consider the 2-dimensional function\[F_{\varepsilon}(x,y) = 0.5(y-\sin(x))^2 + 0.5\varepsilon x^2.\]
The special case $\varepsilon = 0$ is introduced in \cite{apidopoulos2022convergence}. %We have $F_{0} \in \plf{2}$, and its minimizers form a valley following the graph of $(x,\sin(x))$, namely $\arg \min F_0 = \{(x,y)\in \R^2, ~ y = \sin(x)\}$.
For $\varepsilon > 0$, $F_{\varepsilon}$ has $(0,0)$ as unique minimizer, and it belongs to $\pl$ with $\mu \le \frac{2+\varepsilon
-\sqrt{\varepsilon^2+4}}{2}$, see Lemma~\ref{lem:toy_example_pl}. We display $F_{0.001}$ in Figure~\ref{fig:toy_prob}. We also provide a three-dimensional visualization in the appendix, see Figure~\ref{fig:3d_vizualization}. %\jh{appendix 3d vizual}.
At certain points of the landscape, the descent direction is anti-correlated with the direction pointing toward the minimizer, thereby violating the aiming condition.
%In the same time, because there is a unique minimizer--and unique critical point--, we are ensured that momentum will not push the trajectory toward a different solution than \eqref{eq:gd}, such that we can isolate the problem of aiming condition.\jh{nul}
\begin{figure}[h]
    \centering
    % Colonne gauche (2 figures côte à côte)
    \begin{minipage}{0.65\linewidth}
        \centering
        \begin{minipage}{0.48\linewidth}
            \centering
            \includegraphics[width=\linewidth]{figures/toy_example_fig.pdf}
        \end{minipage}
        \begin{minipage}{0.48\linewidth}
            \centering
            \includegraphics[width=\linewidth]{figures/toy_example_traj.pdf}
        \end{minipage}
    \end{minipage}
    % Colonne droite (2 figures empilées)
    \begin{minipage}{0.3\linewidth}
        \centering
        \includegraphics[width=\linewidth]{figures/toy_example_fonc.pdf}

        \vspace{0.3cm}

        \includegraphics[width=\linewidth]{figures/toy_example_corr.pdf}
    \end{minipage}

  \caption{ Left: Heatmap of $F_{0.001}(x,y) = 0.5(y-\sin(x))^2 + 0.001\cdot0.5x^2$. Blue arrows indicate the descent direction $-\nabla F_{0.001}$. Center: First $1000$ iterations of the trajectories of \eqref{eq:gd} and \eqref{eq:nm_prime}, both starting from the initialization point $(0,3)$, for different values of momentum parameter $\alpha$.  Top right: Corresponding decrease of $\log(f)$. Bottom right: Values of the aiming condition along the iterates, zoomed on the first 100 iterations. Early negative aiming condition values cause momentum to drive the trajectory away from the minimizer, leading \eqref{eq:gd} to outperform momentum in early iterations, with increasing effect for larger $\alpha$.}
    \label{fig:toy_prob}
\end{figure}

\noindent
\textbf{Negative aiming condition shoots \eqref{eq:nm_prime} away \:} In Figure~\ref{fig:toy_prob}, we compare \eqref{eq:gd} and \eqref{eq:nm_prime} for different values of momentum parameter $\alpha$, see details in Appendix~\ref{app:toy_prob}. Due to negative aiming condition values during the early phase of the optimization process, iterates of \eqref{eq:nm_prime} are initially driven away from the minimizer. Increasing the momentum parameter amplifies this effect, causing the trajectory to deviate even further. 
%\jh{gradient steps drive the trajectory away from the minimizer, which is amplified by momentum} iterates of \eqref{eq:nm_prime} are initially driven away from the minimizer. This gets even more pronounced as the momentum parameter $\alpha$ is taken larger.
As a result, \eqref{eq:gd} outperforms \eqref{eq:nm_prime} for up to $1000$ iterations in terms of function value decrease. This example illustrates how negative aiming condition values can hinder momentum methods. 
More broadly, designing functions with unfavorable aiming behavior while retaining desirable optimization properties—such as satisfying the PL condition and having a unique minimizer—appears to be an interesting direction for better understanding landscapes on which momentum is ineffective.
% We believe that designing functions with bad aiming condition while retaining favorable optimization properties, \textit{e.g.} being $\pl$ or with unique minimizer, is an interesting question and could help understand which landscapes are unfavorable to momentum.

\paragraph{\eqref{eq:nm_prime} eventually catches up due to favorable averaged aiming condition \:}
Figure~\ref{fig:toy_prob} shows that, although the aiming condition takes negative values at early iterations, it then stabilizes around a positive value close to one. As discussed in Section~\ref{sec:avg}, this suggests that \eqref{eq:nm_prime} should eventually benefit from momentum for a large enough number of iterations. We confirm this behavior in Figure~\ref{fig:convergence_longer} in the appendix.

\paragraph{Aiming Condition and Neural Networks}

Because of the curse of dimensionality and their high level of non-convexity, we could expect that large-scale models such as those involved in deep learning do not verify any aiming condition. This intuition is challenged by the fact that well-designed neural networks are known to have favorable optimization properties. It was observed in \cite{li2018visualizing} that around 90$\%$ of the variations of the path of SGD with large batch size lie in a dimension 2 space, despite the high dimensionality of the models. Of special interest to our work, \cite{guille2024no} performs an extensive empirical evaluation of the aiming condition, which they name the \textit{cosine similarity}. They observe that for various tasks and architectures, $a$ remains between $10^{-1}$ and $10^{-3}$, despite the high dimensionality of the model. This observation was consistent across several stochastic first-order algorithms, including SGD with momentum. Although our theoretical results do not apply to general deep learning models, these experiments tend to highlight the crucial role of the aiming condition as a key ingredient in explaining the effectiveness of momentum.
% \cite{guille2024no} already explored empirically aiming condition. \cite{li2018visualizing} observed that around 90$\%$ of the variations of the path of SGD with large batch lies in a space of dimension 2. 
% \section*{Accessibility}
% Authors are kindly asked to make their submissions as accessible as possible for everyone including people with disabilities and sensory or neurological differences.
% Tips of how to achieve this and what to pay attention to will be provided on the conference website \url{http://icml.cc/}.
\subsection{Heuristic Evaluation of the Acceleration Bound}\label{app:heuristic}
In this section, we provide a heuristic validation of the bound on the aiming condition to provide acceleration, derived in Section~\ref{sec:mom_cv}. To do so, we consider the $\pl$ function introduced in \cite{PLlowerbound}, which serves to obtain a lower bound on the number of gradient call needed by first order algorithms--algorithms that make use of gradient information--to obtain a point $\hat x \in \R^d$ such that $f(\hat x)-\min f \le \varepsilon$. Precisely, it is shown that the number of needed gradient calls scales at best as $\bigO(\frac{L}{\mu}\log(\frac{1}{\varepsilon}))$ for any first-order algorithms. This bound corresponds to a $\mu/L$ convergence rate, inducing that, up to constants, the convergence rate of \eqref{eq:gd} is optimal, see Table~\ref{table:sc_sqc_rates}. The definition of the hard instance is rather involved see details in Appendix~\ref{app:hard_fonc_sqc}. 

\textbf{Non acceleration of momentum on the hard function \:}
We observe empirically that, as expected, momentum does not allow improvement over gradient descent when minimizing the hard function of \cite{PLlowerbound}. We consider \eqref{eq:gd} and \eqref{eq:nm} with the \hyperlink{cont_param}{continuized parameterization}. We chose a parameterization of the hard function such that it belongs to $\pl \cap \ls$ with $\mu=10^{-4}$, $L = 10^3$, and the dimension is $4550$ (see details in Appendix~\ref{app:hard_fonc_sqc}). For both \eqref{eq:gd} and \eqref{eq:nm}, we tune the parameters such that the decrease at iteration $1000$ is maximized. We plot the decrease of function values of the hard function in Figure~\ref{fig:hard_instance}. We observe that tuning of \eqref{eq:nm} only permits to catch up with the decrease of \eqref{eq:gd}, but provides no acceleration. 

\textbf{Heuristic test of the acceleration bound \:}
In Section~\ref{sec:mom_cv}, we stated that for $f \in \pl \cap \ac \cap \ls$, momentum provide acceleration as long as 
\[a \ge \bpar{\frac{L_0}{\mu_0}}^{1/4}\sqrt{\frac{\mu}{L}}.\]
It is natural to test whether this bound is verified or not on the aforementioned hard function. Importantly, the following discussion should be regarded as a \textbf{heuristic}, rather than a rigorous derivation.

For our test, we do not consider the global parameters $\mu, a ,L$ such that $f$ belongs to $\pl \cap \ac \cap \ls$. This is because (i) we do not know the optimal parameters $a, \mu_0$ and $L_0$, and (ii) in the lower bound provided in \cite{PLlowerbound}, there is a hidden constant with values $\approx 10^6$, which is not negligible. Moreover, the decrease observed in Figure~\ref{fig:hard_instance} is significantly faster than what is predicted by the theoretical bound derived under the global parameters, such that it actually poorly describes the algorithm's behavior. This motivates to use a perspective that has a similar spirit to the one described in Section~\ref{sec:avg}, such that we will consider the values of PL and aiming condition taken along the optimization path. Along the iterates of $\xdisc$ generated by \eqref{eq:gd} or \eqref{eq:nm}, we compute the point-wise PL and aiming condition values, namely
\[\mu(\tilde x_i) :=  2\frac{\norm{\nabla f(\tilde x_i)}^2}{f(\tilde x_i)-f^\ast}, \quad a(\tilde x_i) := \frac{\dotprod{\nabla f(\tilde x_i),\tilde x_i-x^\ast}}{\norm{\nabla f(\tilde x_i)}\norm{\tilde x_i-x^\ast}}.\] Empirically, we observe that $\mu(\tilde x_i) \approx 1$ (see Figure~\ref{fig:mu_corr_values}), while the function is $\pl$ with $\mu = 10^{-4}$. This suggests that the algorithm stays in favorable regions compared to those existing on the global landscape. As a result, in the case of \eqref{eq:gd}, compared to the rate $\mu/L$, we expect that its convergence rate is more accurately described by 
\[\tilde \mu_{gd}/L,\]

with $\tilde \mu_{gd} = \frac{1}{N}\sum_{i=1}^N \mu(\tilde x_i)$ is the averaged point-wise PL values along the $N$ iterations of the algorithm. Moreover, because we observed that \eqref{eq:gd} achieves a better decrease with stepsize $\gamma_{gd} \ge 1/L$, running the algorithm with this stepsize suggests that an even more accurate rate is
\begin{equation}\label{cvmodel_bound_gd}
    \tilde \mu_{gd}\gamma_{gd}.
\end{equation}
We plot the curve associated to this theoretical rate in Figure~\ref{fig:hard_instance}. We observe that, if not matching perfectly the decrease of \eqref{eq:gd}, it is significantly more accurate compared with the rate $\mu/L$.
Following the same reasoning, we will model the rate for \eqref{eq:nm} as (see Theorem~\ref{thm:pl_accel} (ii))
\begin{equation}\label{cvmodel_bound_nm}
    \tilde a (\tilde \mu_0/ \tilde L_0)^{1/4} \sqrt{\tilde \mu_{nm}\gamma_{nm}},
\end{equation}

where $\tilde a = \frac{1}{N}\sum_{i=1}^N a(\tilde x_i)$ and $\tilde \mu_{nm} = \frac{1}{N}\sum_{i=1}^N \mu(\tilde x_i)$ are respectively the point-wise aiming condition value and point-wise PL values along the algorithm, and $\tilde \mu_0$ and $\tilde L_0$ are estimation of the optimal constant involved in $\qginf$ and $\qg$, and $\gamma_{nm}\ge L$ is the used stepsize for \eqref{eq:nm} that arises from the tuning of parameters. The parameters used in \eqref{eq:nm} result from grid-search rather than the theoretical recommendation of Theorem~\ref{thm:pl_accel}, and we observe that \eqref{eq:nm} has a slightly better decreasing rate compared to the bound \eqref{cvmodel_bound_nm}.
Up to associating the empirical decreasing rate of \eqref{eq:gd} and \eqref{eq:nm} with the bounds \eqref{cvmodel_bound_gd} and \eqref{cvmodel_bound_nm}, we thus test the acceleration condition of Section~\ref{sec:mom_cv} by comparing $\tilde a $ with

\[\frac{ \tilde \mu_{gd}\gamma_{gd} }{\sqrt{\gamma_{nm}\tilde \mu_{nm}}}  \bpar{\frac{\tilde L_0}{\tilde \mu_0}}^{\frac{1}{4}}.\]
We find numerically that $\tilde a \approx 0.16$, while $ \frac{ \tilde \mu_{gd}\gamma_{gd} }{\sqrt{\gamma_{nm}\tilde \mu_{nm}}}  \bpar{\frac{\tilde L_0}{\tilde \mu_0}}^{\frac{1}{4}} \approx 0.19$ such that they are roughly the same. So, in this setting where \eqref{eq:nm} does not accelerate over \eqref{eq:gd}, we observe that the acceleration condition of Section~\ref{sec:mom_cv} is indeed not verified. This gives heuristic weights to the relevance of this bound. Finally, this suggests that the function that serves as lower bound for PL functions, which actually belongs to $\pl \cap \ac \cap \ls$, is such that the aiming condition values do not scale properly with other geometrical constants to permit acceleration using momentum.
\begin{figure}
    \centering
    \includegraphics[width=0.8\linewidth]{figures/pl_worst_instance.pdf}
    \caption{Function values decrease when running \eqref{eq:gd} and \eqref{eq:nm} on the hard instance of \cite{PLlowerbound}. We also plot the log of the theoretical convergence upper bound $\log((1-\mu \gamma)^k(f(x_0)-f^\ast)$ with $\mu$ the global PL constant, and the log of the theoretical convergence upper bound using parameter values across the path, see \eqref{cvmodel_bound_gd} and \eqref{cvmodel_bound_nm} in Section~\ref{app:heuristic}. It can be observed that \eqref{eq:gd} and \eqref{eq:nm} decrease at same speed, and that the bounds proposed at \eqref{cvmodel_bound_gd} and \eqref{cvmodel_bound_nm} are significantly more accurate compared with the global worst case bound.}
    \label{fig:hard_instance}
\end{figure}

\begin{figure}[t!]
        \centering
  \includegraphics[width=0.45\linewidth]{figures/mu_values.pdf}
    \includegraphics[width=0.45\linewidth]{figures/corr_values.pdf}

    \caption{Left: Values of $\mu(\tilde x_i)$ along the iterates of \eqref{eq:gd} and \eqref{eq:nm}. Values of $a(\tilde x_i)$ along the iterates of \eqref{eq:nm}. The underlying function is the hard instance of \cite{PLlowerbound}, see details in Appendix~\ref{app:hard_fonc_sqc}.} %Note that we can ensure to grasp more information if we also consider the constant $\mu_0$ such that $f \in \qginf$, which necessarily characterizes the red zone.} 
    \label{fig:mu_corr_values}
\end{figure}
\section*{Reproducibility Statement}
Our figures and experiments are reproducible, see our code \url{https://github.com/J-Hermant/Nesterov_Acceleration_Aiming_Condition/tree/main}.
\section*{Acknowledgment}
This work is supported by PEPR PDE-AI and the ANR SOS2ID (grant ANR-24-CE40-3786). We thank Jean-François Aujol, Erell Gachon and Marien Renaud for their feedback, which helped improve the presentation and clarity of this work. 
% \section*{Software and Data}

% If a paper is accepted, we strongly encourage the publication of software and data with the
% camera-ready version of the paper whenever appropriate. This can be
% done by including a URL in the camera-ready copy. However, \textbf{do not}
% include URLs that reveal your institution or identity in your
% submission for review. Instead, provide an anonymous URL or upload
% the material as ``Supplementary Material'' into the OpenReview reviewing
% system. Note that reviewers are not required to look at this material
% when writing their review.

% Acknowledgements should only appear in the accepted version.
% \section*{Acknowledgements}
%remercier erell, JF, Marien
% \textbf{Do not} include acknowledgements in the initial version of
% the paper submitted for blind review.

% If a paper is accepted, the final camera-ready version can (and
% usually should) include acknowledgements.  Such acknowledgements
% should be placed at the end of the section, in an unnumbered section
% that does not count towards the paper page limit. Typically, this will 
% include thanks to reviewers who gave useful comments, to colleagues 
% who contributed to the ideas, and to funding agencies and corporate 
% sponsors that provided financial support.
% \section{Concluding Remarks and Discussion}
% After highlighting the flaws of the class of strongly quasar convex functions to understand acceleration using momentum, we considered an aiming condition as a different viewpoint to analyze gradient descent with and without momentum. 
% \begin{itemize}
%     \item As an alternative viewpoint to address the issues presented in Section~\ref{sec:sqc_pitfalls}, one could consider $\sqcf{\tau}{0}\cap \qginf$, namely the class of quasar-convex functions with global quadratic lower bound. In particular, it solves the problem of parameter interpretation and optimality. However, $\sqcf\tau{0}$ carries more information than the alignment of the descent direction with the minimizers.
%     \item Our characterization with aiming condition is suboptimal for very simple functions, such as the strongly convex quadratic case. For those functions,  This is hinted by Proposition~\ref{lem:reparam_sqc}. \textit{e.g.} Considering the characterization of 
% \end{itemize}
% \section*{Impact Statement}
% This paper presents work whose goal is to advance the field of 
% Machine Learning. There are many potential societal consequences 
% of our work, none which we feel must be specifically highlighted here.




% In the unusual situation where you want a paper to appear in the
% references without citing it in the main text, use \nocite
% \nocite{langley00}
\bibliographystyle{plain} 
\bibliography{bib}
% \bibliographystyle{icml2025}


%%%%%%%%%%%%%%%%%%%%%%%%%%%%%%%%%%%%%%%%%%%%%%%%%%%%%%%%%%%%%%%%%%%%%%%%%%%%%%%
%%%%%%%%%%%%%%%%%%%%%%%%%%%%%%%%%%%%%%%%%%%%%%%%%%%%%%%%%%%%%%%%%%%%%%%%%%%%%%%
% APPENDIX
%%%%%%%%%%%%%%%%%%%%%%%%%%%%%%%%%%%%%%%%%%%%%%%%%%%%%%%%%%%%%%%%%%%%%%%%%%%%%%%
%%%%%%%%%%%%%%%%%%%%%%%%%%%%%%%%%%%%%%%%%%%%%%%%%%%%%%%%%%%%%%%%%%%%%%%%%%%%%%%
\newpage
\appendix
\onecolumn
\section{Related Works and Background Details}
In this section, we first provide bibliographical details about Table~\ref{table:sc_sqc_rates} (Appendix~\ref{app:related_work_tab}), then give more details about a specific related work in Appendix~\ref{app:related:uaac}.
% \jh{parler des systems dynamiques pré SUBOYDCANDES, type Bolte, jendoubi}
% \subsection{Other non-convex accelerations works}
% Hess lip, Third derivative Lip, detail SQC, quasi strong convexity (a priori, pas de preuve theoriques)

\subsection{Details for Table~\ref{table:sc_sqc_rates}}\label{app:related_work_tab}
In this section, we detail Table~\ref{table:sc_sqc_rates}, giving in particular references for the results. We note that some of the results do not exist in the literature. Before, we state the following result that we will use in this section.
\begin{proposition}\label{prop:equiv_hb_nmo}
    Let $\xz$ verify \eqref{eq:nmo}. Then, it verifies the Heavy Ball equation with Hessian damping
    \begin{equation}
    \ddot{x}_t
    + \bpar{\eta_t - \frac{\dot{\eta}_t}{\eta_t} + \eta'_t }\dot{x}_t
    + \gamma \nabla^2 f(x_t)\dot{x}_t
    + \bpar{ \eta'_t \gamma + \eta_t \gamma' -\frac{\dot{\eta}_t}{\eta_t}\gamma }\nabla f(x_t)
    = 0.
\end{equation}
\end{proposition}
Proposition~\ref{prop:equiv_hb_nmo} states that \eqref{eq:nmo} can be written as the Heavy Ball equation with Hessian damping \cite{alvarez2002second}. We will use results established for this equation in the case of \eqref{eq:nmo}.


\textbf{About $\strongconv$ \:} The rate for \eqref{eq:gf} can be deduced from the one in the $\pl$ case. The rate stated for \eqref{eq:nmo} can be found in the case of heavy ball with Hessian damping in \cite{shi2018understanding}, from which we deduce the same rate for \eqref{eq:nmo}.  The rates for \eqref{eq:gd} and \eqref{eq:nm} can be found in \cite{nesterovbook}.


\textbf{About $\sqc$ \:} See Proposition~\ref{prop:gf_sqc} for the rate of \eqref{eq:gf}. The rate of \eqref{eq:gd} can be found in \cite[Proposition 3]{hermant2024study}. The rate stated for \eqref{eq:nmo} can be found in the case of heavy ball with Hessian damping in \cite[Theorem 5]{hermant2024study}, from which we deduce the same rate for \eqref{eq:nmo}. The rate for \eqref{eq:nm} requires some care. It can be found in \cite[Theorem 1]{hinder2020near}, but involves a line-search procedure that requires computing evaluations of $f$ and $\nabla f$ at each iteration to compute one of the parameters. It can also be found in \cite[Theorem 2]{hermant2024study} up to a supplementary assumption of weak-convexity, namely $f + \rho/2\norm{\cdot}^2$ is convex for some $\rho \in \R$, with $\rho \le L$. Finally, it can be found in \cite[Theorem 2]{wang2023continuized} up to non-deterministic results, because of the use of the continuized parameterization.


\textbf{About $\pl$ \:} The rate $\mu$ for \eqref{eq:gf} can be found in \cite[Theorem 3]{polyak2017lyapunov}. The rate $\mu/L$ for \eqref{eq:gd} has been known for a long time \cite{Polyak1963,bolte2017error}, see \cite[Theorem 1]{karimi2016linear} for a direct proof. The rate for \eqref{eq:nm} can be found in \cite[Proposition 7]{hermant2025continuized}, but holds with high probability because of the use of the continuized framework. We note that using the Heavy Ball algorithm \eqref{eq:hb}, this rate can be improved to a rate close to $\sqrt{\mu/L}$ if assuming that $f$ is four times differentiable, and if the algorithm is initialized close enough to the minimizer \cite{kassing2024polyak}. The case of \eqref{eq:nmo} is more complicated. For Polyak's momentum (Heavy Ball equation), a result exists assuming $f$ is also in $\ls$, yielding a convergence rate $\approx \mu/\sqrt{L}$ \cite{apidopoulos2022convergence}. For the same equation, the rate $\sqrt{\mu}$ holds asymptotically, provided that $f$ is four times differentiable, providing strong local regularity around minimizers \cite{kassing2024polyak}. In the case of Heavy Ball with Hessian damping, a rate $\sqrt{\mu}$ is deduced \cite{apidopoulos2025heavy}, but the gradient term in the equation is rescaled by $\bigO(1/\sqrt{\mu})$, which means we are minimizing a rescaled version of the function. In Proposition~\ref{prop:nmo_pl}, we actually show without assuming $f \in \ls$ or smoother than $C^1$--and with a considerably simpler proof than the previously mentioned works--that \eqref{eq:nmo} achieves a rate $\gamma \mu$, where $\gamma$ is a parameter of \eqref{eq:nmo}. It highlights that we can achieve whatever rate we want, as long as we take $\gamma$ big enough. However, in the discrete case, $\gamma$ corresponding intuitively to the stepsize, it appears clearly that the descent lemma induced by $f \in \ls$ will impose a bound on $\gamma$, namely $\gamma \le 1/L$, yielding a rate $\mu/L$.

\subsubsection{Proof of not existing rates of Table~\ref{table:sc_sqc_rates}}\label{app:table_proof}
% \begin{proposition}\label{prop:gf_sqc}
%     Let $f\in\sqc$, $\xcont \sim~ $\eqref{eq:gf}. Then,
%     \[f(x_t)-f^\ast = \bigO(e^{-\tau \mu t}).\]
% \end{proposition}
% \begin{proof}
%       Let be
% \begin{equation}
%     E_t = f(x_t)-f^\ast + \frac{\mu}{2}\norm{x_t-x^\ast}^2.
% \end{equation}
% We have
% \begin{align}
%     \dot E_t = \dotprod{\nabla f(x_t),\dot x_t} + \mu \dotprod{x_t-x^\ast,\dot x_t} &\overset{(i)}=  -\norm{\nabla f(x_t)}^2 - \mu\dotprod{\nabla f(x_t),x_t-x^\ast} \\
%     &\le - \mu\dotprod{\nabla f(x_t),x_t-x^\ast}\\
%    &\overset{(ii)}\le -\mu \tau\bpar{f(x_t)-f^\ast + \frac{\mu}{2}\norm{x_t-x^\ast}^2}\\
%     &= -\mu \tau E_t,
% \end{align} 
% where (i) is by definition of \eqref{eq:gd} and (ii) is because $f \in \sqc$.
% By integration, we deduce
% $$E_t \le e^{-\tau \mu t}E_0,$$
% thus the result.
% \end{proof}
\begin{proposition}\label{prop:gf_sqc}
Let $f\in\sqc$, $\xcont\sim~\eqref{eq:gf}$. Then,
\[
f(x_t)-f^\ast = \bigO(e^{-\tau \mu t}).
\]
\end{proposition}
\begin{proof}
Let
\[
E_t =f(x_t)-f^\ast + \frac{\mu}{2}\norm{x_t-x^\ast}^2.
\]
Then
\begin{align*}
\dot E_t
&= \dotprod{\grad{x_t},\dot x_t} + \mu \dotprod{x_t-x^\ast,\dot x_t}\\
&\overset{(i)}= -\norm{\grad{x_t}}^2 - \mu\dotprod{\grad{x_t},x_t-x^\ast}\\
&\le - \mu\dotprod{\grad{x_t},x_t-x^\ast}\\
&\overset{(ii)}\le -\mu \tau\Bigl(f(x_t)-f^\ast + \frac{\mu}{2}\norm{x_t-x^\ast}^2\Bigr)
= -\mu\tau E_t,
\end{align*}
where (i) follows from \eqref{eq:gf} and (ii) from $f\in\sqc$.
Integrating yields $E_t\le e^{-\mu\tau t}E_0$, hence $f(x_t)-f^\ast\le E_t=\bigO(e^{-\mu\tau t})$.
\end{proof}

\begin{proposition}\label{prop:nmo_pl}
Let $f \in \pl$. Let $\xcont \sim$ \eqref{eq:nmo} with $\eta_t\equiv \eta$ and $\eta'_t\equiv \eta'$ constants and under the constraint $\gamma = (\eta+\eta')/\mu$. Then, we have
\[
f(x_t)-f^\ast = \bigO(e^{-2\mu \gamma t}).
\]
\end{proposition}
\begin{proof}
Let
\[
E_t = f(x_t)-f^\ast + \frac{\delta}{2}\norm{x_t - z_t}^2.
\]
Then
\begin{align*}
\dot{E}_t
&= \dotprod{\grad{x_t}, \dot{x}_t} + \delta \dotprod{x_t - z_t,\dot{x}_t - \dot{z}_t}\\
&\overset{(i)}=
\dotprod{\grad{x_t},  \eta(z_t-x_t) - \gamma \grad{x_t)}}
+ \delta \dotprod{x_t - z_t, (\eta+\eta')(z_t-x_t) + (\gamma'-\gamma)\grad{x_t}}\\
&=
\bigl(\eta+\delta(\gamma-\gamma')\bigr)\dotprod{\grad{x_t},  z_t-x_t}
- \gamma \norm{\grad{x_t}}^2
- \delta(\eta + \eta')\norm{z_t-x_t}^2,
\end{align*}
where (i) follows from \eqref{eq:nmo}.
Using $f\in\pl$, we have $-\gamma \norm{\grad{x_t}}^2 \le -2\mu\gamma (f(x_t)-f^\ast)$, hence
\[
\dot E_t \le
\bigl(\eta+\delta(\gamma-\gamma')\bigr)\dotprod{\grad{x_t},  z_t-x_t}
-2\mu\gamma (f(x_t)-f^\ast)
- \delta(\eta + \eta')\norm{z_t-x_t}^2.
\]
Impose $\eta+\delta(\gamma-\gamma')=0$ (\textit{i.e.} $\delta=\eta/(\gamma'-\gamma)$) and $\eta+\eta'=\mu\gamma$.
Then $\dot E_t \le -2\mu\gamma E_t$, so $E_t \le e^{-2\mu\gamma t}E_0$.
Since $f(x_t)-f^\ast \le E_t$, the claim follows.
\end{proof}
\subsubsection{Proof of Proposition~\ref{prop:equiv_hb_nmo}}
The following proof is borrowed to \cite[Proposition 27]{hermant2025continuized}.
 From the first line of
\eqref{eq:nmo}, we have
\begin{equation}\label{eq:z_as_x_constgam}
    z_t = x_t + \frac{1}{\eta_t}\bpar{\dot{x}_t +  \gamma \nabla f(x_t)}.
\end{equation}
Differentiating \eqref{eq:z_as_x_constgam} yields
\begin{align}
    \dot{z}_t
    &= \dot{x}_t - \frac{\dot{\eta}_t}{\eta_t^2}\bpar{\dot{x}_t +  \gamma \nabla f(x_t)}
    + \frac{1}{\eta_t}\bpar{\ddot{x}_t +  \gamma \nabla^2 f(x_t)\dot{x}_t}.
\end{align}
Using this expression and the second line of \eqref{eq:nmo}, we obtain
\begin{align}
     \dot{x}_t - \frac{\dot{\eta}_t}{\eta_t^2}\bpar{\dot{x}_t +  \gamma \nabla f(x_t)}
     + \frac{1}{\eta_t}\bpar{\ddot{x}_t +  \gamma \nabla^2 f(x_t)\dot{x}_t}
     &=  \eta'_t(x_t-z_t) - \gamma' \nabla f(x_t) \\
     &= -\frac{\eta'_t}{\eta_t}\dot{x}_t -  \frac{\eta'_t \gamma}{\eta_t}\nabla f(x_t) - \gamma' \nabla f(x_t),
\end{align}
where we used $x_t-z_t = -\frac{1}{\eta_t}\bpar{\dot{x}_t +  \gamma \nabla f(x_t)}$.
Multiplying by $\eta_t$ and rearranging, we get
\begin{equation}\label{eq:second_order_x_constgam}
    \ddot{x}_t
    + \bpar{\eta_t - \frac{\dot{\eta}_t}{\eta_t} + \eta'_t }\dot{x}_t
    + \gamma \nabla^2 f(x_t)\dot{x}_t
    + \bpar{ \eta'_t \gamma + \eta_t \gamma' -\frac{\dot{\eta}_t}{\eta_t}\gamma }\nabla f(x_t)
    = 0.
\end{equation}

\subsection{Comparison with Theorem 7 of \cite{hermant2024study}}\label{app:related:uaac}
\begin{theorem}[\cite{hermant2024study}]\label{thm:uaac}
    Let $f \in \pl$, with a unique minimizer $x^\ast$. Then, there exists $(\tau,\mu') \in (0,1] \times \mathbb{R}^\ast_+$ such that $f\in \sqcf{\tau}{\mu'}$ if and only if
          \begin{equation}\label{eq:uaac}
                          \exists a > 0,  \forall x \in \mathbb{R}^d, \quad 1 \geqslant \frac{\langle \nabla f(x),x - x^\ast \rangle}{\lVert  \nabla f(x) \rVert \lVert x - x^\ast \rVert} \geqslant a> 0,
          \end{equation}
  In particular if (\ref{eq:uaac}) holds, then as long as $\tau < \frac{2 \mu a}{L}$, $F$ is $(\gamma, \frac{\mu a}{\tau} - \frac{L}{2})$-strongly quasar-convex.
\end{theorem}
 The parameters derived in Theorem~\ref{thm:uaac} are significantly different compared to those we deduce from our analysis, see Appendix~\ref{app:link}. Theorem~\ref{thm:uaac} ensures that, at best, we can take $\tau = \bigO(a\frac{\mu}{L})$, inducing $\mu' = \bigO(L)$. We are in the setting where the condition number interpretation fails for strongly quasar-convex functions, see Section~\ref{sec:sqc_pitfalls}. Also, plugging these choices of parameters in Table~\ref{table:sc_sqc_rates} in the line of convergence rate for $\sqc$, if the convergence rates of \eqref{eq:gf} and \eqref{eq:gd} remain the same, namely $a \mu$ and $a \mu /L$ respectively, the convergence rates of \eqref{eq:nmo} and \eqref{eq:nm} however also reduce to $a \mu$ and $a \mu /L$ respectively, such that we do not obtain accelerated bound with Nesterov's momentum.
\subsection{Choice of Parameters for $\sqc$.}\label{app:choice_param_sqc}
In Section~\ref{sec:sqc_pitfalls}, we discuss that $ \sup_\tau  \sup_\mu \{(\tau,\mu) \in (0,1] \times \R_+: f \in \sqc \}$ and $\sup_\mu \sup_\tau \{(\tau,\mu) \in (0,1] \times \R_+: f \in \sqc \}$ leads to different choices of parameters. To justify this statement, we assume there exists $\tau, \mu \in (0,1]\times \R_+$ such that $f \in \sqc$. We first recall
$f \in \sqcf{\tau}{\mu} \Rightarrow f \in \sqcf{\theta \tau}{\mu/\theta}$, for all $\theta \in (0,1]$ \cite[Observation 5]{hinder2020near}. Letting $\theta \to 0$, it follows that
$$\arg \sup_\tau \arg \sup_\mu \{ (\tau,\mu) \in (0,1] \times \R_+:f \in \sqc \} = (0,+\infty).$$
On the other hand, it is clear that
$\arg \sup_\mu \arg \sup_\tau \{ (\tau,\mu)\in (0,1] \times \R_+ : f \in \sqc \} \neq (0,+\infty)$.
% \begin{align}\left\{
%     \begin{array}{ll}
%         \tilde{y}_k &= (\tilde{z}_{k} - \tilde{x}_k)\frac{\eta}{\eta+\eta'}\bpar{1 -e^{-(\eta + \eta')(T_{k+1}-T_k)} } + \tilde{x}_k \\
%         \tilde{x}_{k+1} &=  \tilde{y}_k  - \gamma \nabla f(\tilde{y}_k,\xi_{k+1})\\
%         \tilde{z}_{k+1} &=   \tilde{z}_{k} + \eta' \frac{\bpar{1- e^{-(\eta+\eta')(T_{k+1}-T_k)}}}{\eta' + \eta e^{-(\eta+\eta')(T_{k+1}-T_k)}}(\tilde{y}_k - \tilde{z}_{k}) - \gamma' \nabla f(\tilde{y}_k,\xi_{k+1})
%     \end{array}
% \right.
% \end{align}
\section{More Details on Figures and Numerical Experiments}
In this section, we provide details about our method to generate Figure~\ref{fig:pl_sqc_low_corr}, Figure~\ref{fig:sqc_flaws} (Appendix~\ref{app:param_details}) and Figure~\ref{fig:toy_prob} (Appendix~\ref{app:toy_prob}). We provide in Appendix~\ref{app:heuristic} a heuristic discussion of the bound over aiming condition to have acceleration derived in Section~\ref{sec:mom_cv}. We also provide in Appendix~\ref{app:conceptual_figures} additional figures, useful for the main text but that could not be included because of space limitations. Our figures and experiments are reproducible; see the code in the supplement.
\subsection{Details about Figure~\ref{fig:pl_sqc_low_corr} and Figure~\ref{fig:sqc_flaws}}\label{app:param_details}
In this section, we detail how we numerically compute the convergence rate of the functions displayed in Figure~\ref{fig:pl_sqc_low_corr} and Figure~\ref{fig:sqc_flaws}. We split our presentation between the two underlying functions.

\textbf{One-dimensional example \:}
The function displayed on the left in Figure~\ref{fig:pl_sqc_low_corr} is
\[f(t) = 5(t+0.19\sin(5t))^2, \]
defined on $[-2,2]$. It has $0$ as its unique critical point and minimizer, which satisfies $f(0) = 0$. We define a grid of $10^{5}$ points in $[-2,2]$, equally spaced. This grid has been refined until the subsequent numerical values are not impacted by further refinement. For the $i$-th point $x_i$ in this grid, we compute the associated point-wise PL value
\[ \mu_i = \frac{f'(x_i)^2}{2f(x_i)},  \]
and point-wise $L$-smooth value
\[ L_i = |f''(x_i)|.\]
We define $\mu := \min_{1\le i\le 10^5} \mu_i$ and $L := \max_{1\le i \le 10^5} L_i$ the constant such that numerically, $f \in \pl \cap \ls$ on $[-2,2]$. On this domain, we thus obtain a convergence rate $\mu/L$ for \eqref{eq:gd}.

Because of the double parameterization, the case of $\sqc$ is slightly more involved. We define a grid of $1000$ potential $\tau$ parameters, equally spaced in $[10^{-5},0.1]$. Note that we take $10^{-5}$ as the lower bound for this interval instead of zero, to avoid dividing by zero. For a given $\tau_j$ in the grid on $[10^{-5},0.1]$ , and each $x_i$ belonging to the grid on $[-2,2]$, we compute the point-wise strong quasar convexity constant $\mu'_{i,j}$
\[ \mu'_{i,j} = -2\frac{\Big( f(x_i) - f'(x_i)x_i \Big)}{x_i^2}.  \]
For this $\tau_j$, we define $\mu'_j = \min_{1\le i \le 10^5} \mu_{i,j}$ the constant such that numerically, $f \in \sqcf{\tau_j}{\mu_j}$. In the case where $\mu_j \le 0$, we consider $\tau_j$ is not admissible, which occurs for $\tau_j \ge 0.1$. The pairs $(\tau_j,\mu_j)_{1\le j\le 1000}$ thus generate Figure~\ref{fig:sqc_flaws}, by plotting $(\tau_j \mu_j/L)_{1\le j \le 1000 }$ for \eqref{eq:gd} and $(\tau_j \sqrt{\mu_j/L})_{1\le j \le 1000 }$ for \eqref{eq:nm}.  As a supplement, we plot in Figure~ $(\tau_j \mu_j)_{1\le j \le 1000 }$ and $(\tau_j \sqrt{\mu_j})_{1\le j \le 1000 }$ which are the convergence rates of \eqref{eq:gf} and \eqref{eq:nmo}. We observe a similar behavior compared with the discrete case. Finally, we obtain the convergence rate values of Figure~\ref{fig:pl_sqc_low_corr} by considering $\max_{1\le j \le 1000} \tau_j \mu_j/L$ for \eqref{eq:gd}, and $\max_{1\le j \le 1000} \tau_j \sqrt{\mu_j/L}$ for \eqref{eq:nm}. As illustrated in Figure~\ref{fig:sqc_flaws}, the pair $(\tau_{j_0},\mu_{j_0})$ that maximizes the rate of \eqref{eq:gd} differs from the pair $(\tau_{j_1},\mu_{j_1})$ that maximizes the rate of \eqref{eq:nm}.

\textbf{Two-dimensional example \:}
The function displayed on the left in Figure~\ref{fig:pl_sqc_low_corr} is
\[f(x,y) = 0.5(0.5x^2 -y)^2 + 0.05x^2,\]
defined on the square $[-1.2638,1.2638]^2$. It has a unique minimizer and critical point at $(0,0)$. To compute the convergence rate for this function, we proceed similarly to the one-dimensional case. The only difference is that we define a two-dimensional grid on $[-1.2638,1.2638]^2$ of $10^6$ points, equally spaced. The grid has been refined until the numerical values we compute are not impacted by further refinement.



\subsection{Details about Figure~\ref{fig:toy_prob}}\label{app:toy_prob}

In Figure~\ref{fig:toy_prob}, we optimize
\[F_{\varepsilon}(x,y) = 0.5(y-\sin(x))^2 + 0.5\varepsilon x^2.,\]
with $\varepsilon =10^{-3}$. It is clear that $F_{\varepsilon}(0,0) = 0$ is a global minimizer, and for all $(x,y) \neq (0,0)$, $F_{\varepsilon}(x,y) > 0$. This function has no other critical point, which is induced by the fact that it is PL.

\begin{lemma}\label{lem:toy_example_pl}
    For $x,y\in\R,\ \varepsilon>0$ the function $F_{\varepsilon}(x,y)=\tfrac12\,(y-\sin x)^2+\tfrac12\,\varepsilon x^2$ belongs to $\pl$ with
  $$\mu \le \frac{2+\varepsilon
-\sqrt{\varepsilon^2+4}}{2}.$$

\end{lemma}
\begin{proof}
    Consider the function
\[
F_{\varepsilon}(x,y)=\tfrac12\,(y-\sin x)^2+\tfrac12\,\varepsilon x^2,
\qquad x,y\in\R,\ \varepsilon>0.
\]

Let
\[
u(x,y) := y - \sin x.
\]
Then
\[
F_{\varepsilon}=\tfrac12\,(u^2+\varepsilon x^2).
\]

The gradient is
\[
\nabla F_{\varepsilon}=
\begin{pmatrix}
-u\cos x + \varepsilon x\\[3pt]
u
\end{pmatrix},
\qquad
\|\nabla F_{\varepsilon}\|^2 = u^2 + (-u\cos x + \varepsilon x)^2.
\]


Writing both sides as quadratic forms in \(z=(u(x,y),x)^\top\), we obtain
\[
2F_{\varepsilon}
= z^\top 
\begin{pmatrix}
1 & 0\\[3pt]
0 & \varepsilon
\end{pmatrix}z,
\qquad
\|\nabla F_{\varepsilon}\|^2
= z^\top
\begin{pmatrix}
1+\cos^2 x & -\varepsilon\cos x\\[3pt]
-\varepsilon\cos x & \varepsilon^2
\end{pmatrix}z.
\]

Define
\[
N=\begin{pmatrix}1&0\\0&\varepsilon\end{pmatrix},\qquad
M(x)=\begin{pmatrix}
1+\cos(x)^2 & -\varepsilon \cos(x)\\[3pt]
-\varepsilon \cos(x) & \varepsilon^2
\end{pmatrix},
\qquad c:=\cos x\in[-1,1].
\]
So, $f$ belong to $\pl$ with
\begin{equation}
    \begin{aligned}
        \mu = \min_{\substack{z=(u(x,y),x)\neq 0,\\(x,y)\in \R^2}}~ \frac{z^T M(x)z}{z^T N z} &\ge \min_{\substack{z=(z_1,z_2),\\(z_1,z_2)\in \R^2\backslash (0,0)}}~ \frac{z^T M(x)z}{z^T N z} \\&\overset{(i)}{=}  \min_{\substack{z=(z_1,z_2),\\(z_1,z_2)\in \R^2\backslash (0,0)}}~ \frac{z^T N^{-1/2}M(x)N^{-1/2}z}{z^Tz} \\
        &\overset{(ii)}{=} \lambda_{\min}(N^{-1/2}M(x)N^{-1/2}),
    \end{aligned}
\end{equation}
where (i) holds with the reparameterization $z' = N^{-1/2}z$, and (ii) is a property of the Rayleigh coefficient. We consider the matrix
\[
N^{-1/2}M(x)N^{-1/2}
=
\begin{pmatrix}
1+\cos(x)^2 & -\sqrt{\varepsilon}\,\cos(x)\\[3pt]
-\sqrt{\varepsilon}\,\cos(x) & \varepsilon
\end{pmatrix}
:= \begin{pmatrix}
1+c^2 & -\sqrt{\varepsilon}\,c\\[3pt]
-\sqrt{\varepsilon}\,c & \varepsilon
\end{pmatrix}=: B(c),\quad c:= \cos(x) \in [-1,1].
\]
So, we have
$$\lambda_{\min}(N^{-1/2}M(x)N^{-1/2}) \ge \inf_{c\in[-1,1]}\ 
\lambda_{\min}(B(c)).$$

The eigenvalues of \(B(c)\) are
\[
\lambda_\pm(c)
=
\frac{\operatorname{tr}(B(c))
\pm
\sqrt{\operatorname{tr}(B(c))^2 - 4\det(B(c))}}{2}.
\]

We have
\[
\det(B(c))=\varepsilon,
\qquad
\operatorname{tr}(B(c)) = 1 + c^2 + \varepsilon.
\]
Such that we solve
$$\min_{c \in [-1,1]} \frac{1+c^2 + \varepsilon - 
\sqrt{(1+c^2 + \varepsilon)^2 - 4\varepsilon}}{2} =\min_{t \in [0,1]} \frac{1+t + \varepsilon - 
\sqrt{(1+t + \varepsilon)^2 - 4\varepsilon}}{2} := \phi(t).$$
We have for $t \in [0,1]$,
$$2\phi'(t) = 1 - \frac{1+t+\varepsilon}{\sqrt{(1+t + \varepsilon)^2 - 4\varepsilon}} \le 0,$$
such that $\phi$ is decreasing on $[0,1]$. So, the minimum is reached for $t = 1 \Rightarrow c^2 = 1$, which leads to the eigenvalue
$$\inf_{c\in[-1,1]}\ 
\lambda_{\min}(B(c)) = \lambda_{\min}(B(1)) = \frac{2+\varepsilon
-\sqrt{\varepsilon^2+4}}{2}
>0.$$


Therefore, for all \((x,y)\in\R^2\),
\[
\|\nabla F_{\varepsilon}(x,y)\|^2 \;\ge\; 2\mu\, F_{\varepsilon}(x,y),
\]
with 
$$\mu \le \frac{2+\varepsilon
-\sqrt{\varepsilon^2+4}}{2}.$$
\end{proof}
The Hessian matrix of $F_{\varepsilon}$ at $(x,y)$ is
\[ \nabla^2  F_{\varepsilon}(x,y) = \begin{pmatrix}
   \cos(x)^2 + \sin(x)(y-\sin(x) +\varepsilon & -\cos(x) \\ -\cos(x) & 1
\end{pmatrix}.\]
The function is not globally $L$-smooth, because of the $y$ factor in $\nabla^2_{1,1} F_{\varepsilon}(x,y)$. However, for a fixed $y$, the eigenvalues are bounded for any $x\in \R$.

\textbf{Implementation Details \:} We fix $\varepsilon = 10^{-3}$, and chose $(0,3)$ as the initialization point. We compute the $L$-smooth value of $F_{0.001}$ on the domain $[-2\pi,2\pi] \times [-3,3]$, using a grid of $10^6$ equally spaced points. We then fix $\gamma = 1/L$ for both \eqref{eq:gd} and \eqref{eq:nm_prime}.

\subsection{Proof of Proposition~\ref{prop:hard_fonc_sqc}}\label{app:hard_fonc_sqc}
 The definition of the hard instance is rather involved. The core object of its definition is the following function
\[ g_{T,t}(x) = q_{T,t}(x) + \sum_{i=1}^{Tt}v_{y_i}(x_i), \]
where $x_i$ is the $i$-th coordinate of $x \in \R^{Tt} \to \R$, where $q_{T,t}(x)$ is defined as follows
\[q_{T,t}(x) = \frac{1}{2}\sum_{i=0}^{t-1}\left[ \Big(\frac{7}{8}x_{iT}-x_{iT+1}\Big)^2 + \sum_{j=1}^{T-1}(x_{iT+j+1}-x_{iT+j})^2 \right],\]
and for all $y\in \R$, $v_y(t)$ is the following one-dimensional function
\begin{equation*}
     v_y(t) = \left\{ \begin{array}{ll}
     \frac{1}{2}t,& t\le \frac{31}{32}y,  \\
     \frac{1}{2}t^2 - 16(t-\frac{31}{32}y)^2,& \frac{31}{32}y < t \le y,\\
     \frac{1}{2}t^2 - \frac{y^2}{32} + 16(t-\frac{33}{32}y)^2, &y<t\le \frac{33}{32}y,\\
     \frac{1}{2}t^2 - \frac{y^2}{32},&t>\frac{33}{32}y.\end{array} \right.  
\end{equation*}
The function $q_{T,t}$ belongs to the class of zero-chain functions \cite{nesterov2004introductory,lowerboundI,carmon2021lower}, which are hard to minimize with first-order algorithms. The key idea is that for these functions, if starting from $0$, first-order algorithms will be able to explore only one new coordinate at each iteration for a number of iterations equal to the dimension of the problem. Here, the $v_y(\cdot )$ components introduce non-convexity. Up to a rescaling, the function $g_{T,t}$ belongs to $\pl \cap \ls$ for desired constants $\mu$ and $L$, see all details in \cite{PLlowerbound}.


\textit{Proof of Proposition~\ref{prop:hard_fonc_sqc}.}
 The function $q_{T,t}$ is a strongly convex quadratic function with unique minimizer $0$, such that it belongs to $\sqc$ for some parameters. The $v_y(\cdot )$ components introduces non-convexity.  For any $y\in \R$, the function $v_y(\cdot)$ has zero as its unique minimizer, and has it belongs to $\pl \cap \ls$ for some parameter, which implies it belongs to $\sqc \cap \ls $ because the function is one-dimensional. It follows that as a sum of one-dimensional functions in $\sqc \cap \ls$, the function $ \sum_{i=1}^{Tt}v_{y_i}(x_i)$ belongs to $\sqc \cap \ls$ for some parameters, with unique minimizer $0 \in \R^d$. Then, by sum of strongly quasar-convex functions with same minimizer, $g_{t,T}$ is strongly quasar convex \cite[Appendix D.3]{hinder2020near}. Finally, to have the function $g_{T,t}$ belongs to $\pl \cap \ls$ for desired constants $\mu$ and $L$, it remains to rescale it as $\tilde g(x) = a_0g_{t,T}(y-a_1x)$ for some constants $a_0,a_1 \in \R$ and where $y = (y_1,\dots,y_{Tt})$, see all details in \cite{PLlowerbound}. Such rescaling preserve the strong quasar convexity property, see \cite[Appendix D.3]{hinder2020near}. 

\subsection{Supplementary Figures}
In this section, we provide some figures that are not in the main text, due to available space. 

\textbf{About Figure~\ref{fig:figure_sqc} \:}\label{app:conceptual_figures} We display in Figure~\ref{fig:figure_sqc} two synthetic examples of functions belonging to $\sqc$, in dimension one and two. The one-dimensional example is
\[f(t) =  0.5\cdot(t +0.15\sin(5t))^2.\]
The two dimensional example is $h: \mathbb{R}^2 \to \mathbb{R}$, built as follows:
  \begin{equation}
    h(x) = f(\lVert x \rVert)g\left(\frac{x}{\lVert x \rVert} \right)
\end{equation}
where $f(t) = t^2$ and 
\begin{equation}
    g(x_1,x_2) = \frac{1}{4N} \sum_{i=1}^N \left(a_i \sin(b_i x_1)^2 + c_i \cos(d_i x_2)^2 \right) + 1
\end{equation}
with $N = 10$ and the $\{ a_i \}_i$, $\{ c_i \}_i$ are independently and uniformly distributed on $[0,20]$, and the $\{ b_i \}_i$, $\{ d_i \}_i$ are independently and uniformly distributed on $[-25,25]$. It is borrowed from \cite{hermant2024study}.
\begin{figure}[ht]
    \centering
    % Colonne gauche (2 figures empilées)
    \begin{minipage}[b]{0.4\textwidth}
        \centering
        \includegraphics[width=\textwidth]{figures/f_sqc.pdf}\\[1ex]
        \includegraphics[width=\textwidth]{figures/fprime_sqc.pdf}
    \end{minipage}%
     % \hfill
    % Colonne droite (grande figure)
    \begin{minipage}[b]{0.38\textwidth}
        \centering
        \includegraphics[width=\textwidth]{figures/3d_sqc_cropped.pdf}
    \end{minipage}
    \caption{Visualization of functions belonging in $\sqc$, for some parameters. The figure is borrowed from \cite{hermant2025continuized}, see details in Appendix~\ref{app:conceptual_figures}.}
    \label{fig:figure_sqc}
\end{figure}


\textbf{About Figure~\ref{fig:quad_bounds} \:} Figure~\ref{fig:quad_bounds} displays a simple one-dimensional example where $f \in \ls \cap \qg$ with $L$ significantly larger than $L_0$, and similarly for $f \in \pl \cap \qginf$. Such differences are created by local variations outside of the vicinity of the minimizer. This is because the parameterizing $\qg$ and $\qginf$ is mainly affected by the geometry in the vicinity of the minimizer, while $\ls$ and $\pl$ are strongly sensitive to variations at any points.

\textbf{About Figure~\ref{fig:worst_case} \:} Figure~\ref{fig:worst_case} illustrates that, for a function belonging to $\pl$, the optimal $\mu$ may depend on the local smaller curvature at a minimizer $x^\ast$, or by plateaus. 
\begin{figure}
    \centering
    \includegraphics[width=0.5\linewidth]{figures/quad_bound_diff.pdf}
    \caption{ The grey curve is $f(t) = 0.5 \cdot  5(t+0.07\sin(13t))^2$. The blue (red) curve is a quadratic upper (lower) bound. We compute the largest $\mu$ and smallest $L$ such that  $f \in \pl \cap \ls$ on the displayed domain, numerically evaluated at $\mu \approx 4\cdot10^{-2}$, $L \approx 6\cdot 10^2$. Also, the parameters $L_0$ and $\mu_0$ that parameterize the quadratic bounds are $\mu_0 \approx 3$ and $L_0 \approx 18$. On this prototype 1-dimensional example, $\frac{\mu_0}{L_0} \approx 0.2$ while $\frac{\mu}{L} \approx 7\cdot 10^{-5}$. This highlight the possibility of a significant gap between these two ratios.}
    \label{fig:quad_bounds}
\end{figure}
\begin{figure}[t!]
        \centering
  \includegraphics[width=0.4\linewidth]{figures/worst_case_quad.pdf}
    \includegraphics[width=0.4\linewidth]{figures/worst_case_pl.pdf}

    \caption{Left: Simple $1$-dimensional functions in $\strongconv$, so also in $\pl$. The worst approaching direction toward the minimizer $t^\ast$ is the red zone, which determines the optimal $\mu$ constant. Right: Simple non-convex function that belongs to $\pl$ for some $\mu > 0$. The $\mu$ constant could be determined either by the local flatness around the minimizer $t^\ast$ (red zone), or by the flatness of the plateau (blue zone).} %Note that we can ensure to grasp more information if we also consider the constant $\mu_0$ such that $f \in \qginf$, which necessarily characterizes the red zone.} 
    \label{fig:worst_case}
\end{figure}
\begin{figure}
    \centering
    \includegraphics[width=0.4\linewidth]{figures/toy_example_fig_longer.pdf}
    \caption{ $3$-dimensional visualization of $F_{0.001}(x,y)= 0.5(y-\sin(x))^2 + 0.001\cdot0.5x^2$, considered in Section~\ref{sec:numerical}}
    \label{fig:3d_vizualization}
\end{figure}

\begin{figure}
    \centering

        \begin{minipage}{0.48\linewidth}
            \centering
            \includegraphics[width=\linewidth]{figures/toy_example_traj_long.pdf}
        \end{minipage}
    % Colonne droite (2 figures empilées)
    \hfill
    \begin{minipage}{0.5\linewidth}
        \centering
        \includegraphics[width=\linewidth]{figures/toy_example_fonc_long.pdf}

        \vspace{0.3cm}

        \includegraphics[width=\linewidth]{figures/toy_example_corr_long.pdf}
    \end{minipage}
    \caption{ Plot of the $3000$ first iterations of the trajectories of \eqref{eq:gd} and \eqref{eq:nm_prime} on $f$, when minimizing $F_{0.001}(x,y) = 0.5(y-\sin(x))^2 + 0.001\cdot0.5x^2$. This is a version of Figure~\ref{fig:toy_prob} with $3000$ iterations instead of $1000$. Observe that the decrease of the momentum method \eqref{eq:nm_prime} catches up with the one of \eqref{eq:gd}, which is coherent with the values of the aiming condition.
     }
    \label{fig:convergence_longer}
\end{figure}
\begin{figure}
    \centering
    \includegraphics[width=0.8\linewidth]{figures/sqc_bound_cont.pdf}
    \caption{For a range of parameters $\tau$, we compute the highest admissible $\mu$ ensuring $f \in \sqc $, and plot the associated theoretical convergence rates of \eqref{eq:gf} and \eqref{eq:nmo} under $\sqc$, namely $\tau \mu$ and $\tau \sqrt{\mu}$, with $f(t) = 5(t+0.19\sin(5t))^2$. See details in Appendix~\ref{app:param_details}. This is the continuous counterpart of Figure~\ref{fig:sqc_flaws}.
    %For each $\tau$, we empirically find the highest $\mu$ such that the function belongs to $\sqc$. We only consider value of $\tau$ such that the associated $\mu$ is positive. Note that the pair 
Similarly to the discrete setting, the pair $(\tau, \mu)$ that maximizes the convergence rate differ in both cases.}
    \label{fig:bound_cont}
\end{figure}

\section{Technical Lemmas}
The following lemmas serve as technical tools for our convergence proofs.
% \begin{lemma}\label{lem:scalar_prod_bound}
% Let $f \in \pl \cap \ac \cap \qg $. Let $
% \mu_0 := \sup\{\mu' \ge \mu : f \in \qginff{\mu'}\}$.
% Then, for all $x\in\R^d$,
% \[
% \dotprod{\nabla f(x), x^\ast - x}
% \le  -a \sqrt{\frac{\mu}{L_0}}\,(f(x)-f^\ast)
%       - \frac{a}{2}\sqrt{\mu \mu_0}\,\norm{x-x^\ast}^2.
% \]
% \end{lemma}




\begin{lemma}\label{lem:prod}
Let $f \in \pl \cap \qg$. Noting $
\mu_0 := \sup \{\mu'\ge \mu: f \in  \qginff{\mu'} \}$, 
we have for all $x\in\R^d$,
\[
\norm{x-x^\ast}\norm{\nabla f(x)}
\ge \sqrt{\frac{\mu}{L_0}}(f(x)-f^\ast) + \frac{\sqrt{\mu \mu_0}}{2}\norm{x-x^\ast}^2.
\]
\end{lemma}

\begin{proof}
If $x = x^\ast$, the statement is trivial. Assume $x \ne x^\ast$ and write
\[
\norm{x-x^\ast}\norm{\nabla f(x)}
= \norm{x-x^\ast}\sqrt{f(x)-f^\ast}\,\frac{\norm{\nabla f(x)}}{\sqrt{f(x)-f^\ast}}.
\]

\medskip
\noindent\textbf{(1) Step 1 ($\pl$ and $\qg$).}
Since $f\in\pl$, $\norm{\nabla f(x)}\ge \sqrt{2\mu}\sqrt{f(x)-f^\ast}$, and since $f\in\qg$,
\[
f(x)-f^\ast \le \frac{L_0}{2}\norm{x-x^\ast}^2
\quad\Rightarrow\quad
\norm{x-x^\ast}\ge \sqrt{\frac{2}{L_0}}\sqrt{f(x)-f^\ast}.
\]
Therefore,
\begin{equation}\label{eq:prod_0_corr}
\norm{x-x^\ast}\norm{\nabla f(x)}
\ge \sqrt{\frac{2}{L_0}}(f(x)-f^\ast)\sqrt{2\mu}
= 2\sqrt{\frac{\mu}{L_0}}(f(x)-f^\ast).
\end{equation}

\medskip
\noindent\textbf{Step 2 ($\pl$ and $\qginff{\mu_0}$).}
By definition of $\mu_0$
\[
f(x)-f^\ast \ge \frac{\mu_0}{2}\norm{x-x^\ast}^2
\quad\Rightarrow\quad
\sqrt{f(x)-f^\ast}\ge \sqrt{\frac{\mu_0}{2}}\norm{x-x^\ast}.
\]
Combining with $\norm{\nabla f(x)}\ge \sqrt{2\mu}\sqrt{f(x)-f^\ast}$ gives
\begin{equation}\label{eq:prod_1_corr}
\norm{x-x^\ast}\norm{\nabla f(x)}
\ge \sqrt{\mu\mu_0}\,\norm{x-x^\ast}^2.
\end{equation}
\medskip
\noindent\textbf{Combine
 the bounds}
From \eqref{eq:prod_0_corr} and the previous bound,
\[
\norm{x-x^\ast}\norm{\nabla f(x)}
\ge \max\Bigl\{2\sqrt{\frac{\mu}{L_0}}(f(x)-f^\ast),\ \sqrt{\mu\mu_0}\,\norm{x-x^\ast}^2\Bigr\}.
\]
Using $\max\{u,v\}\ge \frac{u+v}{2}$ for $u,v\ge 0$, we obtain
\[
\norm{x-x^\ast}\norm{\nabla f(x)}
\ge \sqrt{\frac{\mu}{L_0}}(f(x)-f^\ast) + \frac{\sqrt{\mu\mu_0}}{2}\norm{x-x^\ast}^2,
\]
which concludes the proof.
\end{proof}







% \begin{lemma}\label{lem:lyapunov_0_bound}
%     Let $f \in \qg$. Then.
%     $$f(x_0)-f^\ast + \frac{ \sqrt{\mu_0 L_0}}{2}\norm{x_0-x^\ast}^2 \le \frac{\sqrt{L_0}(\sqrt{L_0} + \sqrt{\mu_0})}{2}\norm{x-x^\ast}^2.$$
% \end{lemma}
% \begin{proof}
% By assumption, we have
%     $$f(x_0)-f^\ast \le \frac{L_0}{2}\norm{x_0-x^\ast}^2,$$ such that
%     $$f(x_0)-f^\ast + \frac{ \sqrt{\mu_0 L_0}}{2}\norm{x_0-x^\ast}^2 \le \bpar{\frac{L_0}{2} +  \frac{ \sqrt{\mu_0 L_0}}{2}}\norm{x-x^\ast}^2 = \frac{\sqrt{L_0}(\sqrt{L_0} + \sqrt{\mu_0})}{2}\norm{x-x^\ast}^2.$$
    
% \end{proof}

\begin{lemma}\label{lem:tec:exact_scalar_prod_bound}
Assume $x\neq x^\ast$ and $f(x)\neq f^\ast$. Define the point-wise quantities
\[
\mu(x) = \frac{1}{2}\frac{\norm{\nabla f(x)}^2}{f(x)-f^\ast},\qquad
a(x) =\frac{\dotprod{\nabla f(x),x-x^\ast}}{\norm{\nabla f(x)}\norm{x-x^\ast}},\qquad
\mu_0(x) = 2\frac{f(x)-f^\ast}{\norm{x-x^\ast}^2}.
\]
Then,
\[
\dotprod{\nabla f(x), x^\ast - x}
=  -a(x) \sqrt{\frac{\mu(x)}{\mu_0(x)}}(f(x)-f^\ast)
   - \frac{a(x)}{2}\sqrt{\mu(x)\mu_0(x)}\,\norm{x-x^\ast}^2.
\]

If $f \in \pl$, it follows
\[
\dotprod{\nabla f(x), x^\ast - x}
=  -a(x) \sqrt{\frac{\mu}{\mu_0}}(f(x)-f^\ast)
   - \frac{a(x)}{2}\sqrt{\mu\mu_0}\,\norm{x-x^\ast}^2,
\]
with $
\mu_0 := \sup \{\mu'\ge \mu: f \in  \qginff{\mu'} \}$.
\end{lemma}

\begin{proof}

By definition of $a(x)$,
\[
\dotprod{\nabla f(x), x^\ast - x} = -a(x)\norm{\nabla f(x)}\norm{x-x^\ast}.
\]
Moreover, by the definitions of $\mu(x)$ and $\mu_0(x)$,
\[
\norm{\nabla f(x)} = \sqrt{2\mu(x)}\sqrt{f(x)-f^\ast},
\qquad
\norm{x-x^\ast} = \sqrt{\frac{2}{\mu_0(x)}}\sqrt{f(x)-f^\ast}.
\]
Hence
\[
\norm{\nabla f(x)}\norm{x-x^\ast}
= 2\sqrt{\frac{\mu(x)}{\mu_0(x)}}\,(f(x)-f^\ast).
\]
Also, using $\norm{\nabla f(x)} = \sqrt{2\mu(x)}\sqrt{f(x)-f^\ast} =\sqrt{2\mu(x)} \sqrt{\frac{\mu_0(x)}{2}}\norm{x-x^\ast}$,
\[
\norm{\nabla f(x)}\norm{x-x^\ast}
= \sqrt{\mu(x)\mu_0(x)}\,\norm{x-x^\ast}^2.
\]
Combining the two equal expressions and using $\xi = \frac12\xi+\frac12\xi$ with
$\xi= \norm{\nabla f(x)}\norm{x-x^\ast}$ gives
\[
\norm{\nabla f(x)}\norm{x-x^\ast}
=
\sqrt{\frac{\mu(x)}{\mu_0(x)}}\,(f(x)-f^\ast)
+\sqrt{\mu(x)\mu_0(x)}\,\frac{1}{2}\norm{x-x^\ast}^2,
\]
and multiplying by $-a(x)$ yields the claim.
\end{proof}




% The following lemma is classical, see by instance \cite{nesterov2004introductory}.
% \begin{lemma}[Descent Lemma]\label{lem:descent_lemma}
%     Let $f\in\ls$, $x\in \R^d$. Then
%     \[ f\Big(x-\frac{1}{L}\nabla f(x)\Big) - f(x) \le \frac{1}{2L}\norm{\nabla f (x)}^2. \]
% \end{lemma}
\section{Proofs of Section~\ref{sec:cv_results}}\label{app:cv_results}
This section is dedicated to proving the results of Section~\ref{sec:cv_results}. In the discrete setting, we will actually prove a more general result, replacing the gradient with a stochastic estimator, under a strong growth condition, also sometimes named multiplicative noise assumption.\begin{assumption}[Strong growth condition]\label{ass:sgc}
    We assume we have access to stochastic approximations $\nabla f(x,\xi)$ of the real gradient $\nabla f(x)$, where $\xi \sim \mathcal{P}(\Xi)$, such that
    \begin{enumerate}[label=(\roman*)]
        \item (Unbiased estimator) $\forall x \in \R^d, ~ \E{\nabla f(x,\xi)} = \nabla f(x)$. 
        \item (Strong Growth Condition) $\exists \rho \ge 1, \forall x \in \R^d,~ \mathbb{E}_{\xi}\left[ \norm{\nabla f(x,\xi)}^2 \right] \le \rho \norm{\nabla f(x)}^2$. 
    \end{enumerate} 
\end{assumption}
To our knowledge, this assumption was introduced in \cite{polyak1987introduction}, and further used by \cite{cevher2019linear} and \cite{vaswani2019fast} as a relaxation of the maximal strong growth condition \cite{tseng1998incremental,solodov1998incremental}.
 We avoid this more general setting in the main text, to discard adding a layer of complexity. The deterministic case is a special instance of Assumption~\ref{ass:sgc}, corresponding to $\rho = 1$.
For functions in $\ls$ that verify Assumption~\ref{ass:sgc}, we have a stochastic descent lemma, see \cite[Lemma 16]{gupta2024nesterov}.
\begin{lemma}\label{lem:descent_sgc}
If $f \in \ls$ and Assumption~\ref{ass:sgc} holds, then we have

\begin{equation}
     \mathbb{E}_\xi\left[ f\Big(x-\frac{1}{\rho L}\nabla f(x,\xi)\Big) - f(x) \right] \leq -\frac{1}{2\rho L}\lVert  \nabla f(x) \rVert^2.
\end{equation}
\end{lemma} 
In the subsequent sections, for an algorithm $\xdisc$ we denote $\condE{\cdot} = \E{\cdot \: | \: \mathcal{F}_k }$. For a sequence $\{\xi_k \}_{k\in \N}$ of random variables, we note $\mathcal{F}_k = \sigma(\xi_0,\dots,\xi_k)$ the filtration generated by $\xi_0,\dots,\xi_k$.
\subsection{Proof of Theorem~\ref{thm:gf_pl_&_ac}}\label{app:gf_pl_&_ac}


\textbf{(i) $f \in \pl$ \:} The following proof is classical \cite{polyak2017lyapunov}. We provide a proof the sake of completeness.
Let
\[
E_t = f(x_t)-f^\ast.
\]
By definition of \eqref{eq:gf}, we have
\[
\dot E_t = \dotprod{\nabla f(x_t),\dot x_t}
= -\norm{\nabla f(x_t)}^2
\le -2\mu\,(f(x_t)-f^\ast)
= -2\mu\,E_t,
\]
where we used $f\in\pl$. By integration,
\[
f(x_t)-f^\ast \le e^{-2\mu t}(f(x_0)-f^\ast).
\]

\textbf{(ii) $f \in \pl \cap \ac \cap \qg$ \:}
Let
\[
E_t = f(x_t)-f^\ast + \frac{\beta}{2}\norm{x_t-x^\ast}^2.
\]
Then
\begin{align*}
\dot E_t
&= \dotprod{\nabla f(x_t),\dot x_t} + \beta \dotprod{x_t-x^\ast,\dot x_t}\\
&= -\norm{\nabla f(x_t)}^2 - \beta \dotprod{\nabla f(x_t),x_t-x^\ast}\\
&\le - \beta \dotprod{\nabla f(x_t),x_t-x^\ast}
= \beta \dotprod{\nabla f(x_t),x^\ast-x_t}.
\end{align*}
By Lemma~\ref{lem:scalar_prod_bound}, for all $t\ge 0$,
\[
\dotprod{\nabla f(x_t),x^\ast-x_t}
\le -a\sqrt{\frac{\mu}{L_0}}\,(f(x_t)-f^\ast)
     -\frac{a}{2}\sqrt{\mu\mu_0}\,\norm{x_t-x^\ast}^2,
\]
hence
\[
\dot E_t \le
-\beta a\sqrt{\frac{\mu}{L_0}}\,(f(x_t)-f^\ast)
-\frac{\beta a}{2}\sqrt{\mu\mu_0}\,\norm{x_t-x^\ast}^2.
\]
Choose $\beta$ so that the two coefficients match:
\[
\beta a\sqrt{\frac{\mu}{L_0}} = a\sqrt{\mu\mu_0}
\quad\Longrightarrow\quad
\beta = \sqrt{\mu_0 L_0}.
\]
Then
\[
\dot E_t \le -a\sqrt{\mu\mu_0}\Bigl(f(x_t)-f^\ast + \frac{\beta}{2}\norm{x_t-x^\ast}^2\Bigr)
= -a\sqrt{\mu\mu_0}\,E_t,
\]
and by integration,
\[
E_t \le e^{-a\sqrt{\mu\mu_0}\,t}E_0.
\]
Since $f\in\qginff{\mu_0}$, we have
\[
\frac{\beta}{2}\norm{x_0-x^\ast}^2
\le \beta\,\frac{1}{\mu_0}\,(f(x_0)-f^\ast)
= \sqrt{\frac{L_0}{\mu_0}}\,(f(x_0)-f^\ast),
\]
so $E_0 \le \bpar{1+\sqrt{\frac{L_0}{\mu_0}}}\,(f(x_0)-f^\ast)$.
Finally, since $f(x_t)-f^\ast\le E_t$, we deduce
\[
f(x_t)-f^\ast
\le \bpar{1+\sqrt{\frac{L_0}{\mu_0}}}\,e^{-a\sqrt{\mu\mu_0}\,t}\,(f(x_0)-f^\ast).
\]

\subsection{Proof of Theorem~\ref{thm:gd_pl_&_ac}}\label{app:gd_pl_ac}
We prove the following more general result.

\begin{theorem}\label{thm:sto_gd_pl_&_ac}
Let $f \in \ls$. Let $\xdisctilde \sim \eqref{eq:gd}$ with $\nabla f(\tilde x_k)$ replaced by
$\nabla f(\tilde x_k,\xi_k)$, where $\xi_k$ is independent of $\tilde x_k$, and assume Assumption~\ref{ass:sgc}.
Let $\gamma = \frac{1}{\rho L}$.

(i) If $f \in \pl$, then
\[
\E{ f(\tilde x_k)-f^\ast} \le \bpar{1-\frac{\mu}{\rho L}}^k\,(f(\tilde x_0)-f^\ast).
\]

(ii) If $f \in \pl \cap \ac \cap \qg$, then
\[
\E{f(\tilde x_k)-f^\ast}
\le
\bpar{1+\sqrt{\frac{L_0}{\mu_0}}}\,
\bpar{1-\frac{a\sqrt{\mu \mu_0}}{\rho L}}^k\,(f(\tilde x_0)-f^\ast),
\]
with $\mu_0 := \sup\{\mu' \ge \mu: f \in \qginff{\mu'}\}$ and
$L_0 :=\inf\{L' \le  L:f \in \qgf{L'}\}$.
\end{theorem}



The result for Theorem~\ref{thm:gd_pl_&_ac} is deduced in the case $\rho =1$, which removes the expectations. 

\begin{proof}[Proof of Theorem~\ref{thm:sto_gd_pl_&_ac}]
\textbf{(i) $f \in \pl$ \:} This result and proof is well known, see by instance \cite{karimi2016linear} in the deterministic setting, \cite{vaswani2019fast} for the stochastic case. We provide a proof for completeness.
Define \[E_k = f(\tilde x_k)-f^\ast.\] By Lemma~\ref{lem:descent_sgc},
\[
\condE{E_{k+1}-E_k}
=\condE{f(\tilde x_{k+1})-f(\tilde x_k)}
\le -\frac{1}{2\rho L}\norm{\nabla f(\tilde x_k)}^2.
\]
Since $f\in\pl$, $\norm{\nabla f(\tilde x_k)}^2 \ge 2\mu\,(f(\tilde x_k)-f^\ast)=2\mu E_k$, hence
\[
\condE{E_{k+1}} \le \bpar{1-\frac{\mu}{\rho L}}E_k.
\]
By induction and taking expectation,
\[
\E{E_k}\le \bpar{1-\frac{\mu}{\rho L}}^k E_0,
\]
which is the claim.

\medskip
\textbf{(ii) $f \in \pl \cap \ac \cap \qg$ \:}
Set $\beta := \sqrt{\mu_0 L_0}$ and define
\[
E_k =f(\tilde x_k)-f^\ast + \frac{\beta}{2}\norm{\tilde x_k-x^\ast}^2.
\]
Using the update
$\tilde x_{k+1}=\tilde x_k-\gamma\nabla f(\tilde x_k,\xi_k)$, we expand
\[
\norm{\tilde x_{k+1}-x^\ast}^2-\norm{\tilde x_k-x^\ast}^2
= -2\gamma \dotprod{\tilde x_k-x^\ast,\nabla f(\tilde x_k,\xi_k)}
+ \gamma^2\norm{\nabla f(\tilde x_k,\xi_k)}^2.
\]
Taking $\condE{\cdot}$, using Assumption~\ref{ass:sgc} (unbiasedness + SGC) and Lemma~\ref{lem:descent_sgc}, we obtain
\begin{align*}
\condE{E_{k+1}-E_k}
&\le -\frac{1}{2\rho L}\norm{\nabla f(\tilde x_k)}^2
+ \frac{\beta}{2}\Bigl(
-2\gamma \dotprod{\tilde x_k-x^\ast,\nabla f(\tilde x_k)}
+ \gamma^2\rho \norm{\nabla f(\tilde x_k)}^2
\Bigr)\\
&= \Bigl(-\frac{1}{2\rho L}+\frac{\beta}{2}\gamma^2\rho\Bigr)\norm{\nabla f(\tilde x_k)}^2
-\beta\gamma \dotprod{\tilde x_k-x^\ast,\nabla f(\tilde x_k)}.
\end{align*}
With $\gamma=\frac{1}{\rho L}$, we have $\gamma^2\rho=\frac{1}{\rho L^2}$, hence
\[
-\frac{1}{2\rho L}+\frac{\beta}{2}\gamma^2\rho
= -\frac{1}{2\rho L}+\frac{\beta}{2\rho L^2}
= -\frac{1}{2\rho L}\bpar{1-\frac{\beta}{L}} \le 0,
\]
because $\beta=\sqrt{\mu_0L_0}\le L_0\le L$. Therefore,
\[
\condE{E_{k+1}-E_k}
\le -\frac{\beta}{\rho L}\dotprod{\tilde x_k-x^\ast,\nabla f(\tilde x_k)}
= \frac{\beta}{\rho L}\dotprod{\nabla f(\tilde x_k),x^\ast-\tilde x_k}.
\]
By Lemma~\ref{lem:scalar_prod_bound},
\[
\dotprod{\nabla f(\tilde x_k),x^\ast-\tilde x_k}
\le -a\sqrt{\frac{\mu}{L_0}}\,(f(\tilde x_k)-f^\ast)
-\frac{a}{2}\sqrt{\mu\mu_0}\,\norm{\tilde x_k-x^\ast}^2.
\]
Multiplying by $\beta/(\rho L)$ and using $\beta\sqrt{\mu/L_0}=\sqrt{\mu\mu_0}$ (since $\beta=\sqrt{\mu_0L_0}$), we get
\[
\condE{E_{k+1}-E_k}
\le -\frac{a\sqrt{\mu\mu_0}}{\rho L}\,(f(\tilde x_k)-f^\ast)
-\frac{a\sqrt{\mu\mu_0}}{\rho L}\,\frac{\beta}{2}\norm{\tilde x_k-x^\ast}^2
= -\frac{a\sqrt{\mu\mu_0}}{\rho L}E_k.
\]
Hence
\[
\condE{E_{k+1}}\le \bpar{1-\frac{a\sqrt{\mu\mu_0}}{\rho L}}E_k,
\]
taking the expectation and by induction,
\[
\E{E_k}\le \bpar{1-\frac{a\sqrt{\mu\mu_0}}{\rho L}}^k E_0.
\]
Finally, since $f(\tilde x_k)-f^\ast\le E_k$, and using $f\in\qginff{\mu_0}$,
\[
\frac{\beta}{2}\norm{\tilde x_0-x^\ast}^2
\le \beta\,\frac{1}{\mu_0}(f(\tilde x_0)-f^\ast)
= \sqrt{\frac{L_0}{\mu_0}}(f(\tilde x_0)-f^\ast),
\]
so $E_0\le \bpar{1+\sqrt{\frac{L_0}{\mu_0}}}(f(\tilde x_0)-f^\ast)$. This yields the claim.
\end{proof}

\subsection{Proof of Theorem~\ref{thm:pl_accel}}\label{app:pl_accel}

\textbf{(i) \:}
Assume that $\eta_t\equiv\eta$, $\eta'_t\equiv\eta'$ are constants.
Let
\begin{equation}\label{eq:E_pl_accel}
    E_t = e^{\theta t}\Bigl(f(x_t)-f^\ast +\beta \frac{1}{2}\norm{z_t -  x^\ast}^2\Bigr).
\end{equation}
Then $E_t = E_0 + \int_0^t \dot{E}_s\,ds$, where
\begin{align}
    \dot{E}_t
    &= \theta  e^{\theta t}(f(x_t)-f^\ast)  +\theta \beta \frac{ e^{\theta t}}{2}\norm{z_t -  x^\ast}^2
    + e^{\theta t}\dotprod{\nabla f(x_t),\dot{x}_t} + \beta  e^{\theta t}\dotprod{z_t -  x^\ast,\dot{z}_t}
   \nonumber \\
    &= \theta  e^{\theta t}(f(x_t)-f^\ast)  +\theta \beta \frac{ e^{\theta t}}{2}\norm{z_t -  x^\ast}^2
    + e^{\theta t}\dotprod{\nabla f(x_t), \eta(z_t-x_t) - \gamma \nabla f(x_t)}\nonumber\\
    &\hspace{2cm}+ \beta   e^{\theta t}\dotprod{z_t - x^\ast,\eta'(x_t-z_t) - \gamma' \nabla f(x_t) } \nonumber\\
    &= \theta  e^{\theta t}(f(x_t)-f^\ast)  +\theta \beta \frac{ e^{\theta t}}{2}\norm{z_t -  x^\ast}^2
    - \gamma e^{\theta t}\norm{\nabla f(x_t)}^2 \nonumber\\
    &\hspace{2cm}+ e^{\theta t}\eta \dotprod{\nabla f(x_t), z_t-x_t}
    + \beta\eta' e^{\theta t}\dotprod{z_t - x^\ast,x_t - z_t}
    - \beta\gamma' e^{\theta t}\dotprod{\nabla f(x_t),z_t-x^\ast}.\label{eq:nest_ode_0_rew}
\end{align}
Using the decomposition
\[
\dotprod{\nabla f(x_t), z_t-x_t}
=
\dotprod{\nabla f(x_t), x^\ast-x_t}+\dotprod{\nabla f(x_t), z_t-x^\ast},
\]
and the identity
\[
\dotprod{z_t-x^\ast,x_t - z_t}
= \frac{1}{2}\norm{x_t-x^\ast}^2 -\frac{1}{2}\norm{z_t-x^\ast}^2-\frac{1}{2}\norm{x_t-z_t}^2
\le \frac{1}{2}\norm{x_t-x^\ast}^2 -\frac{1}{2}\norm{z_t-x^\ast}^2,
\]
we obtain from \eqref{eq:nest_ode_0_rew}
\begin{align}
    e^{-\theta t}\dot{E}_t
    &\le \theta(f(x_t)-f^\ast)
    + \frac{\beta}{2}\bigl(\theta-\eta'\bigr)\norm{z_t-x^\ast}^2
    + \frac{\beta\eta'}{2}\norm{x_t-x^\ast}^2 \nonumber\\
    &\hspace{1.5cm}+ \eta \dotprod{\nabla f(x_t), x^\ast-x_t}
    + (\eta-\beta\gamma')\dotprod{\nabla f(x_t), z_t-x^\ast}
    - \gamma \norm{\nabla f(x_t)}^2.\label{eq:pl_corr_nest_1_alt}
\end{align}

We choose $\eta'=\theta$ to cancel the $\norm{z_t-x^\ast}^2$ term, and we choose $\eta=\beta\gamma'$ to cancel the
$\dotprod{\nabla f(x_t),z_t-x^\ast}$ term. Dropping the negative term $-\gamma\norm{\nabla f(x_t)}^2$, we get
\begin{equation}\label{eq:pl_corr_nest_1_alt_simpl}
    e^{-\theta t}\dot{E}_t
    \le \theta(f(x_t)-f^\ast)
    + \eta \dotprod{\nabla f(x_t), x^\ast-x_t}
    + \frac{\beta\theta}{2}\norm{x_t-x^\ast}^2.
\end{equation}
By Lemma~\ref{lem:scalar_prod_bound},
\[
\dotprod{\nabla f(x_t), x^\ast-x_t}
\le  -a \sqrt{\frac{\mu}{L_0}}(f(x_t)-f^\ast) - \frac{a}{2}\sqrt{\mu \mu_0}\norm{x_t-x^\ast}^2.
\]
Plugging into \eqref{eq:pl_corr_nest_1_alt_simpl} yields
\begin{align}
    e^{-\theta t}\dot{E}_t
    &\le \Bigl(\theta-\eta a\sqrt{\frac{\mu}{L_0}}\Bigr)(f(x_t)-f^\ast)
    + \frac12\Bigl(\beta\theta-\eta a\sqrt{\mu\mu_0}\Bigr)\norm{x_t-x^\ast}^2.\label{eq:pl_corr_nest_4_alt}
\end{align}
We impose the cancellation of both coefficients in \eqref{eq:pl_corr_nest_4_alt}:
\[
\theta=\eta a\sqrt{\frac{\mu}{L_0}}
\qquad\text{and}\qquad
\beta\theta=\eta a\sqrt{\mu\mu_0}.
\]
These two equalities are equivalent to $\beta=\sqrt{\mu_0L_0}$ and $\theta=\eta a\sqrt{\mu/L_0}$.
With the additional choice $\eta=\beta\gamma'$, we may take
\[
\beta=\sqrt{\mu_0L_0},\qquad
\gamma'=\beta^{-1/2}=(\mu_0L_0)^{-1/4},\qquad
\eta=\beta\gamma'=(\mu_0L_0)^{1/4},
\qquad
\eta'=\theta=a\Bigl(\frac{\mu_0}{L_0}\Bigr)^{1/4}\sqrt{\mu}.
\]
With these parameters, \eqref{eq:pl_corr_nest_4_alt} gives $e^{-\theta t}\dot E_t\le 0$, hence $E_t\le E_0$.
Therefore, using \eqref{eq:E_pl_accel},
\[
f(x_t)-f^\ast \le e^{-\theta t}E_t \le e^{-\theta t}E_0
= \exp\!\Bigl(-a\Bigl(\frac{\mu_0}{L_0}\Bigr)^{1/4}\sqrt{\mu}\,t\Bigr)\,E_0.
\]
Finally, using $x_0 = z_0$ and $f\in\qginff{\mu_0}$,
\[
\frac{\beta}{2}\norm{x_0-x^\ast}^2
\le \beta\,\frac{1}{\mu_0}(f(x_0)-f^\ast)
= \sqrt{\frac{L_0}{\mu_0}}(f( x_0)-f^\ast),
\]
so $E_0\le \bpar{1+\sqrt{\frac{L_0}{\mu_0}}}(f( x_0)-f^\ast)$, yields the claim.


 
\textbf{(ii) \:} The proof in the discrete case exploits the continuized framework. This implies that the following analysis differs from the previous ones. We refer to Appendix~\ref{app:continuized_intro} to a brief introduction of the tools needed when carrying such analysis, or \cite{hermant2025continuized} for a more detailed introduction. Also, we will prove the following more general result.
\begin{theorem}\label{thm:sto_pl_accel}
Let $f\in \pl \cap \ac \cap \ls$.  Denote $L_0 := \inf \{L'\le L:  f \in \qgf{L'} \}$, $\mu_0 := \sup\{\mu' \ge \mu : f \in \qginff{\mu'}\}$.
 Let $\xdisctilde$ verify \eqref{eq:nm} with the \hyperlink{cont_param}{continuized parameterization} and with $\nabla f(\tilde x_k)$ replaced by $\nabla f(\tilde x_k,\xi_k)$ for random variables $\xi_k$ independent of $\tilde x_k$ and under Assumption~\ref{ass:sgc}. Let $\eta' = \frac{a}{\rho}\bpar{\frac{\mu_0}{L_0}}^{1/4}\sqrt{\frac{\mu}{L}},\quad \gamma = \frac{1}{\rho L},\quad \gamma' = \frac{1}{\rho} \frac{1}{( \mu_0 L_0)^{1/4}\sqrt{ L}}, $ and $\eta = \frac{1}{\rho}\frac{(\mu_0 L_0)^{1/4}}{\sqrt{L}}$. Then,  
    with probability  $1-\frac{1}{c_0} - \exp(-(c_1-1-\log(c_1))k)$, for some $c_0 > 1$ and $c_1 \in (0,1)$ we have
$$ f(\tilde x_k)-f^\ast \le K_1e^{-\frac{a}{\rho}\bpar{\frac{\mu_0}{L_0}}^{1/4}\sqrt{\frac{\mu}{L}}(1-c_1)k}, $$
where $K_1 := c_0\bpar{1+\sqrt\frac{L_0}{\mu_0}}(f(x_0)-f^\ast)$.
\end{theorem}
The proof of Theorem~\ref{thm:pl_accel} (ii) is deduced in the case $\rho = 1$.
% \begin{proof}[Proof of Theorem~\ref{thm:sto_pl_accel}]
%    Let $\overline{x}_t= (t,x_t,z_t)$. It satisfies $d\overline{x}_t = b(\overline{x}_t)dt + G(\overline{x}_t,\xi)dN(t,\xi)$, where
% \begin{equation*}
%     b(\overline{x}_t) = \begin{pmatrix}
%     1 \\ \eta_t(z_t-x_t) \\ \eta_t'(x_t-z_t)
%     \end{pmatrix},\quad G(\overline{x}_t,\xi) = \begin{pmatrix}
%     0 \\ -\gamma_t \nabla f(x_t,\xi) \\ -\gamma_t' \nabla f(x_t,\xi).
%     \end{pmatrix}
% \end{equation*}
% We apply Proposition~\ref{prop:sto_calc} to $\phi_t= \varphi(\overline{x}_t)$, where
% \begin{equation*}
%     \varphi(t,x,z) =  A_t(f(x)-f^\ast) + \frac{B_t}{2}\norm{z - x^\ast}^2.
% \end{equation*}
% so 
% \begin{equation}\label{eq:prop_calc_sto}
%     \varphi(\overline{x}_t) = \varphi(\overline{x}_0) + \int_{0}^t \dotprod{\nabla \varphi(\overline{x}_s),b(\overline{x}_s)}ds + \int_{[0,t]} \mathbb{E}_\xi( \varphi(\overline{x}_s + G(\overline{x}_s,\xi)) - \varphi(\overline{x}_s))ds + M_t,
% \end{equation}
% where $M_t$ is a martingale.
% We have 
% \begin{align*}
%     \frac{\partial \varphi}{\partial s} &= \frac{dA_s}{ds}(f(x)-f^\ast + \frac{1}{2}\frac{dB_s}{ds}\norm{z - x^\ast}^2, \\
%     \frac{\partial \varphi}{\partial x} &= A_s\nabla f(x),\\
%     \quad  \frac{\partial \varphi}{\partial z} &= B_s(z - x^\ast).
% \end{align*}
% So, 
% \begin{align}
%     \dotprod{\nabla \varphi(\overline{x}_s),b(\overline{x}_s)} &=\frac{dA_s}{ds}(f(x_s)-f^\ast) + \frac{1}{2}\frac{dB_s}{ds}\norm{z_s - x^\ast}^2 \nonumber \\
%     &+\dotprod{A_s \nabla f(x_s),\eta_s(z_s-x_s)} + \dotprod{B_s(z_s-x^\ast),\eta_s'(x_s - z_s)} \nonumber\\
%     &=\frac{dA_s}{ds}(f(x_s)-f^\ast) + \frac{1}{2}\frac{dB_s}{ds}\norm{z_s - x^\ast}^2\label{eq:sqc_eq_1}\\
%     &+A_s \eta_s \dotprod{ \nabla f(x_s) ,z_s-x_s}+ B_s \eta_s' \dotprod{z_s-x^\ast,x_s - z_s}\label{eq:sqc_eq_2}
% \end{align}
% Also, 
% \begin{align}
%    \mathbb{E}_{\xi} \varphi(\overline{x}_s+ G(\overline{x}_s)) - \varphi(\overline{x}_s) &= \mathbb{E}_{\xi}[A_s(f(x_s - \gamma_s \nabla f(x_s,\xi)) - f(x_s))] \nonumber\\
%    &+ \frac{B_s}{2}\mathbb{E}_{\xi}[\norm{ z_s - \gamma'_s\nabla f(x_s,\xi) - x^\ast}^2 - \norm{z_s -x^\ast}^2]\nonumber\\
%     &= A_s\mathbb{E}_{\xi}[f(x_s - \gamma_s \nabla f(x_s,\xi)) - f(x_s)]  + \frac{B_s \gamma_s^{'2}}{2}\mathbb{E}_{\xi}[\norm{\nabla f(x_s,\xi)}^2] \label{eq:sqc_eq_3}\\
%     &+ B_s\gamma'_s \dotprod{\nabla f(x_s),x^\ast - z_s} \label{eq:sqc_eq_4}
% \end{align}
% where we used $\mathbb{E}_{\xi}[\nabla f(x_s,\xi)] = \nabla f(x_s)$.
% As $f \in \ls$ and as Assumption~\ref{ass:sgc} holds, by Lemma~\ref{lem:descent_sgc}, we have
% \begin{align}\label{eq:pl_avg_pre_1}
%     \eqref{eq:sqc_eq_3} \le \bpar{\rho B_s \gamma_s^{'2} + A_s \gamma_s( \rho\gamma_s L-2)}\frac{1}{2}\norm{\nabla f(x_s)}^2.
% \end{align}
% We rewrite the right term in \eqref{eq:sqc_eq_2}
% \begin{equation}\label{eq:pl_avg_pre_2}
%     \dotprod{z_s-x^\ast,x_s - z_s} = \frac{1}{2}\norm{x_s-x^\ast}^2 -\frac{1}{2}\norm{z_s-x^\ast}^2-\frac{1}{2}\norm{x_s-z_s}^2 \le\frac{1}{2}\norm{x_s-x^\ast}^2 -\frac{1}{2}\norm{z_s-x^\ast}^2.  
% \end{equation}
% Then, focusing on the left term in \eqref{eq:sqc_eq_2} and \eqref{eq:sqc_eq_4}
% \begin{align}\label{eq:pl_avg_pre_3}
%     A_s \eta_s \dotprod{ \nabla f(x_s) ,z_s-x_s} +  B_s\gamma'_s \dotprod{\nabla f(x_s),x^\ast - z_s} &= A_s \eta_s \dotprod{ \nabla f(x_s) ,x^\ast-x_s} + \bpar{A_s\eta_s - B_s \gamma'_s} \dotprod{ \nabla f(x_s) ,z_s-x_s}.
% \end{align}
% Using Lemma~\ref{lem:scalar_prod_bound}, we have
% \begin{equation*}
%      \dotprod{ \nabla f(x_s) ,x^\ast-x_s} \le  -a \sqrt{\frac{\mu}{L_0}}(f(x_t)-a \sqrt{\mu\mu_0}\norm{x_t-x^\ast}^2f^\ast).
% \end{equation*}
% Gathering the previous computations with \eqref{eq:sqc_eq_1}-\eqref{eq:sqc_eq_4}, we obtain that
% \begin{align*}
%     &\dotprod{\nabla \varphi(x_s),b(x_s)} +  \mathbb{E}_{\xi} \varphi(\overline{x}_s+ G(\overline{x}_s)) - \varphi(\overline{x}_s) \\
%     &\le \frac{1}{2}\bpar{ \frac{dB_s}{ds}-B_s \eta_s'}\norm{z_s-x^\ast}^2 +  \frac{1}{2}\bpar{  B_s \eta'_s - A_s \eta_sa \sqrt{\mu\mu_0}}\norm{x_s-x^\ast}^2 + \bpar{\frac{dA_s}{ds}- A_s \eta_sa \sqrt{\frac{\mu}{L_0}}}(f(x_s)-f^\ast)\\
%     & + \bpar{A_s\eta_s - B_s \gamma'_s} \dotprod{ \nabla f(x_s) ,z_s-x_s}+\bpar{\rho B_s \gamma_s^{'2} + A_s \gamma_s( \rho\gamma_s L-2)}\frac{1}{2}\norm{\nabla f(x_s)}^2.
% \end{align*}

% \noindent
% \textbf{Parameter tuning \:}
% We set $A_s = e^{\theta s}$, $B_s = \sqrt{\mu_0 L_0}e^{\theta_s}$, with $ \theta = \eta' = \frac{a}{\rho}\bpar{\frac{\mu_0}{L_0}}^{1/4}\sqrt{\frac{\mu}{L}},\quad \gamma_s = \frac{1}{\rho L},\quad \gamma' = \frac{1}{\rho} \frac{1}{( \mu_0 L_0)^{1/4}\sqrt{ L}}, $ and $\eta = \frac{1}{\rho}\frac{(\mu_0 L_0)^{1/4}}{\sqrt{L}}$. Precisely, we chose the same set of algorithm's parameter, simply normalizing by $1/\sqrt{L}$, and choosing the stepsize $\gamma = \frac{1}{\rho L}$ to optimize the descent lemma. This choice implies
% \begin{align}\label{eq:cancel_param:0}
%     A_s\eta_s - B_s \gamma'_s = 0
% \end{align}
% \begin{align}\label{eq:cancel_param:1}
%     \rho B_s \gamma_s^{'2} + A_s \gamma_s( \rho\gamma_s L-2) = \frac{e^{\theta s}}{\rho L} - \frac{e^{\theta s}}{\rho L}=0, 
% \end{align}
% \begin{align}\label{eq:cancel_param:2}
%     \frac{dB_s}{ds}-B_s \eta_s' = \theta \sqrt{\mu_0 L_0} e^{\theta s} - \sqrt{\mu_0 L_0} \theta e^{\theta s} = 0,
% \end{align}
% \begin{align}\label{eq:cancel_param:3}
%     \frac{dA_s}{ds}- A_s \eta a \sqrt{\frac{\mu}{L_0}} = \theta e^{\theta s} -  \frac{a}{\rho} \bpar{\frac{\mu_0}{L_0}}^{1/4}\sqrt{\frac{\mu}{L}}e^{\theta s} = 0 
% \end{align}
% \begin{align}\label{eq:cancel_param:4}
%    B_s \eta'_s - A_s \eta_s a \sqrt{\mu\mu_0} =  \sqrt{\mu_0 L_0}\frac{a}{\rho}\bpar{\frac{\mu_0}{L_0}}^{1/4}\sqrt{\frac{\mu}{L}}e^{\theta_s} -  \frac{1}{\rho}\frac{(\mu_0 L_0)^{1/4}}{\sqrt{L}} a \sqrt{\mu\mu_0}e^{\theta s}=0.
% \end{align}
% This set of parameters implies $\dotprod{\nabla \varphi(x_s),b(x_s)} +  \mathbb{E}_{\xi} \varphi(\overline{x}_s+ G(\overline{x}_s)) - \varphi(\overline{x}_s) \le 0$, which in turn implies
% \begin{equation}\label{eq:pl_corr_cont_supermart}
%     \phi(t) \le \phi(0) + M_t.
% \end{equation}
% We evaluate \eqref{eq:pl_corr_cont_supermart} at $t = T_k$, take expectation and use Theorem~\ref{thm:martingal_stopping} and Proposition~\ref{prop:discretization_constant_param}, to get
% \begin{equation}
% \E{e^{ \frac{a}{\rho}\bpar{\frac{\mu_0}{L_0}}^{1/4}\sqrt{\frac{\mu}{L}}T_k}f(\tilde x_k)-f^\ast)} \le \phi(0).
% \end{equation}
% Using $f(x_0)-f^\ast \ge \frac{\mu_0}{2}\norm{x_0-x^\ast}^2$, we deduce $\phi(0) = f(\tilde x_0)-f^\ast + \frac{\sqrt{\mu_0 L_0}}{2}\norm{\tilde x_0-x^\ast}^2 \le \bpar{1+\sqrt\frac{L_0}{\mu_0}}(f(x_0)-f^\ast)$
% %by Lemma~\ref{lem:lyapunov_0_bound}. 
% So, we deduce
% $$\E{e^{ \frac{a}{\rho}\bpar{\frac{\mu_0}{L_0}}^{1/4}\sqrt{\frac{\mu}{L}}T_k}f(\tilde x_k)-f^\ast)} \le  \bpar{1+\sqrt\frac{L_0}{\mu_0}}(f(x_0)-f^\ast). $$

% Finally, by Lemma~\ref{lem:conv_high_proba}, with probability  $1-\frac{1}{c_0} - e^{-(c_1-1-\log(c_1))k}$, for some $c_0 > 1$ and $c_1 \in (0,1)$ we have
% $$ f(\tilde x_k)-f^\ast \le \bpar{1+\sqrt\frac{L_0}{\mu_0}}(f(x_0)-f^\ast)e^{-\frac{a}{\rho}\bpar{\frac{\mu_0}{L_0}}^{1/4}\sqrt{\frac{\mu}{L}}(1-c_1)k}. $$
% \end{proof}











% \begin{theorem}\label{thm:sto_pl_accel}
% Let $f\in \pl \cap \ac \cap \ls$. Denote $L_0 := \inf \{L'\le L:  f \in \qgf{L'} \}$,
% $\mu_0 := \sup\{\mu' \ge \mu : f \in \qginff{\mu'}\}$.
% Let $\xdisctilde$ verify \eqref{eq:nm} with the \hyperlink{cont_param}{continuized parameterization} and with
% $\nabla f(\tilde x_k)$ replaced by $\nabla f(\tilde x_k,\xi_k)$ for random variables $\xi_k$ independent of $\tilde x_k$
% and under Assumption~\ref{ass:sgc}. Let
% \[
% \eta' = \frac{a}{\rho}\Bigl(\frac{\mu_0}{L_0}\Bigr)^{1/4}\sqrt{\frac{\mu}{L}},
% \quad
% \gamma = \frac{1}{\rho L},
% \quad
% \gamma' = \frac{1}{\rho}\frac{1}{(\mu_0 L_0)^{1/4}\sqrt{L}},
% \quad
% \eta = \frac{1}{\rho}\frac{(\mu_0 L_0)^{1/4}}{\sqrt{L}}.
% \]
% Then, with probability $1-\frac{1}{c_0} - \exp(-(c_1-1-\log(c_1))k)$, for some $c_0 > 1$ and $c_1 \in (0,1)$,
% we have
% \[
% f(\tilde x_k)-f^\ast \le K_1\exp\!\Bigl(-a\Bigl(\frac{\mu_0}{L_0}\Bigr)^{1/4}\sqrt{\frac{\mu}{L}}(1-c_1)k\Bigr),
% \]
% where one may take
% \[
% K_1 := c_0\Bigl(1+\sqrt{\frac{L_0}{\mu_0}}\Bigr)\,(f(\tilde x_0)-f^\ast).
% \]
% \end{theorem}

\begin{proof}[Proof of Theorem~\ref{thm:sto_pl_accel}]
Let $\overline{x}_t= (t,x_t,z_t)$. It satisfies
$d\overline{x}_t = b(\overline{x}_t)dt + G(\overline{x}_t,\xi)dN(t,\xi)$, where
\[
b(\overline{x}_t) = \begin{pmatrix}
1 \\ \eta(z_t-x_t) \\ \eta'(x_t-z_t)
\end{pmatrix},
\qquad
G(\overline{x}_t,\xi) = \begin{pmatrix}
0 \\ -\gamma \nabla f(x_t,\xi) \\ -\gamma' \nabla f(x_t,\xi)
\end{pmatrix}.
\]
We apply Proposition~\ref{prop:sto_calc} to $\phi(t)= \varphi(\overline{x}_t)$, where
\[
\varphi(t,x,z) =  A_t(f(x)-f^\ast) + \frac{B_t}{2}\norm{z - x^\ast}^2.
\]
Hence
\begin{equation}\label{eq:prop_calc_sto}
\varphi(\overline{x}_t) = \varphi(\overline{x}_0)
+ \int_{0}^t \dotprod{\nabla \varphi(\overline{x}_s),b(\overline{x}_s)}ds
+ \int_{0}^t \mathbb{E}_{\xi}\!\left[ \varphi(\overline{x}_s + G(\overline{x}_s,\xi)) - \varphi(\overline{x}_s)\right]ds
+ M_t,
\end{equation}
where $M_t$ is a martingale.

We have
\[
\frac{\partial \varphi}{\partial s}(\overline{x}_s)
= \frac{dA_s}{ds}\,(f(x_s)-f^\ast) + \frac{1}{2}\frac{dB_s}{ds}\norm{z_s - x^\ast}^2,
\qquad
\frac{\partial \varphi}{\partial x}(\overline{x}_s) = A_s\nabla f(x_s),
\qquad
\frac{\partial \varphi}{\partial z}(\overline{x}_s) = B_s(z_s - x^\ast).
\]
So,
\begin{align}
\dotprod{\nabla \varphi(\overline{x}_s),b(\overline{x}_s)}
&=\frac{dA_s}{ds}(f(x_s)-f^\ast) + \frac{1}{2}\frac{dB_s}{ds}\norm{z_s - x^\ast}^2 \nonumber\\
&\quad + A_s \eta \dotprod{ \nabla f(x_s) ,z_s-x_s}+ B_s \eta' \dotprod{z_s-x^\ast,x_s - z_s}.
\label{eq:pl_sqc_eq_12}
\end{align}
Also,
\begin{align}
\mathbb{E}_{\xi}\!\left[\varphi(\overline{x}_s+ G(\overline{x}_s,\xi)) - \varphi(\overline{x}_s)\right]
&= \mathbb{E}_{\xi}\!\left[A_s(f(x_s - \gamma \nabla f(x_s,\xi)) - f(x_s))\right]\nonumber\\
&\quad + \frac{B_s}{2}\mathbb{E}_{\xi}\!\left[\norm{ z_s - \gamma'\nabla f(x_s,\xi) - x^\ast}^2 - \norm{z_s -x^\ast}^2\right]\nonumber\\
&= A_s\mathbb{E}_{\xi}\!\left[f(x_s - \gamma \nabla f(x_s,\xi)) - f(x_s)\right]
+ \frac{B_s \gamma^{'2}}{2}\mathbb{E}_{\xi}\!\left[\norm{\nabla f(x_s,\xi)}^2\right]\label{eq:pl_sqc_eq_3}\\
&\quad + B_s\gamma' \dotprod{\nabla f(x_s),x^\ast - z_s}, \label{eq:pl_sqc_eq_4}
\end{align}
where we used $\mathbb{E}_{\xi}[\nabla f(x_s,\xi)] = \nabla f(x_s)$.

As $f \in \ls$ and Assumption~\ref{ass:sgc} holds, by Lemma~\ref{lem:descent_sgc} (with $\gamma=\frac{1}{\rho L}$),
\[
A_s\mathbb{E}_{\xi}\!\left[f(x_s - \gamma \nabla f(x_s,\xi)) - f(x_s)\right]
\le -\frac{A_s}{2\rho L}\norm{\nabla f(x_s)}^2.
\]
Moreover, by the strong growth condition,
$\mathbb{E}_{\xi}\!\left[\norm{\nabla f(x_s,\xi)}^2\right]\le \rho \norm{\nabla f(x_s)}^2$. Hence
\begin{equation}\label{eq:pl_avg_pre_1_corr}
\eqref{eq:pl_sqc_eq_3}
\le \frac{1}{2}\Bigl(\rho B_s \gamma^{'2} - \frac{A_s}{\rho L}\Bigr)\norm{\nabla f(x_s)}^2.
\end{equation}

We use the identity
\begin{equation}\label{eq:pl_avg_pre_2}
\dotprod{z_s-x^\ast,x_s - z_s}
= \frac{1}{2}\norm{x_s-x^\ast}^2 -\frac{1}{2}\norm{z_s-x^\ast}^2-\frac{1}{2}\norm{x_s-z_s}^2
\le\frac{1}{2}\norm{x_s-x^\ast}^2 -\frac{1}{2}\norm{z_s-x^\ast}^2.
\end{equation}
Also,
\begin{align}\label{eq:pl_avg_pre_3_corr}
A_s \eta \dotprod{ \nabla f(x_s) ,z_s-x_s} +  B_s\gamma' \dotprod{\nabla f(x_s),x^\ast - z_s}
&= A_s \eta \dotprod{ \nabla f(x_s) ,x^\ast-x_s}
+ \bpar{A_s\eta - B_s \gamma'} \dotprod{ \nabla f(x_s) ,z_s-x^\ast}.
\end{align}

By Lemma~\ref{lem:scalar_prod_bound},
\[
\dotprod{ \nabla f(x_s) ,x^\ast-x_s}
\le  -a \sqrt{\frac{\mu}{L_0}}\bigl(f(x_s)-f^\ast\bigr)
      - \frac{a}{2}\sqrt{\mu\mu_0}\norm{x_s-x^\ast}^2.
\]

Gathering \eqref{eq:pl_sqc_eq_12}, \eqref{eq:pl_sqc_eq_4}, \eqref{eq:pl_avg_pre_1_corr}, \eqref{eq:pl_avg_pre_2},
and \eqref{eq:pl_avg_pre_3_corr}, we obtain
\begin{align*}
&\dotprod{\nabla \varphi(\overline{x}_s),b(\overline{x}_s)}
+ \mathbb{E}_{\xi}\!\left[\varphi(\overline{x}_s+ G(\overline{x}_s,\xi)) - \varphi(\overline{x}_s)\right] \\
&\le \frac{1}{2}\bpar{ \frac{dB_s}{ds}-B_s \eta'}\norm{z_s-x^\ast}^2
+  \frac{1}{2}\bpar{  B_s \eta' - A_s \eta\,a \sqrt{\mu\mu_0}}\norm{x_s-x^\ast}^2
+ \bpar{\frac{dA_s}{ds}- A_s \eta\,a \sqrt{\frac{\mu}{L_0}}}\bigl(f(x_s)-f^\ast\bigr)\\
&\quad + \bpar{A_s\eta - B_s \gamma'} \dotprod{ \nabla f(x_s) ,z_s-x^\ast}
+ \frac{1}{2}\bpar{\rho B_s \gamma^{'2} - \frac{A_s}{\rho L}}\norm{\nabla f(x_s)}^2.
\end{align*}

\noindent
\textbf{Parameter tuning \:}
Set $A_s = e^{\theta s}$ and $B_s = \sqrt{\mu_0 L_0}\,e^{\theta s}$, with
\[
\theta=\eta'=\frac{a}{\rho}\Bigl(\frac{\mu_0}{L_0}\Bigr)^{1/4}\sqrt{\frac{\mu}{L}},
\quad
\gamma=\frac{1}{\rho L},
\quad
\gamma'=\frac{1}{\rho}\frac{1}{(\mu_0 L_0)^{1/4}\sqrt{L}},
\quad
\eta=\frac{1}{\rho}\frac{(\mu_0 L_0)^{1/4}}{\sqrt{L}}.
\]
These parameters imply the cancellations
\begin{align}
A_s\eta - B_s \gamma' &= 0,\label{eq:cancel_param:0}\\
\rho B_s \gamma^{'2} - \frac{A_s}{\rho L} &= 0,\label{eq:cancel_param:1}\\
\frac{dB_s}{ds}-B_s \eta' &= 0,\label{eq:cancel_param:2}\\
\frac{dA_s}{ds}- A_s \eta\,a \sqrt{\frac{\mu}{L_0}} &= 0,\label{eq:cancel_param:3}\\
B_s \eta' - A_s \eta\,a \sqrt{\mu\mu_0} &= 0.\label{eq:cancel_param:4}
\end{align}
Therefore,
\[
\dotprod{\nabla \varphi(\overline{x}_s),b(\overline{x}_s)}
+ \mathbb{E}_{\xi}\!\left[\varphi(\overline{x}_s+ G(\overline{x}_s,\xi)) - \varphi(\overline{x}_s)\right] \le 0,
\]
so from \eqref{eq:prop_calc_sto},
\[
\varphi(\overline{x}_t) \le \varphi(\overline{x}_0) + M_t.
\]

We evaluate at $t=T_k$, take expectation, and use Theorem~\ref{thm:martingal_stopping} and
Proposition~\ref{prop:discretization_constant_param} to get
\[
\E{e^{\theta T_k}\bigl(f(\tilde x_k)-f^\ast\bigr)} \le \varphi(\overline{x}_0).
\]
Finally, as $x_0 = z_0$ and $f\in \qginf$, we have
\[
\varphi(\overline{x}_0)
= f(\tilde x_0)-f^\ast + \frac{\sqrt{\mu_0 L_0}}{2}\norm{\tilde x_0-x^\ast}^2
\le \Bigl(1+\sqrt{\frac{L_0}{\mu_0}}\Bigr)\,(f(\tilde x_0)-f^\ast),
\]
where we used $\tilde x_0 = \tilde z_0$,
and we obtain \eqref{eq:a.s.:2}, namely
\begin{equation*}   \mathbb{E}\left[e^{\beta T_k}(f(\tilde{x}_k)-f^\ast)\right] \le K_0,
         \end{equation*}with $\beta=\theta$ and $K_0=\Bigl(1+\sqrt{\frac{L_0}{\mu_0}}\Bigr)(f(\tilde x_0)-f^\ast)$.
By Lemma~\ref{lem:conv_high_proba}, with probability $1-\frac{1}{c_0}-\exp(-(c_1-1-\log c_1)k)$,
\[
f(\tilde x_k)-f^\ast
\le c_0\Bigl(1+\sqrt{\frac{L_0}{\mu_0}}\Bigr)(f(\tilde x_0)-f^\ast)\,
\exp\!\Bigl(-\theta(1-c_1)k\Bigr),
\]
which is the claim.
\end{proof}


\section{Trajectory-adapted convergence}\label{app:adptative_condition}
In this section, we address a setting in which we assume the aiming condition is verified on average along the trajectory (Assumption~\ref{ass:avg_corelation}).
We start by proving Theorem~\ref{thm:gf_pl_&_avg_cor}, that addresses the dynamical system case. Then, we will state and prove similar results in the discrete case.
  \subsection{Proof of Theorem~\ref{thm:gf_pl_&_avg_cor}}\label{app:gf_pl_&_avg_cor}
\paragraph{Proof with Gradient Flow \eqref{eq:gf}}

Let
\[
E_t = e^{a_{avg}\sqrt{\mu \mu_0}\, t}\left[f(x_t)-f^\ast + \frac{\beta}{2}\norm{x_t-x^\ast}^2 \right].
\]
Then
\begin{align*}
\dot E_t
&= a_{avg}\sqrt{\mu\mu_0}\,e^{a_{avg}\sqrt{\mu\mu_0}t}
\left[f(x_t)-f^\ast + \frac{\beta}{2}\norm{x_t-x^\ast}^2 \right]\\
&\quad + e^{a_{avg}\sqrt{\mu\mu_0}t}\left[\dotprod{\nabla f(x_t),\dot x_t}
+ \beta \dotprod{x_t-x^\ast,\dot x_t}\right]\\
&\overset{(i)}{=} a_{avg}\sqrt{\mu\mu_0}\,e^{a_{avg}\sqrt{\mu\mu_0}t}
\left[f(x_t)-f^\ast + \frac{\beta}{2}\norm{x_t-x^\ast}^2 \right]\\
&\quad + e^{a_{avg}\sqrt{\mu\mu_0}t}\left[-\norm{\nabla f(x_t)}^2
- \beta \dotprod{x_t-x^\ast,\nabla f(x_t)}\right]\\
&\le a_{avg}\sqrt{\mu\mu_0}\,e^{a_{avg}\sqrt{\mu\mu_0}t}
\left[f(x_t)-f^\ast + \frac{\beta}{2}\norm{x_t-x^\ast}^2 \right]
- \beta e^{a_{avg}\sqrt{\mu\mu_0}t}\dotprod{x_t-x^\ast,\nabla f(x_t)},
\end{align*}
where (i) follows from $\dot x_t=-\nabla f(x_t)$.

Integrating from $0$ to $t$ yields
\begin{equation}\label{eq:gf_ac_average_1_corr}
E_t \le E_0
+ \int_0^t a_{avg}\sqrt{\mu\mu_0}\,e^{a_{avg}\sqrt{\mu\mu_0}s}
\left[f(x_s)-f^\ast + \frac{\beta}{2}\norm{x_s-x^\ast}^2 \right]ds
- \int_0^t \beta e^{a_{avg}\sqrt{\mu\mu_0}s}\dotprod{x_s-x^\ast,\nabla f(x_s)}\,ds.
\end{equation}

By Assumption~\ref{ass:avg_corelation} (with $\theta=a_{avg}\sqrt{\mu\mu_0}$), we have
\[
\int_0^t e^{\theta s}\dotprod{\nabla f(x_s),x_s-x^\ast}\,ds
\ge a_{avg}\int_0^t e^{\theta s}\norm{\nabla f(x_s)}\norm{x_s-x^\ast}\,ds,
\]
hence
\[
\int_0^t \beta e^{\theta s}\dotprod{x_s-x^\ast,\nabla f(x_s)}\,ds
\ge
\int_0^t \beta a_{avg} e^{\theta s}\norm{x_s-x^\ast}\norm{\nabla f(x_s)}\,ds.
\]
Plugging into \eqref{eq:gf_ac_average_1_corr} gives
\begin{equation}\label{eq:gf_ac_average_2_corr}
E_t \le E_0
+ \int_0^t a_{avg}\sqrt{\mu\mu_0}\,e^{\theta s}
\left[f(x_s)-f^\ast + \frac{\beta}{2}\norm{x_s-x^\ast}^2 \right]ds
- \int_0^t \beta a_{avg} e^{\theta s}\norm{x_s-x^\ast}\norm{\nabla f(x_s)}\,ds.
\end{equation}

Using Lemma~\ref{lem:prod}, for all $s\ge 0$,
\[
-\norm{x_s-x^\ast}\norm{\nabla f(x_s)}
\le -\sqrt{\frac{\mu}{L_0}}(f(x_s)-f^\ast) - \frac{\sqrt{\mu \mu_0}}{2}\norm{x_s-x^\ast}^2.
\]
Multiplying by $\beta a_{avg}e^{\theta s}$ and inserting into \eqref{eq:gf_ac_average_2_corr}, we obtain
\begin{align*}
E_t
&\le E_0 + \int_0^t e^{\theta s}\Bigl[
a_{avg}\sqrt{\mu\mu_0}\,(f(x_s)-f^\ast)
- \beta a_{avg}\sqrt{\frac{\mu}{L_0}}\,(f(x_s)-f^\ast)\\
&\hspace{2.2cm}
+ \frac{\beta}{2}a_{avg}\sqrt{\mu\mu_0}\,\norm{x_s-x^\ast}^2
- \frac{\beta}{2}a_{avg}\sqrt{\mu\mu_0}\,\norm{x_s-x^\ast}^2
\Bigr]ds\\
&= E_0 + a_{avg}\int_0^t e^{\theta s}
\Bigl(\sqrt{\mu\mu_0}-\beta\sqrt{\frac{\mu}{L_0}}\Bigr)\,(f(x_s)-f^\ast)\,ds.
\end{align*}
Choose $\beta=\sqrt{\mu_0L_0}$, so that
$\sqrt{\mu\mu_0}-\beta\sqrt{\mu/L_0}=0$. Then $E_t\le E_0$ for all $t\ge 0$.
Since $f(x_t)-f^\ast \le e^{-\theta t}E_t$, we get
\[
f(x_t)-f^\ast \le e^{-\theta t}E_0
= e^{-a_{avg}\sqrt{\mu\mu_0}\,t}\Bigl(f(x_0)-f^\ast + \frac{\beta}{2}\norm{x_0-x^\ast}^2\Bigr).
\]
Finally, using $f(x_0)-f^\ast \ge \frac{\mu_0}{2}\norm{x_0-x^\ast}^2$ and $\beta=\sqrt{\mu_0L_0}$,
\[
\frac{\beta}{2}\norm{x_0-x^\ast}^2
\le \sqrt{\frac{L_0}{\mu_0}}\,(f(x_0)-f^\ast),
\]
hence
\[
f(x_t)-f^\ast \le \Bigl(1+\sqrt{\frac{L_0}{\mu_0}}\Bigr)e^{-a_{avg}\sqrt{\mu\mu_0}\,t}\,(f(x_0)-f^\ast).
\]


% Using Lemma~\ref{lem:lyapunov_0_bound}, we conclude that
% $$f(x_t)-f^\ast \le e^{-a_{avg}\sqrt{\mu \mu_0} t}E_0 = e^{-a_{avg}\sqrt{\mu \mu_0} t}\frac{\sqrt{L_0}(\sqrt{L_0}- \sqrt{\mu_0})}{2}\norm{x_0-x^\ast}^2.$$
% By \eqref{eq:pl_corr_nest_2} and \eqref{eq:pl_corr_nest_3}, we have
% \begin{align}
%     \dot E_t \le -\beta  a\sqrt{\frac{\mu_0}{L_0}}(f(x_t)-f^\ast) - \frac{\beta}{2}a\sqrt{\mu \mu_0}\norm{x_t-x^\ast}^2.
% \end{align}
% To obtain a linear rate, we want
% $$ \beta a \sqrt{\frac{\mu_0}{L_0}} = a\sqrt{\mu \mu_0},$$
% which imposes $\beta = \sqrt{\mu_0 L_0}$. We then deduce
% $$\dot E_t \le -a\sqrt{\mu \mu_0}\bpar{f(x_t)-f^\ast + \frac{\beta}{2}\norm{x_t-x^\ast}^2},  $$
% and by integration
% $$E_t \le e^{-a\sqrt{\mu \mu_0}}E_0.$$


% \paragraph{Proof with Nesterov ODE momentum \eqref{eq:nmo}}

% Let be
% \begin{equation}
%     E_t = e^{\theta t}(f(x_t)-f^\ast) +\beta \frac{ e^{\theta t}}{2}\norm{z_t -  x^\ast}^2
% \end{equation}
% The beginning of the proof is similar with the proof of Theorem~\ref{thm:pl_accel} (i). In particular, choosing $\eta' = \theta$ and $\eta = \beta\gamma'$, we get \eqref{eq:pl_corr_nest_1_alt_simpl}, namely
% \begin{align}\label{eq:pl_corr_avg_nest_1}
%  \dot{E}_t \le   e^{\theta t}\bpar{ \theta(f(x_t)-f^\ast) +\eta\dotprod{\nabla f(x_t), x^\ast - x_t}+  \frac{\beta  \theta }{2}\norm{x^\ast-x_t}^2}.
% \end{align}
% We have
% $$E_t = E_0 + \int_0^t \dot E_sds,$$
% in which we can plug \eqref{eq:pl_corr_avg_nest_1} to deduce
% \begin{equation}
% \begin{aligned}\label{eq:pl_corr_avg_nest_2}
%     E_t &\le  E_0 + \int_0^t  e^{\theta s}\bpar{ \theta(f(x_s)-f^\ast) +\eta\dotprod{\nabla f(x_s), x^\ast - x_s}+  \frac{\beta  \theta }{2}\norm{x^\ast-x_s}^2}ds\\
%     &= E_0 + \int_0^t  e^{\theta s}\bpar{ \theta(f(x_s)-f^\ast)+  \frac{\beta  \theta }{2}\norm{x^\ast-x_s}^2}ds + \eta \int_0^te^{\theta s} \dotprod{\nabla f(x_s), x^\ast - x_s}ds.
% \end{aligned}
% \end{equation}
% By Assumption~\ref{ass:avg}, we have
% $$ \int_0^t  e^{\theta s} \dotprod{x_s-x^\ast,\nabla f(x_s)}ds \ge \int_0^t  a_{avg} e^{\theta s} \norm{x_s-x^\ast}\norm{\nabla f(x_s)}ds. $$
% such that \eqref{eq:pl_corr_avg_nest_2} becomes
% \begin{equation}\label{eq:pl_corr_avg_nest_3}
%     E_0 + \int_0^t  e^{\theta s}\bpar{ \theta(f(x_s)-f^\ast)+  \frac{\beta  \theta }{2}\norm{x^\ast-x_s}^2}ds - \eta a_{avg} \int_0^te^{\theta s} \norm{x_s-x^\ast}\norm{\nabla f(x_s)}ds.
% \end{equation}
% Using Lemma~\ref{lem:prod}, we have
%         $$  -\norm{x_s-x^\ast}\norm{\nabla f(x_s)} \le -\sqrt{\frac{\mu}{L_0}}(f(x_s)-f^\ast) - \frac{\sqrt{\mu \mu_0}}{2}\norm{x_s-x^\ast}^2,$$
% which we use in \eqref{eq:pl_corr_avg_nest_3} to obtain 
% \begin{equation}\label{eq:pl_corr_avg_nest_4}
%      E_t \le E_0 + \int_0^t e^{\theta s}\left[ \bpar{\theta - \eta a_{avg}\sqrt{\frac{\mu}{L_0}}}(f(x_s)-f^\ast) + \frac{1}{2}\bpar{\beta \theta - \eta a_{avg}\sqrt{\mu \mu_0}} \right]ds.
% \end{equation}
% Fixing
% \begin{equation*}
%     \theta = \eta' = a_{avg}\bpar{\frac{\mu_0}{L_0}}^{\frac{1}{4}}\sqrt{\mu},\quad \gamma' =  (\mu_0 L_0)^{-1/4}, \quad \eta = (\mu_0 L_0)^{1/4}, \quad \gamma \ge 0, \quad \beta = \sqrt{\mu_0 L_0},
% \end{equation*}
% we cancel all the term inside the integral in \eqref{eq:pl_corr_avg_nest_4}, and we have
% $$E_t \le E_0.$$
% As $E_t \ge e^{ a_{avg}\bpar{\frac{\mu_0}{L_0}}^{\frac{1}{4}}\sqrt{\mu}}(f(x_t)-f^\ast)$, we deduce
% $$f(x_t)-f^\ast \le e^{- a_{avg}\bpar{\frac{\mu_0}{L_0}}^{\frac{1}{4}}\sqrt{\mu}}E_0= e^{- a_{avg}\bpar{\frac{\mu_0}{L_0}}^{\frac{1}{4}}\sqrt{\mu}}\bpar{f(x_0)-f^\ast + \frac{\sqrt{\mu_0 L_0}}{2}\norm{x_0-x^\ast}^2}.$$
% Using $f(x_0)-f^\ast \ge \frac{\mu_0}{2}\norm{x_0-x^\ast}^2$, we conclude that
% $$f(x_t)-f^\ast \le e^{- a_{avg}\bpar{\frac{\mu_0}{L_0}}^{\frac{1}{4}}\sqrt{\mu}} \bpar{1+\sqrt\frac{L_0}{\mu_0}}(f(x_0)-f^\ast)$$







\paragraph{Proof with Nesterov ODE momentum \eqref{eq:nmo}}

Let
\[
E_t = e^{\theta t}\Bigl(f(x_t)-f^\ast\Bigr)+\frac{\beta}{2}e^{\theta t}\norm{z_t-x^\ast}^2.
\]
As in the proof of Theorem~\ref{thm:pl_accel} (i), choose $\eta'=\theta$ and $\eta=\beta\gamma'$.
Then we get (see \eqref{eq:pl_corr_nest_1_alt_simpl})
\begin{equation}\label{eq:pl_corr_avg_nest_1_corr}
 \dot{E}_t \le e^{\theta t}\Bigl(
 \theta(f(x_t)-f^\ast)
 +\eta\dotprod{\nabla f(x_t), x^\ast - x_t}
 +\frac{\beta\theta}{2}\norm{x_t-x^\ast}^2
 \Bigr).
\end{equation}

Integrating,
\begin{align}
E_t
&\le E_0 + \int_0^t e^{\theta s}\Bigl(
 \theta(f(x_s)-f^\ast)
 +\frac{\beta\theta}{2}\norm{x_s-x^\ast}^2
 \Bigr)\,ds
 + \eta\int_0^t e^{\theta s}\dotprod{\nabla f(x_s),x^\ast-x_s}\,ds \nonumber\\
&= E_0 + \int_0^t e^{\theta s}\Bigl(
 \theta(f(x_s)-f^\ast)
 +\frac{\beta\theta}{2}\norm{x_s-x^\ast}^2
 \Bigr)\,ds
 - \eta\int_0^t e^{\theta s}\dotprod{\nabla f(x_s),x_s-x^\ast}\,ds.\label{eq:pl_corr_avg_nest_2_corr}
\end{align}

By Assumption~\ref{ass:avg_corelation},
\[
\int_0^t e^{\theta s}\dotprod{\nabla f(x_s),x_s-x^\ast}\,ds
\ge a_{avg}\int_0^t e^{\theta s}\norm{\nabla f(x_s)}\norm{x_s-x^\ast}\,ds,
\]
hence
\begin{equation}\label{eq:pl_corr_avg_nest_3_corr}
E_t
\le E_0 + \int_0^t e^{\theta s}\Bigl(
 \theta(f(x_s)-f^\ast)
 +\frac{\beta\theta}{2}\norm{x_s-x^\ast}^2
 \Bigr)\,ds
 - \eta a_{avg}\int_0^t e^{\theta s}\norm{\nabla f(x_s)}\norm{x_s-x^\ast}\,ds.
\end{equation}

Using Lemma~\ref{lem:prod},
\[
-\norm{x_s-x^\ast}\norm{\nabla f(x_s)}
\le -\sqrt{\frac{\mu}{L_0}}(f(x_s)-f^\ast)
 -\frac{\sqrt{\mu\mu_0}}{2}\norm{x_s-x^\ast}^2,
\]
and inserting into \eqref{eq:pl_corr_avg_nest_3_corr}, we obtain
\begin{equation}\label{eq:pl_corr_avg_nest_4_corr}
E_t
\le E_0 + \int_0^t e^{\theta s}\left[
\Bigl(\theta-\eta a_{avg}\sqrt{\frac{\mu}{L_0}}\Bigr)(f(x_s)-f^\ast)
+\frac{1}{2}\Bigl(\beta\theta-\eta a_{avg}\sqrt{\mu\mu_0}\Bigr)\norm{x_s-x^\ast}^2
\right]ds.
\end{equation}

Choose $\beta=\sqrt{\mu_0L_0}$ and impose the cancellations
\[
\theta=\eta a_{avg}\sqrt{\frac{\mu}{L_0}}
\qquad\text{and}\qquad
\beta\theta=\eta a_{avg}\sqrt{\mu\mu_0}.
\]
These two equalities are equivalent to $\beta=\sqrt{\mu_0L_0}$ and
\[
\theta=a_{avg}\Bigl(\frac{\mu_0}{L_0}\Bigr)^{1/4}\sqrt{\mu},
\qquad
\eta=(\mu_0L_0)^{1/4}.
\]
With $\eta'=\theta$ and $\eta=\beta\gamma'$, we can take
\[
\gamma'=\frac{\eta}{\beta}=(\mu_0L_0)^{-1/4},
\qquad
\eta'=\theta=a_{avg}\Bigl(\frac{\mu_0}{L_0}\Bigr)^{1/4}\sqrt{\mu},
\qquad
\gamma\ge 0.
\]
Then the integrand in \eqref{eq:pl_corr_avg_nest_4_corr} is identically zero and we get $E_t\le E_0$.
Since $f(x_t)-f^\ast \le e^{-\theta t}E_t$, we deduce
\[
f(x_t)-f^\ast \le e^{-\theta t}E_0
= e^{-\theta t}\Bigl(f(x_0)-f^\ast+\frac{\beta}{2}\norm{z_0-x^\ast}^2\Bigr).
\]
As $z_0=x_0$, then
\[
f(x_t)-f^\ast \le e^{-\theta t}\Bigl(f(x_0)-f^\ast+\frac{\sqrt{\mu_0L_0}}{2}\norm{x_0-x^\ast}^2\Bigr).
\]
Using $f(x_0)-f^\ast \ge \frac{\mu_0}{2}\norm{x_0-x^\ast}^2$, we obtain
\[
f(x_t)-f^\ast
\le \Bigl(1+\sqrt{\frac{L_0}{\mu_0}}\Bigr)e^{-\theta t}\,(f(x_0)-f^\ast),
\qquad
\theta=a_{avg}\Bigl(\frac{\mu_0}{L_0}\Bigr)^{1/4}\sqrt{\mu}.
\]



\subsection{Discrete case}\label{app:pl_accel_avg}
We address the discrete case. As in Appendix~\ref{app:cv_results}, we consider stochastic gradients under Assumption~\ref{ass:sgc}.  
\begin{assumption}\label{ass:avg}
  For some $\{ A_i \}_{i =1,\dots,k}$ and $a_{avg} > 0$, $\{\tilde y_k \}_{k\ge 0}$ is such that
    $$\sum_{i = 1}^k A_i  \left[\dotprod{\nabla
     f(\tilde y_i,\xi_i),\tilde y_i-x^\ast}-a_{avg} \norm{\nabla
     f(\tilde y_i,\xi_i)}\norm{x^\ast - \tilde y_i}\right] \ge 0.$$
\end{assumption}
We consider \eqref{eq:nm}, a similar and simpler analysis can be carried for \eqref{eq:gd}.
% \begin{theorem}\label{thm:pl_accel_avg}
% Let $f\in \pl  \cap \ls$.  Denote $L_0 := \inf \{L'\le L:  f \in \qgf{L'} \}$, $\mu_0 := \sup\{\mu' \ge \mu : f \in \qginff{\mu'}\}$.
% Let $\xdisctilde$ verify \eqref{eq:nm} with the \hyperlink{cont_param}{continuized parameterization} and with $\nabla f(\tilde x_k)$ replaced by $\nabla f(\tilde x_k,\xi_k)$ for random variables $\xi_k$ independent of $\tilde x_k$, under Assumption~\ref{ass:sgc}  and such that Assumption~\ref{ass:avg_corelation} is verified with $ \theta = \frac{a_{avg}}{\rho}\bpar{\frac{\mu_0}{L_0}}^{1/4}\sqrt{\frac{\mu}{L}}$ and $A_i = e^{\theta T_i}$. Let
% $ \eta' = \frac{a_{avg}}{\rho}\bpar{\frac{\mu_0}{L_0}}^{1/4}\sqrt{\frac{\mu}{L}},\quad \gamma = \frac{1}{\rho L},\quad \gamma' = \frac{1}{\rho( \mu_0 L_0)^{1/4}\sqrt{ L}}$, $\eta = \frac{1}{\rho}\frac{(\mu_0 L_0)^{1/4}}{\sqrt{L}}$, Then,  
%     with probability  $1-\frac{1}{c_0} - e^{-(c_1-1-\log(c_1))k}$, for some $c_0 > 1$ and $c_1 \in (0,1)$ we have
% $$ f(\tilde x_k)-f^\ast \le \bpar{1+\sqrt\frac{L_0}{\mu_0}}(f(x_0)-f^\ast)e^{-\frac{a_{avg}}{\rho}\bpar{\frac{\mu_0}{L_0}}^{1/4}\sqrt{\frac{\mu}{L}}(1-c_1)k}.$$

% \end{theorem}
% \begin{proof}
% %\jh{replace $A_s$ by $e^{\theta s}$ in this and previous ?}
% The beginning of the proof can be written identically as for Theorem~\ref{thm:pl_accel} (ii), proven in the previous Appendix~\ref{app:pl_accel_avg}. So, we start by combining \eqref{eq:sqc_eq_1}-\eqref{eq:pl_avg_pre_3}, which leads to 
% \begin{align*}
%     &\dotprod{\nabla \varphi(x_s),b(x_s)} +  \mathbb{E}_{\xi} \varphi(\overline{x}_s+ G(\overline{x}_s)) - \varphi(\overline{x}_s) \\
%     &\le \frac{1}{2}\bpar{ \frac{dB_s}{ds}-B_s \eta_s'}\norm{z_s-x^\ast}^2 +  \frac{B_s \eta'_s }{2} \norm{x_s-x^\ast}^2 + \frac{dA_s}{ds}(f(x_s)-f^\ast) +A_s \eta_s \dotprod{\nabla
%      f(x_s),x^\ast -x_s}\\
%     & + \bpar{A_s\eta_s - B_s \gamma'_s} \dotprod{ \nabla f(x_s) ,z_s-x_s}+\bpar{\rho B_s \gamma_s^{'2} + A_s \gamma_s( \rho\gamma_s L-2)}\frac{1}{2}\norm{\nabla f(x_s)}^2.
% \end{align*}
% From \eqref{eq:prop_calc_sto}, we have
% \begin{equation}
%     \begin{aligned}\label{eq:avg_continuized:0}
%         \varphi(t,x_t,z_t) &= \varphi(0,x_0,z_0) + \int_0^t \left[  \frac{1}{2}\bpar{ \frac{dB_s}{ds}-B_s \eta_s'}\norm{z_s-x^\ast}^2 +  \frac{B_s \eta'_s }{2} \norm{x_s-x^\ast}^2 + \frac{dA_s}{ds}(f(x_s)-f^\ast)\right. \\
%         &+ \left. \bpar{A_s\eta_s - B_s \gamma'_s} \dotprod{ \nabla f(x_s) ,z_s-x_s}+\bpar{\rho B_s \gamma_s^{'2} + A_s \gamma_s( \rho\gamma_s L-2)}\frac{1}{2}\norm{\nabla f(x_s)}^2 \right]ds\\
%         &+ \int_0^t A_s \eta_s \dotprod{\nabla
%      f(x_s),x^\ast -x_s}ds + M_t. 
%     \end{aligned}
% \end{equation}
% We have
% \begin{equation}\label{eq:avg_continuized:1}
%     \int_0^t A_s \eta_s \dotprod{\nabla
%      f(x_s),x^\ast -x_s}ds  = \int_0^t A_s \eta_s \dotprod{\nabla
%      f(x_s),x^\ast -x_s}dNs+\underbrace{\int_0^t A_s \eta_s \dotprod{\nabla
%      f(x_s),x^\ast -x_s}(ds-dN_s)}_{:=-U_t},  
% \end{equation}
% where $-U_t$ is a martingale with $\E{-U_t} =0$, for all $t\in \R_+$. By Assumption~\ref{ass:avg}, we have
% $$\int_0^t A_s  \dotprod{\nabla
%      f(x_s),x_s-x^\ast}dN_s \ge \int_0^t A_s a_{avg} \norm{\nabla
%      f(x_s)}\norm{x^\ast -x_s}dN_s.$$
%      Note also that
%      \begin{equation}\label{eq:avg_continuized:2}
%          \int_0^t A_s a_{avg} \norm{\nabla
%      f(x_s)}\norm{x^\ast -x_s}dN_s = \int_0^t A_s a_{avg} \norm{\nabla
%      f(x_s)}\norm{x^\ast -x_s}ds + \underbrace{\int_0^t A_s a_{avg} \norm{\nabla
%      f(x_s)}\norm{x^\ast -x_s}(dN_s-ds)}_{:= -V_t},
%      \end{equation}
%      with $-V_t$ is a martingale satisfying $\E{-V_t} = 0$, for all $t\in \R_+$. Then, using Lemma~\ref{lem:prod}, we have
%      \begin{equation}\label{eq:avg_continuized:3}
%          \int_0^t A_s a_{avg} \norm{\nabla
%      f(x_s)}\norm{x^\ast -x_s}dN_s \ge \int_0^t A_s a_{avg}\bpar{\sqrt{\frac{\mu}{L_0}}(f(x)-f^\ast) + \frac{\sqrt{\mu \mu_0}}{2}\norm{x-x^\ast}^2}dN_s
%      \end{equation}
%      So, combining \eqref{eq:avg_continuized:0}-\eqref{eq:avg_continuized:3}, we have
%      \begin{equation}
%     \begin{aligned}\label{eq:avg_continuized:4}
%         \varphi(t,x_t,z_t) &= \varphi(0,x_0,z_0) + \int_0^t \left[  \frac{1}{2}\bpar{ \frac{dB_s}{ds}-B_s \eta_s'}\norm{z_s-x^\ast}^2 +  \frac{1}{2}\bpar{B_s \eta'_s - A_s a_{avg}\sqrt{\mu \mu_0}}\norm{x_s-x^\ast}^2 \right. \\&+ \bpar{\frac{dA_s}{ds}- A_s a_{avg}\sqrt{\frac{\mu}{L_0}}}(f(x_s)-f^\ast)
%         +  \bpar{A_s\eta_s - B_s \gamma'_s} \dotprod{ \nabla f(x_s) ,z_s-x_s}\\
%         & \left.+\bpar{\rho B_s \gamma_s^{'2}  + A_s \gamma_s( \rho\gamma_s L-2)}\frac{1}{2}\norm{\nabla f(x_s)}^2 \right]ds+ \underbrace{U_t + V_t + M_t}_{:= M'_t},
%     \end{aligned}
% \end{equation}
% where $M'_t$ is a martingale such that $\E{M'_t} = 0$ for all $t \in \R_+$. In order to cancel the integrand in \eqref{eq:avg_continuized:4}, we can perform the same computation as for Theorem~\ref{thm:pl_accel} (i), proved in Appendix~\ref{app:cv_results}. Precisely, fixing $A_s = e^{\theta s}$, $B_s = \sqrt{\mu_0 L_0}e^{\theta_s}$, with $ \theta = \eta' = \frac{a_{avg}}{\rho}\bpar{\frac{\mu_0}{L_0}}^{1/4}\sqrt{\frac{\mu}{L}},\quad \gamma = \frac{1}{\rho L},\quad \gamma' = \frac{1}{\rho} \frac{1}{( \mu_0 L_0)^{1/4}\sqrt{ L}}, $ and $\eta = \frac{1}{\rho}\frac{(\mu_0 L_0)^{1/4}}{\sqrt{L}}$, one can verify that we have \eqref{eq:cancel_param:0}-\eqref{eq:cancel_param:4}, replacing $a$ by $a_{avg}$. This set of parameters implies $\dotprod{\nabla \varphi(x_s),b(x_s)} +  \mathbb{E}_{\xi} \varphi(\overline{x}_s+ G(\overline{x}_s)) - \varphi(\overline{x}_s) \le 0$, which in turn implies
% \begin{equation}\label{eq:pl_corr_cont_supermart_avg}
%     \phi(t) \le \phi(0) + M'_t.
% \end{equation}
% Then, we evaluate \eqref{eq:pl_corr_cont_supermart_avg} at $t = T_k$, take expectation and use Theorem~\ref{thm:martingal_stopping} and Proposition~\ref{prop:discretization_constant_param}, to get
% \begin{equation}
% \E{e^{ \frac{a_{avg}}{\rho}\bpar{\frac{\mu_0}{L_0}}^{1/4}\sqrt{\frac{\mu}{L}}T_k}f(\tilde x_k)-f^\ast)} \le \phi(0).
% \end{equation}
% Using $f(x_0)-f^\ast \ge \frac{\mu_0}{2}\norm{x_0-x^\ast}^2$, we deduce $\phi(0) = f(\tilde x_0)-f^\ast + \frac{\sqrt{\mu_0 L_0}}{2}\norm{\tilde x_0-x^\ast}^2 \le \bpar{1+\sqrt\frac{L_0}{\mu_0}}(f(x_0)-f^\ast)$
% %by Lemma~\ref{lem:lyapunov_0_bound}. 
% So, we have
% $$\E{e^{ \frac{a_{avg}}{\rho}\bpar{\frac{\mu_0}{L_0}}^{1/4}\sqrt{\frac{\mu}{L}}T_k}f(\tilde x_k)-f^\ast)} \le  \bpar{1+\sqrt\frac{L_0}{\mu_0}}(f(x_0)-f^\ast). $$ 
% Finally, by Lemma~\ref{lem:conv_high_proba}, with probability  $1-\frac{1}{c_0} - e^{-(c_1-1-\log(c_1))k}$, for some $c_0 > 1$ and $c_1 \in (0,1)$ we have
% $$ f(\tilde x_k)-f^\ast \le c_0 \bpar{1+\sqrt\frac{L_0}{\mu_0}}(f(x_0)-f^\ast)e^{-\frac{a_{avg}}{\rho}\bpar{\frac{\mu_0}{L_0}}^{1/4}\sqrt{\frac{\mu}{L}}(1-c_1)k}. $$
% \end{proof}




\begin{theorem}\label{thm:pl_accel_avg}
Let $f\in \pl \cap \ls$. Denote
$L_0 := \inf \{L'\le L:  f \in \qgf{L'} \}$,
$\mu_0 := \sup\{\mu' \ge \mu : f \in \qginff{\mu'}\}$.
Let $\xdisctilde$ verify \eqref{eq:nm} with the \hyperlink{cont_param}{continuized parameterization} and with
$\nabla f(\tilde x_k)$ replaced by $\nabla f(\tilde x_k,\xi_k)$ for random variables $\xi_k$ independent of $\tilde x_k$,
under Assumption~\ref{ass:sgc}, and assume that Assumption~\ref{ass:avg} holds with $A_i = e^{\theta T_i}$, where
\[
\theta = \eta' = \frac{a_{avg}}{\rho}\Bigl(\frac{\mu_0}{L_0}\Bigr)^{1/4}\sqrt{\frac{\mu}{L}}.
\]
Let
\[
\gamma = \frac{1}{\rho L},
\quad
\gamma' = \frac{1}{\rho(\mu_0 L_0)^{1/4}\sqrt{L}},
\quad
\eta = \frac{1}{\rho}\frac{(\mu_0 L_0)^{1/4}}{\sqrt{L}}.
\]
Then, with probability $1-\frac{1}{c_0} - e^{-(c_1-1-\log(c_1))k}$ (for some $c_0>1$, $c_1\in(0,1)$),
\[
f(\tilde x_k)-f^\ast \le c_0\Bigl(1+\sqrt{\frac{L_0}{\mu_0}}\Bigr)(f(\tilde x_0)-f^\ast)\,
\exp\!\Bigl(-\theta(1-c_1)k\Bigr).
\]
\end{theorem}

\begin{proof}
We use the same notations as in the proof of Theorem~\ref{thm:sto_pl_accel}.
Let $\overline{x}_t=(t,x_t,z_t)$ and define
\[
\varphi(t,x,z)
=
A_t(f(x)-f^\ast)+\frac{B_t}{2}\norm{z-x^\ast}^2,
\qquad
A_t=e^{\theta t},
\quad
B_t=\sqrt{\mu_0L_0}\,e^{\theta t}.
\]

Applying Proposition~\ref{prop:sto_calc}, we have
\begin{equation}\label{eq:avg_prop_calc}
\varphi(\overline{x}_t)
=
\varphi(\overline{x}_0)
+
\int_0^t
\Bigl[
\dotprod{\nabla \varphi(\overline{x}_s),b(\overline{x}_s)}
+
\mathbb{E}_{\xi}\!\bigl(
\varphi(\overline{x}_s+G(\overline{x}_s,\xi))-\varphi(\overline{x}_s)
\bigr)
\Bigr]ds
+M_t,
\end{equation}
where $M_t$ is a martingale with $\E{M_t}=0$.

From the same algebraic computations as in
Theorem~\ref{thm:pl_accel} (ii),
we obtain
\begin{align}\label{eq:avg_drift_1}
&\dotprod{\nabla \varphi(\overline{x}_s),b(\overline{x}_s)}
+
\mathbb{E}_{\xi}\!\bigl(
\varphi(\overline{x}_s+G(\overline{x}_s,\xi))-\varphi(\overline{x}_s)
\bigr)
\\
&\le
\frac{1}{2}\Bigl(\frac{dB_s}{ds}-B_s\eta'\Bigr)\norm{z_s-x^\ast}^2
+
\frac{B_s\eta'}{2}\norm{x_s-x^\ast}^2
+
\frac{dA_s}{ds}(f(x_s)-f^\ast)
\nonumber\\
&\quad
+
A_s\eta\,\dotprod{\nabla f(x_s),x^\ast-x_s}
+
(A_s\eta-B_s\gamma')\dotprod{\nabla f(x_s),z_s-x^\ast}
\nonumber\\
&\quad
+
\frac12\Bigl(\rho B_s\gamma'^2-\frac{A_s}{\rho L}\Bigr)\norm{\nabla f(x_s)}^2.
\nonumber
\end{align}

With the parameter choices of the theorem and $B_s=\sqrt{\mu_0L_0}A_s$, we have
\[
A_s\eta-B_s\gamma'=0,
\qquad
\rho B_s\gamma'^2-\frac{A_s}{\rho L}=0,
\qquad
\frac{dB_s}{ds}-B_s\eta'=0.
\]
Thus \eqref{eq:avg_drift_1} reduces to
\begin{equation}\label{eq:avg_drift_2}
\dotprod{\nabla \varphi(\overline{x}_s),b(\overline{x}_s)}
+
\mathbb{E}_{\xi}\!\bigl(
\varphi(\overline{x}_s+G(\overline{x}_s,\xi))-\varphi(\overline{x}_s)
\bigr)
\le
\frac{dA_s}{ds}(f(x_s)-f^\ast)
+
\frac{B_s\eta'}{2}\norm{x_s-x^\ast}^2
+
A_s\eta\,\dotprod{\nabla f(x_s),x^\ast-x_s}.
\end{equation}

We now control the last term using the jump structure.
For any predictable process $h_s$,
\[
\int_0^t h_s\,dN_s
=
\int_0^t h_s\,ds
+
\int_0^t h_s\,(dN_s-ds),
\]
where $\int_0^t h_s(dN_s-ds)$ is a martingale.

Applying this identity with
\[
h_s
=
A_s\eta\Bigl[
\dotprod{\nabla f(x_s,\xi_s),x_s-x^\ast}
-
a_{avg}\norm{\nabla f(x_s,\xi_s)}\norm{x_s-x^\ast}
\Bigr],
\]
and using Assumption~\ref{ass:avg} with $A_i=e^{\theta T_i}$, we obtain for $t=T_k$
\begin{equation}\label{eq:avg_jump_control}
\int_0^{T_k} A_s\eta\,\dotprod{\nabla f(x_s),x^\ast-x_s}\,dN_s
\le
-\eta a_{avg}\int_0^{T_k} A_s\norm{\nabla f(x_s)}\norm{x_s-x^\ast}\,dN_s.
\end{equation}
This holds because by definition of the Poisson integrals, we have
\[\int_0^{T_k} A_s\eta\,\dotprod{\nabla f(x_s),x^\ast-x_s}\,dN_s = \sum_{i=1}^k A_k\eta \dotprod{\nabla f(x_{T_k^-}),x^\ast-x_{T_k^-}} = \sum_{i=1}^k A_k\eta \dotprod{\nabla f(\tilde y_k),x^\ast-\tilde y_k},\]
where the last inequality follows from Proposition~\ref{prop:discretization_constant_param}.

Using Lemma~\ref{lem:prod},
\[
-\norm{\nabla f(x_s)}\norm{x_s-x^\ast}
\le
-\sqrt{\frac{\mu}{L_0}}(f(x_s)-f^\ast)
-
\frac{\sqrt{\mu\mu_0}}{2}\norm{x_s-x^\ast}^2,
\]
and inserting this inequality into \eqref{eq:avg_jump_control},
then back into \eqref{eq:avg_drift_2},
we obtain
\[
\dotprod{\nabla \varphi(\overline{x}_s),b(\overline{x}_s)}
+
\mathbb{E}_{\xi}\!\bigl(
\varphi(\overline{x}_s+G(\overline{x}_s,\xi))-\varphi(\overline{x}_s)
\bigr)
\le 0.
\]

Plugging this bound into \eqref{eq:avg_prop_calc} yields
\[
\varphi(\overline{x}_t)\le \varphi(\overline{x}_0)+M_t,
\]
where $M_t$ is a martingale.
Evaluating at $t=T_k$, taking expectations, and using
Theorem~\ref{thm:martingal_stopping}, we get
\[
\E{e^{\theta T_k}(f(\tilde x_k)-f^\ast)}\le \varphi(\overline{x}_0).
\]

Finally,
\[
\varphi(\overline{x}_0)
=
f(\tilde x_0)-f^\ast+\frac{\sqrt{\mu_0L_0}}{2}\norm{\tilde z_0-x^\ast}^2
\le
\Bigl(1+\sqrt{\frac{L_0}{\mu_0}}\Bigr)(f(\tilde x_0)-f^\ast),
\]
and Lemma~\ref{lem:conv_high_proba} concludes the proof.
\end{proof}











\subsection{Possible Extensions: Adapt to other Conditions and Exact Rates}\label{app:exact}
We could consider other extensions of the proposed viewpoint. First, we could consider other averaged condition, such as an \textit{averaged PL} condition:
\[ \int_0^t e^{\theta t}\norm{\nabla f(x_s)}^2ds \ge 2\mu_{avg}\int_0^t (f(x_s)-f^\ast )ds, \]
for some $\mu_{avg} > 0$.
This inequality is automatically verified if $f \in \pl$, with $\mu_{avg} \ge \mu$. Interestingly, $\xcont$ could cross a saddle point $\hat x$, such that $f \notin \pl$, while having that this inequality is verified. In particular, \eqref{eq:nmo} could cross and escape such saddle point, and then converging to a global minima. This saddle point could be "averaged out".

Alternatively, we also could consider point-wise parameter, as illustrated next.
\begin{proposition}\label{prop:exact_gf}
    Define the point-wise Polyak-Lojasiewicz constant $\mu(x) :=  2\frac{\norm{\nabla f(x)}^2}{f(x)-f^\ast}$. If $\xcont \sim \eqref{eq:gf}$, we have
    $$f(x_t)-f^\ast = e^{-2\int_{0}^t \mu(x_s)ds}(f(x_0)-f^\ast).$$
\end{proposition}
We obtain an exact rate for the decrease of $f(x_t)-f^\ast$. If $f\in \pl$, we have in particular $2\int_0^t \mu(x_s)ds \ge 2\mu t$. We can carry similar analysis with \eqref{eq:nmo}, under a point-wise aiming condition.

\begin{proposition}\label{prop:exact_nmo}
    Define the point-wise: Polyak-Lojasiewicz constant $\mu(x) :=  2\frac{\norm{\nabla f(x)}^2}{f(x)-f^\ast}$, aiming condition constant $a(x) := \frac{\dotprod{\nabla f(x),x-x^\ast}}{\norm{\nabla f(x)}\norm{x-x^\ast}}$ and quadratic growth constant $\mu_0(x) = 0.5\cdot\frac{f(x)-f^\ast}{\norm{x-x^\ast}^2}$.
   If $\xcont \sim \eqref{eq:nmo}$ with $ \eta'_t = a(x_t)\sqrt{\mu(x_t)},\quad \eta_t = \frac{L_0}{\sqrt{\mu_0(x_t)}}, \quad \gamma'_t = \frac{1}{\sqrt{\mu_0(x_t)}},\quad \gamma_t \ge 0$, we have
    $$f(x_t)-f^\ast \le e^{-\int_0^t a(x_s)\sqrt{\mu(x_s)}ds}L_0\norm{x_0-x^\ast}^2,$$
    with $L_0 = \sup_{x\in \R^d} \mu_0(x)$.
\end{proposition}
This rate is however purely conceptual, as it involves exact values of parameters we cannot compute. The averaged conditions such as Assumption~\ref{ass:avg_corelation} are more practical, as it is sufficient to use a small enough value of $a_{avg}$, and not an exact one.
\subsubsection{Proof of Proposition~\ref{prop:exact_gf}}
\textbf{(i) Case of gradient flow \eqref{eq:gf}. \:}
    Set
    $$E_t:= e^{2\int_{0}^t \mu(x_s)ds}(f(x_t)-f^\ast). $$
    We have
    \begin{equation}
        \begin{aligned}\label{eq:pl_avg:1}
                    \dot{E_t} &= 2\mu(x_t)e^{2\int_{0}^t \mu(x_s)ds}(f(x_t)-f^\ast) + e^{2\int_{0}^t \mu(x_s)ds}\dotprod{\nabla f(x_t),\dot x_t}\\
        &\overset{\eqref{eq:gf}}{=} 2\mu(x_t)e^{2\int_{0}^t \mu(x_s)ds}(f(x_t)-f^\ast) - e^{2\int_{0}^t \mu(x_s)ds}\norm{\nabla f(x_t)}^2
        \end{aligned}
    \end{equation}

    Recall $\mu(x) := 2\frac{\norm{\nabla f(x)}^2}{f(x)-f^\ast}$. Plugging it in \eqref{eq:pl_avg:1}, we deduce
    \begin{align}
        \dot{E_t} = 2\mu(x_t)e^{2\int_{0}^t \mu(x_s)ds}(f(x_t)-f^\ast) - 2\mu(x_t)e^{2\int_{0}^t \mu(x_s)ds}(f(x_t)-f^\ast)=0.
    \end{align}
    By integrating we deduce 
    \begin{equation}
        f(x_t)-f^\ast = e^{-2\int_{0}^t \mu(x_s)ds}(f(x_0)-f^\ast).
    \end{equation}


\subsubsection{Proof of Proposition~\ref{prop:exact_nmo}}

Let
\begin{equation}
    E_t
    =
    e^{\int_0^t a(x_s)\sqrt{\mu(x_s)}ds}
    \left(f(x_t)-f^\ast + \frac{\beta}{2}\norm{z_t-x^\ast}^2\right).
\end{equation}
Denote
\[
\Delta_t := e^{\int_0^t a(x_s)\sqrt{\mu(x_s)}ds},
\qquad
u_t := a(x_t)\sqrt{\mu(x_t)}.
\]

We have
\begin{align}
    \Delta_t^{-1}\dot E_t
    &= u_t(f(x_t)-f^\ast)
    + \frac{u_t\beta}{2}\norm{z_t-x^\ast}^2
    + \dotprod{\nabla f(x_t),\dot x_t}
    + \beta\dotprod{z_t-x^\ast,\dot z_t} \nonumber\\
    &= u_t(f(x_t)-f^\ast)
    + \frac{u_t\beta}{2}\norm{z_t-x^\ast}^2 \nonumber\\
    &\quad
    + \dotprod{\nabla f(x_t),\eta_t(z_t-x_t)-\gamma_t\nabla f(x_t)} \nonumber\\
    &\quad
    + \beta\dotprod{z_t-x^\ast,\eta'_t(x_t-z_t)-\gamma'_t\nabla f(x_t)} \nonumber\\
    &= u_t(f(x_t)-f^\ast)
    + \frac{u_t\beta}{2}\norm{z_t-x^\ast}^2 \nonumber\\
    &\quad
    + \eta_t\dotprod{\nabla f(x_t),x^\ast-x_t}
    + (\eta_t-\beta\gamma'_t)\dotprod{\nabla f(x_t),z_t-x^\ast} \nonumber\\
    &\quad
    + \beta\eta'_t\dotprod{z_t-x^\ast,x_t-z_t}
    - \gamma_t\norm{\nabla f(x_t)}^2 .
\end{align}

We choose $\eta'_t = u_t$ to cancel the $\norm{z_t-x^\ast}^2$ term, and
$\eta_t = \beta\gamma'_t$ to cancel the mixed scalar product. Then
\begin{equation}\label{eq:exact:nest_ode_1_corr}
    \Delta_t^{-1}\dot E_t
    =
    u_t(f(x_t)-f^\ast)
    + \beta\gamma'_t\dotprod{\nabla f(x_t),x^\ast-x_t}
    + \frac{\beta u_t}{2}\norm{x_t-x^\ast}^2
    - \gamma_t\norm{\nabla f(x_t)}^2 .
\end{equation}

By Lemma~\ref{lem:tec:exact_scalar_prod_bound},
\[
\dotprod{\nabla f(x_t),x^\ast-x_t}
=
- a(x_t)\sqrt{\frac{\mu(x_t)}{\mu_0(x_t)}}(f(x_t)-f^\ast)
- a(x_t)\sqrt{\mu(x_t)\mu_0(x_t)}\frac{1}{2}\norm{x_t-x^\ast}^2 .
\]

Fixing $\gamma'_t = \frac{1}{\sqrt{\mu_0(x_t)}}$ and ignoring the nonpositive term
$-\gamma_t\norm{\nabla f(x_t)}^2$, we obtain
\begin{equation}
    \Delta_t^{-1}\dot E_t
    =
    u_t\Bigl(1-\frac{\beta}{\mu_0(x_t)}\Bigr)(f(x_t)-f^\ast).
\end{equation}

Choosing $\beta = L_0 = \sup_{x\in\R^d}\mu_0(x)$ ensures $\dot E_t \le 0$ for all $t\ge0$.
Therefore $E_t \le E_0$, and
\[
f(x_t)-f^\ast
\le
e^{-\int_0^t a(x_s)\sqrt{\mu(x_s)}ds}\,E_0.
\]

Finally, since $\mu_0(x_0)\le L_0$, we have
\[
f(x_0)-f^\ast = 2\mu_0(x_0)\norm{x_0-x^\ast}^2 \le 2L_0\norm{x_0-x^\ast}^2,
\]
which yields
\[
f(x_t)-f^\ast
\le
e^{-\int_0^t a(x_s)\sqrt{\mu(x_s)}ds}\,L_0\norm{x_0-x^\ast}^2.
\]







\section{Continuized Toolbox}\label{app:continuized_intro}
The continuized framework \cite{even2021continuized} offers an alternative momentum method compared to the more classic ones. In practice, it results in a stochastic parameterization of \eqref{eq:nm}. It can lead to interesting results in some non-convex setting, see \cite{wang2023continuized,hermant2025continuized}. We refer to \cite[Section 3]{hermant2025continuized} for an methodological introduction of this concept, on the class of function $\pl \cap \ls$. In this section, we will simply state the fundamental objects and tools. The fundamental object are the continuized Nesterov equations:
\begin{align}\label{eq:nest_continuized}\tag{CNE}\left\{
    \begin{array}{ll}
        dx_t &= \eta(z_t-x_t)dt - \gamma \nabla f(x_t) dN_t \\
        dz_t &= \eta'(x_t-z_t)dt - \gamma' \nabla f(x_t) dN_t
    \end{array}
\right.
\end{align} where $dN(t) = \sum_{k\ge 0} \delta_{T_k}(dt)$ is a Poisson point measure with intensity $dt$, with $T_1,T_2,\dots$ random times such that for all $i$, $T_{i+1}-T_i \overset{i.i.d}{\sim} \mathcal{E}(1)$. It mixes a continuous component, the $dt$ factor, with a discrete component that acts at random times, the $dN_t$ factor. $\xz$ cannot be differentiated, but we can still use Lyapunov-approaches through Itô calculus. 
\begin{proposition}[\cite{even2021continuized}, Proposition 2]\label{prop:sto_calc}
Let $x_t \in \R^d$ be a solution of 
\begin{equation*}
    dx_t = b(x_t)dt + \int_\Xi G(x_t,\xi)dN(t,\xi)
\end{equation*}
and $\varphi : \R^d \to \R$ be a smooth function. Then,
\begin{equation}
    \varphi(x_t) = \varphi(x_0) + \int_{0}^t \dotprod{\nabla \varphi(x_s),b(x_s)}ds + \int_{[0,t]}  \mathbb{E}_\xi{ \varphi(x_s + G(x_s,\xi)) - \varphi(x_s)}ds + M_t,
\end{equation}

where $M_t$ is a martingale such that $\E{M_t}=0$, $\forall t\ge 0$.

\end{proposition}
Importantly, the following evaluations $\tilde{y}_k := x_{T_{k+1}^-}$, $\tilde{x}_k := x_{T_{k}}$ and $\tilde{z}_k := z_{T_{k}}$ of \eqref{eq:nest_continuized} can be written as an algorithmic form. 

\begin{proposition}[\cite{hermant2025continuized}, Proposition 5]\label{prop:discretization_constant_param}
Let $\xz$ follow \eqref{eq:nest_continuized}. Define $\tilde{y}_k := x_{T_{k+1}^-}$, $\tilde{x}_k := x_{T_{k}}$ and $\tilde{z}_k := z_{T_{k}}$ as evaluations of this process. Then, $(\tilde{y}_k,\tilde{x}_k, \tilde{z}_k)$ writes as a Nesterov algorithm in the form of \eqref{eq:nm} with stochastic parameters 
\begin{align}\label{alg:constant_param}\left\{
    \begin{array}{ll}
        \tilde{y}_k &= (\tilde{z}_{k} - \tilde{x}_k)\frac{\eta}{\eta+\eta'}\bpar{1 -e^{-(\eta + \eta')(T_{k+1}-T_k)} } + \tilde{x}_k \\
        \tilde{x}_{k+1} &=  \tilde{y}_k  - \gamma \nabla f(\tilde{y}_k,\xi_k)\\
        \tilde{z}_{k+1} &=   \tilde{z}_{k} + \eta' \frac{\bpar{1- e^{-(\eta+\eta')(T_{k+1}-T_k)}}}{\eta' + \eta e^{-(\eta+\eta')(T_{k+1}-T_k)}}(\tilde{y}_k - \tilde{z}_{k}) - \gamma' \nabla f(\tilde{y}_k,\xi_k)
    \end{array}
\right.
\end{align}
\end{proposition}
In order to transfer a result involving $\xz$ to $(\tilde x_k,\tilde z_k)_{k\in\N}$, we use need a stopping theorem, such as the following.
\begin{theorem}[Stopping theorem, \cite{hermant2025continuized}, Theorem 6]\label{thm:martingal_stopping}
    Let $(\varphi_t)_{t\in \R_+}$ be a nonnegative process with cadlag trajectories, such that it verifies
    $$ \varphi_t \le K_0 + M_t, $$
    for some positive random variable $K_0$, some martingale $\mart$ with $M_0 = 0$. Then, for an almost surely finite stopping time $\tau$, one has 
$$ \E{\varphi_\tau} \le \E{K_0}.$$
\end{theorem}
Finally, because the continuized analysis leads to result in expectation, in the case where we obtain linear convergence rates, concentration inequalities such as Chernov inequality allows to deduce a result on the trajectory, with high probability.
\begin{lemma}[\cite{hermant2025continuized}, Proposition 14]\label{lem:conv_high_proba}
     Assume $T_1,\dots,T_k$ are random variables such that $T_{i+1} - T_i$ are  i.i.d  of law $\mathcal{E}(1)$, with convention $T_0 = 0$. If the iterations of (\ref{alg:constant_param}) verify
                     \begin{equation}\label{eq:a.s.:2}       \mathbb{E}\left[e^{\beta T_k}(f(\tilde{x}_k)-f^\ast)\right] \le K_0,
         \end{equation}
     then, with probability  $1-\frac{1}{c_0} - e^{-(c_1-1-\log(c_1))k}$, for some $c_0 > 1$ and $c_1 \in (0,1)$ we have
         \begin{equation*}
    f(\tilde{x}_k)-f^\ast \le c_0K_0e^{-\beta(1-c_1)k}.
    \end{equation*}
\end{lemma}


% \subsection{Link with quasar-convexity}
% The following class of functions relaxes the non-strong convexity.
% \begin{definition}
%     For $\tau \in (0,1]$, we call $\qc$ the functions $f$ that satisfy $\forall x \in \R^d,\: f(x)+\frac{1}{\tau}\dotprod{\nabla f(x),x^\ast-x}\le f^\ast$.
% \end{definition}
% Quasar-convexity arises in applications such as learning dynamical systems
% \cite{hardt2018gradient} and training certain generalized linear models
% \cite{wang2023continuized}. 
% For functions in $\qc$, it also holds that the corelation $\frac{\dotprod{\nabla f(x),x^\ast-x}}{\norm{x-x^\ast}\norm{\nabla f(x)}}$ is non-negative. But the underlying inequality does not directly characterizes this quantity, and actually carries more structure. 
% \begin{proposition}
%     Let $f\in \qc$. If $f$ also satisfies
%     \begin{equation}
%         \norm{\nabla f(x) } \le 2L(f(x)-f^\ast),
%     \end{equation}
%     then $f \in \acf{\frac{\tau}{2L}}$.
% \end{proposition}
% \begin{proof}
%     We have
%     \[\dotprod{\nabla f(x),x^\ast-x} \ge \tau(f(x)-f^\ast) \ge \frac{}{}\]
% % \end{proof}
% \section{stock}
% \begin{align*}
%     \nu(x) &= \frac{\dotprod{\nabla f(x),x-x^\ast}}{\norm{x-x^\ast}^2} = a(x)\frac{\norm{\nabla f(x)}}{\norm{x-x^\ast}}\\
%     =& a(x) \frac{\norm{\nabla f(x)}}{\sqrt{f(x)-f^\ast}}\frac{\sqrt{f(x)-f^\ast}}{\norm{x-x^\ast}} = a(x)\sqrt{\mu(x)L_0(x)}.
% \end{align*}


% \begin{align}\label{eq:snm}\tag{NM}\left\{
%     \begin{array}{ll}
%         \tilde{y}_k &= \tilde{x}_k+ \alpha_k(\tilde{z}_{k} - \tilde{x}_k)  \\
%         \tilde{x}_{k+1} &=  \tilde{y}_k  - \gamma \nabla f(\tilde{y}_k,\xi_{k+1})\\
%         \tilde{z}_{k+1} &=   \tilde{z}_{k} + \beta_k(\tilde{y}_k - \tilde{z}_{k}) - \gamma' \nabla f(\tilde{y}_k,\xi_{k+1})
%     \end{array}
% \right.
% \end{align}
% \section{More insightful comparison: RSI instead of Strong quasar-convexity.}
% \begin{proposition}
%     Define the point-wise Polyak-Lojasiewicz constant $\mu(x) :=  2\frac{\norm{\nabla f(x)}^2}{f(x)-f^\ast}$. If $\xcont \sim \eqref{eq:gf}$, we have
%     $$f(x_t)-f^\ast = e^{-2\int_{0}^t \mu(x_s)ds}(f(x_0)-f^\ast).$$
% \end{proposition}
% If $f \in \pl$, we have $\mu(x) \ge \mu$ for all $x\in\R^d$, such that we deduce from Proposition~\ref{prop:exact} the classical result $f(x_t)-f^\ast \le e^{-2\mu t}(f(x_0)-f^\ast)$. But the rate of Proposition~\ref{prop:exact} describes way more accurately the behavior of \eqref{eq:gf}. It indicates that, if the trajectory does not cross regions with very low PL values or these regions are negligible portions of the trajectory, it will not affect the rate. The rate is less sensitive to pathological regions than the global $\mu$-PL constant.
% \begin{proposition}\label{prop:ode_nest_rsi_traj_depend}
%     Define the point-wise RSI constant $\nu(x) = \frac{\dotprod{\nabla f(x),x-x^\ast}}{\norm{x-x^\ast}^2}$, the point-wise growth constant $L_0(x) :=  2\frac{f(x)-f^\ast}{\norm{x-x^\ast}^2}$, and the upper growth contant $\overline{L_0} := \sup_{x\in \R^d} L_0(x)$. If $\xcont \sim \eqref{eq:nmo}$ with $\eta'_t := \int_0^t \frac{\nu(x_s)}{\sqrt{L_0(x_s)}}ds$, $\eta_t := \overline{L_0}/L_0(x_t)$, $\gamma'_t := 1/L_0(x_t)$ and $\gamma_t \ge 0$, then
%     $$f(x_t)-f^\ast \le e^{-\int_0^t \frac{\nu(x_s)}{\sqrt{L_0(x_s)}}ds}\overline{L_0}\norm{x_0-x^\ast}^2.$$
% \end{proposition}
% Recalling \eqref{eq:mu_rsi_link}, the convergence bound of Proposition~\ref{prop:ode_nest_rsi_traj_depend} becomes
% $$f(x_t)-f^\ast = \bigO \bpar{e^{-\int_0^t a(x_s)\sqrt{\mu(x_s)}ds}}.$$
% Where $\mu(\cdot)$ is the point-wise PL constant, and $a(\cdot)$ is the point-wise correlation between the gradient and the direction pointing towards the minimizer. If this correlation along the trajectory is lower bounded, namely $\inf_{s \in [0,t]} a(x_s) > a > 0$, the bound becomes 
% $$f(x_t)-f^\ast = \bigO \bpar{e^{-a\int_0^t \sqrt{\mu(x_s)}ds}}.$$
% If $a$ is large compared to $\mu$, it means that a similar phenomenon as in the strongly convex case occurs, namely changing $\mu(\cdot)$ into $\sqrt{\mu(\cdot)}$ using inertia. In particular if $f$ is 1-dimensional, $a(x) =1$ for any $x \in \R^d$.

% \begin{assumption}\label{ass:l_smooth}
%     $f$ is such that $\forall x,y\in \R^d,
%  ~ f(x)-f(y)-\dotprod{\nabla f(y),x-y}  \le \frac{L}{2}\norm{x-y}^2$, for some $L>0$.
% \end{assumption}
% \begin{assumption}\label{ass:qg+}
%     $f$ is such that $\forall x\in \R^d,
%  ~ f(x)-f^\ast \le \frac{L_0}{2}\norm{x-x^\ast}^2$, for some $L_0>0$.
% \end{assumption}
% \begin{assumption}[Restricted Secant Inequality]\label{ass:rsi}


%    $\forall x \in \mathbb{R}^d, \quad   \dotprod{\nabla f(x),x-x^\ast} \ge \nu \lVert x- x^\ast \rVert^2.$
% \end{assumption}
% % \begin{assumption}[UAAC]\label{ass:uaac}
% %     \begin{equation}\label{eq:uaac}\tag{UAAC}
% %              \forall x \in \mathbb{R}^d, \quad 1 \geqslant \frac{\langle \nabla F(x),x - x^\ast \rangle}{\lVert  \nabla F(x) \rVert \lVert x - x^\ast \rVert} \geqslant a > 0
% %          \end{equation}
% % \end{assumption}
% RSI is weaker than strong quasar-convexity. To make it "equivalent", we need a supplementary assumption.

% QG+ ensures a quadratic upper bound over the function, parameterized by $L_0$. It can be thought as a \textit{global curvature upper bound}, that is, $f \in \qg$ with small $L_0$ is a flat function. In particular, a $L$-smooth function is $L$-QG+. However, as noted in \cite{guille2021study}, a $L$-smooth function can be $L_0$-QG+ with $L_0\llL$. This is particularly true for non-convex functions.
% Functions that are RSI and QG+ are SQC for some parameters \cite[Lemma 4]{hermant2024study}.
% \begin{proposition}\label{prop:gf}
%     Let $\xcont \sim \eqref{eq:gf}$. If $f \in \pl$
% \begin{equation}
%     E_t \le e^{-2\mu t} E_0.
% \end{equation}
% If $f \in \rsi$
% \begin{equation}
%     E_t \le e^{-\nu t}E_0.
% \end{equation}
% \end{proposition}
% We could believe that, as RSI is a stronger assumption than PL, the rate deduced for this assumption will be better. But this is misleading.




% In order to get intuition, consider the following computation. \paragraph{Link RSI-PL}
% Denote $\nu(x) = \frac{\dotprod{\nabla f(x),x-x^\ast}}{\norm{x-x^\ast}^2}$ the RSI constant at $x \in \R^d$. With straightforward computation, we remark
% \begin{align}\label{eq:mu_rsi_link}
%     \nu(x) = a(x)\sqrt{L_0(x)}\sqrt{\mu(x)}.
% \end{align}
% where  $\mu(x) := 2\frac{\norm{\nabla f(x)}^2}{f(x)-f^\ast}$ is the PL constant at $x$, and $L_0(x) := 2\frac{f(x)-f^\ast}{\norm{x-x^\ast}^2} $ is the QG+ constant at $x$, and $$a(x) :=\frac{\langle \nabla f(x),x - x^\ast \rangle}{\lVert  \nabla f(x) \rVert \lVert x - x^\ast \rVert}.$$
% Equation \eqref{eq:mu_rsi_link} gives interesting insights about the link between PL and RSI functions, as we detail in the subsequent paragraphs.
% \jh{parler de la correlation avant}
% \begin{itemize}
%     \item The value $a(x)$ denotes the correlation between the direction given by the gradient of $f$ at $x$, and the direction $x-x^\ast$, that is the direction that, from $x$, points toward the minimizer $x^\star$. If $f$ is 1 dimensional, $a(x)=1$ for all $x\in \R$. This statement is no longer valid as long as $d>1$. In particular, it indicates that if at some point the correlation is very low, the RSI constant can be very low compared to $\mu(x)$. This idea is conveyed by Figure~\ref{fig:low_corr}. We build a function such that, it is RSI on the displayed domain, but due to vanishing correlation $\dotprod{\nabla f(x),x^\ast-x}$, the RSI constant is approximately 100 times lower than the PL constant. For such an example, although the function belongs to $\rsi$ on the displayed domain, caracterizing its convergence using this assumption is bad.
%     \item  The other factor that appears is $\sqrt{L_0(x)}$. It will be high if the function value grows fast. In particular, if the function grows fast, and then the gradient suddenly vanishes close to zero, making $\mu(x)$ vanishes to zero while $L_0(x)$ is fewly affected, then $\nu(x)$ can be significantly higher than $\mu(x)$. This idea is conveyed by Figure~\ref{fig:rsi_pl_1d} on a 1d function, such that $a(x) = 1$ and $\nu(x) =\sqrt{L_0(x)}\sqrt{\mu(x)}$. The function grows quadratically around its minimizer $0$, then suddenly drops in a plateau. In such a plateau, as the gradient vanishes, $\mu(x)$ does so, while $L_0(x)$ is way less sensitive. This result in the RSI constant being $10^{10}$ times higher than the PL constant.
% \end{itemize}


% \jh{prendre RSI pour comparer. Comparer vitesse CV dans figure faite en graphe avec sin x}
% \begin{theorem}\label{thm:sqc}
%     Let $f \in \sqcf{\tau}{\mu}$. Then:
%     \begin{enumerate}
%         \item If $\xcont \sim \eqref{eq:gf}$, $f(x_t)-f^\ast = e^{-\tau \mu t}$.
%         \item If $\xcont \sim \eqref{eq:nmo}$, $f(x_t)-f^\ast = e^{-\tau \sqrt{\mu} t}$.
%     \end{enumerate}
% \end{theorem}
% In the case of Theorem~\ref{thm:sqc}, contrarily to Theorem~\ref{thm:pl}, it makes no sense to take the sup of two variables.
% Moreover,
% $\arg \sup_\tau \arg \sup_\mu \{ f \in \sqc \}$ and $\arg \sup_\mu \arg \sup_\tau \{ f \in \sqc \}$ would lead to different choices of parameters $(\tau,\mu)$. In particular, we have that $f \in \sqcf{\tau}{\mu} \Rightarrow f \in \sqcf{\theta \tau}{\mu/\theta}$, for all $\theta \in (0,1]$ \cite[Observation 5]{hinder2020near}. This means that 
% $\arg \sup_\tau \arg \sup_\mu \{ f \in \sqc \} = (0,+\infty)$. There is no canonical way to chose $(\tau,\mu)$ such that the bounds are optimal. Worse, the optimal choices depends on the bound we want to optimize, as illustrated in Figure~\ref{fig:sqc_flaws}. We observe that the optimal choice $(\tau,\mu)$ differs from gradient descent/flow and Nesterov method. In that, stating that "we compare $(1-\tau\frac{\mu}{L})$ and $(1-\tau\sqrt{\frac{\mu}{L}})$" is misleading, and it is not clear what it states. \jh{$\frac{\mu}{L}$ analogy with strongly convex is not very good}

% We have the following result.
% \begin{theorem}\label{thm:pl}
%     Let $f$ verify Assumption~\ref{ass:pl}. Then:
%     \begin{enumerate}
%         \item If $\xcont \sim \eqref{eq:gf}$, $f(x_k)-f^\ast = e^{-\mu t}$
%         \item If $\xdisc \sim \eqref{eq:gd}$, $f(x_k)-f^\ast = (1-\frac{\mu}{L})^k.$
%     \end{enumerate}
% \end{theorem}






% \section{Algorithmic result}


% \begin{proposition}\label{prop:gd_rsi}
% Let $\xdisc$ be solution of gradient descent. Under \jh{ref sgc} Assumptions~\ref{ass:eb},\ref{ass:upp_qg} and \ref{ass:l-smooth}, we have  \begin{equation}
%  f(\tilde x_k)-f^\ast \le \bpar{1-a\frac{\nu}{ L}}(f(x_0)-f^\ast)
% \end{equation}
% \end{proposition}

% \begin{proposition}\label{prop:nest_rsi}
% Let Assumptions~\ref{ass:eb}, \ref{ass:l-smooth}, \ref{ass:uaac}, \ref{ass:upp_qg} and \ref{ass:sgc} holds true.
% Assume $\xdisc$ solution of 
% \begin{align}\left\{
%     \begin{array}{ll}
%         \tilde{y}_k &= (\tilde{z}_{k} - \tilde{x}_k)\frac{\eta}{\eta+\eta'}\bpar{1 -e^{-(\eta + \eta')(T_{k+1}-T_k)} } + \tilde{x}_k \\
%         \tilde{x}_{k+1} &=  \tilde{y}_k  - \gamma \nabla f(\tilde{y}_k,\xi_{k+1})\\
%         \tilde{z}_{k+1} &=   \tilde{z}_{k} + \eta' \frac{\bpar{1- e^{-(\eta+\eta')(T_{k+1}-T_k)}}}{\eta' + \eta e^{-(\eta+\eta')(T_{k+1}-T_k)}}(\tilde{y}_k - \tilde{z}_{k}) - \gamma' \nabla f(\tilde{y}_k,\xi_{k+1})
%     \end{array}
% \right.
% \end{align}
% with  $  \eta' = \frac{a}{\rho}\frac{\nu}{\sqrt{L_0}}\frac{1}{\sqrt{L}},\quad \gamma_s = \frac{1}{\rho L}\quad \gamma' = \frac{1}{\rho} \frac{1}{\sqrt{L_0 L}}, $ and $\eta =\frac{1}{\rho}\sqrt{\frac{L_0}{L}}$. Then, we have
% \begin{equation}
% \E{e^{ \frac{a}{\rho}\frac{\nu}{\sqrt{L_0}}\frac{1}{\sqrt{L}}T_k}f(\tilde x_k)-f^\ast)} \le f(x_0)-f^\ast + \frac{L_0}{2}\norm{x_0-x^\ast}^2.
% \end{equation}
% \end{proposition}

% \jh{$\frac{L_0}{L}$ as a non-convexity measurement?}
% \section{Intuition on dynamical system/Continuous time perspective}
% \begin{proposition}\label{prop:gf}
%     Let Assumptions~\ref{ass:pl} hold, and let $f$ verify
% \begin{equation}
%     \dot{x}_t = -\nabla f(x_t).
% \end{equation}
% Then there is a ball such that the function lies in the growther
% \begin{equation}
%     E_t \le e^{-2\mu t} E_0.
% \end{equation}
% \end{proposition}

% We can also deduce a result using Assumption~\ref{ass:uaac}.
% \begin{proposition}\label{prop:gf_angle}
%     Let Assumptions~\ref{ass:rsi}, \ref{ass:upp_qg} hold, and let $f$ verify
% \begin{equation}
%     \dot{x}_t = -\nabla f(x_t).
% \end{equation}
% Then
% \begin{equation}
%     E_t \le e^{-\nu t}E_0.
% \end{equation}
% \end{proposition}

% Let $f \in \rsi$ for some $\nu > 0$. Because $f$ is $L$-smooth, we have $f \in \plf{\nu^2/L_0}$. Because under $L$-smoothness, a function verifying RSI implies it verifies PL, meaning it is stronger, we could be tempted to say the bound of Proposition~\ref{prop:gf_angle} under a stronger assumption will be better.
% However, the inclusion means that there exists $\mu \ge \nu^2/L_0$, such that $f \in \rsi \cap \pl$. It raises the question: is $\nu$ the constant of RSI, better than $\mu$ of PL ? If not, in which case ? 








% \begin{proposition}\label{prop:nest_ode_triple}  Let Assumptions~\ref{ass:rsi}, \ref{ass:upp_qg} hold, and let $f$ verify
% \begin{align}\label{eq:nmo_appendix}\left\{
%     \begin{array}{ll}
%         \dot{x}_t &= \eta(z_t-x_t) - \gamma \nabla f(x_t) \\
%         \dot{z}_t &= \eta'(x_t-z_t) - \gamma' \nabla f(x_t) 
%     \end{array}
% \right.
% \end{align}
% with  $\eta' = \frac{\nu}{\sqrt{L_0}},\quad \gamma' = \frac{1}{\sqrt{L_0}}, \quad \eta = \sqrt{L_0}$ and $\gamma > 0$.
% Then, we have
%     \begin{equation}
%           f(x_t)-f^\ast \le e^{-\frac{\nu}{\sqrt{L_0}}t}E_0.
%     \end{equation}
% \end{proposition}

% \jh{il peut y avoir aussi un aspect lancer les algo/edo sur les fonctions. Dans le cas 1d, on pourrait peut etre voir l'inertiel bloquer dans le plateau à la remontée, ce qui expliquerait le problème du $L_0$}

% \section{Stock figures}
% \begin{figure}
%     \centering
%     \includegraphics[width=0.1\linewidth]{figures/f_1d_graph.pdf}
%     \includegraphics[width=0.1\linewidth]{figures/f_1d_rsi_vs_pl.pdf}
%     \caption{Left: graph of $f(t) = 5(t+0.19\sin(5t))^2$. Right: plot of $\ln(\mu(t)$ and $\log(\nu(t))$. The computed minimum value of $\mu(t)$ is of the order $10^{-23}$, while the computed minimum value of $\nu(t)$ is of the order $10^{-13}$. On this example, the RSI constant is approximately $10^{10}$ higher.}
%     \label{fig:rsi_pl_1d}
% \end{figure}
% \begin{figure}
%     \centering
%     \includegraphics[width=0.1\linewidth]{figures/surface_fx_y.png}
%     \includegraphics[width=0.1\linewidth]{figures/angle.png}
%     \caption{Left: graph of $f(x,y) = 0.5(0.5x^2 -y)^2 + 0.05x^2$. It has $(0,0)$ has unique global minimizer. Right: colormap of $\frac{\mu(x)}{\nu(x)}$, in logscale. Because of the bad correlation, $\nu(x)$ vanishes in the top edges. Globally, $\mu \approx 0.02$, $\nu \approx 10^{-4}$. }
%     \label{fig:low_corr}
% \end{figure}

% \section{Proof with RSI}
% \subsection{Proof of Proposition~\ref{prop:gf_angle}}

%     Let
% \begin{equation}
%     E_t = f(x_t)-f^\ast
%      + \frac{L_0}{2}\norm{x_t-x^\ast}^2.
% \end{equation}
% We have
% \begin{align}
%     \dot E_t = \dotprod{\nabla f(x_t),\dot x_t} &=  -\norm{\nabla f(x_t)}^2 - L_0 \dotprod{\nabla f(x_t),x_t-x^\ast} \\
%     &\le -\norm{\nabla f(x_t)}^2    - L_0  \nu \norm{x_t-x^\ast}^2\\
%     & \le  - L_0  \nu \norm{x_t-x^\ast}^2 \\
%     & = - \frac{L_0 \nu }{2}\norm{x_t-x^\ast}^2 - \frac{L_0  \nu }{2}\norm{x_t-x^\ast}^2   \\
%     & \le - \nu (f(x_t)-f^\ast) - \frac{L_0  \nu }{2}\norm{x_t-x^\ast}^2  \\
%     &\le -\nu E_t.
% \end{align} 
% So, we have
% \begin{equation}
%     E_t \le e^{-\nu t}E_0.
% \end{equation}

% \subsection{Proof of Proposition~\ref{prop:nest_ode_triple}}
%     Let be
% \begin{equation}
%     E_t = e^{\theta t}(f(x_t)-f^\ast) +\beta \frac{ e^{\theta t}}{2}\norm{z_t -  x^\ast}^2
% \end{equation}
% We have
% \begin{equation}
%     E_t = E_0 + \int_0^t \dot{E}_sds,
% \end{equation}
% where
% \begin{align}
%     \dot{E}_t &= \theta  e^{\theta t}(f(x_t)-f^\ast)  +\theta \beta \frac{ e^{\theta t}}{2}\norm{z_t -  x^\ast}^2\nonumber\\
%     &+ e^{\theta t}\dotprod{\nabla f(x_t),\dot{x}_t} + \beta  e^{\theta t}\dotprod{z_t -  x^\ast,\dot{z}_t}
%    \nonumber \\&= \theta  e^{\theta t}(f(x_t)-f^\ast)  +\theta \beta \frac{ e^{\theta t}}{2}\norm{z_t -  x^\ast}^2\nonumber\\
%     &+ e^{\theta t}\dotprod{\nabla f(x_t), \eta_t(z_t-x_t) - \gamma_t \nabla f(x_t)} +\beta   e^{\theta t}\dotprod{z_t - x^\ast,\eta'_t(x_t-z_t) - \gamma'_t \nabla f(x_t) }\nonumber\\
%     &=\theta  e^{\theta t}(f(x_t)-f^\ast)  +\theta \beta \frac{ e^{\theta t}}{2}\norm{z_t -  x^\ast}^2\nonumber\\
%     &+ e^{\theta t}\eta \dotprod{\nabla f(x_t), x^\ast - x_t} + e^{\theta t}(\eta -\beta \gamma')\dotprod{\nabla f(x_t),z_t - x^\ast}\nonumber\\
%     &+ \beta  \eta' e^{\theta t}\dotprod{z_t - x^\ast,x_t - z_t} - \gamma e^{\theta t}\norm{\nabla f(x_t)}^2\nonumber\\
%     &= \theta  e^{\theta t}(f(x_t)-f^\ast)  +\frac{ e^{\theta t}\beta }{2} \bpar{\theta  - \eta'}\norm{z_t -  x^\ast}^2\nonumber\\
%     &+ e^{\theta t}\eta \dotprod{\nabla f(x_t), x^\ast - x_t} + e^{\theta t}(\eta -\beta \gamma')\dotprod{\nabla f(x_t),z_t - x^\ast}\nonumber\\
%     &+ \frac{\beta  \eta' e^{\theta t}}{2}\norm{x^\ast-x_t}^2 - \gamma e^{\theta t}\norm{\nabla f(x_t)}^2.\label{eq:nest_ode_0}
% \end{align}
% We chose $\eta' = \theta$ to cancel the $\norm{z_t-x^\ast}^2$ term. We choose $\eta = \beta\gamma'$ to cancel the $\dotprod{\nabla f(x_t),z_t - x^\ast}$ term. Then,
% \begin{align}
%     e^{-\theta t}\dot{E}_t = \theta(f(x_t)-f^\ast) +\beta \gamma' \dotprod{\nabla f(x_t), x^\ast - x_t}+  \frac{\beta  \theta }{2}\norm{x^\ast-x_t}^2 - \gamma \norm{\nabla f(x_t)}^2.
% \end{align}
% Combining Assumptions~\ref{ass:uaac} and \ref{ass:eb}, we get
% \begin{equation}\label{eq:nest_ode_1}
%     \beta \gamma' \dotprod{\nabla f(x_t), x^\ast - x_t} \le -a\beta \gamma' \norm{\nabla f(x_t)}\norm{x_t-x^\ast} \le -a\nu\beta \gamma'\norm{x_t-x^\ast}^2.
% \end{equation}
% Using Assumption~\ref{ass:upp_qg}, we have
% \begin{equation}\label{eq:nest_ode_2}
%     \theta(f(x_t)-f^\ast) \le \frac{\theta L_0}{2}\norm{x_t-x^\ast}^2.
% \end{equation}
% Combining \eqref{eq:nest_ode_0}, \eqref{eq:nest_ode_1}, \eqref{eq:nest_ode_2}, and ignoring the non positive $\norm{\nabla f(x_t)}^2$ term, we obtain
% \begin{equation}
%     e^{-\theta t}\dot E_t \le  \bpar{\frac{\theta L_0}{2} + \frac{\beta \theta}{2} - a\nu\beta \gamma'}\norm{x_t-x^\ast}.
% \end{equation}
% If we add the constraint $\gamma' \eta = 1 \Rightarrow \gamma' = \frac{1}{\sqrt{L_0}}$, we then need to cancel
% \begin{align}
%     \bpar{\frac{\theta L_0}{2} + \frac{\beta \theta}{2} - a\nu\sqrt{\beta}} = 0 \Leftrightarrow \theta = 2\frac{a\nu \sqrt{\beta}}{L_0 + \beta}.
% \end{align}
% Basic derivative computations leads that $2\frac{a\nu \sqrt{\beta}}{L_0 + \beta}$ is maximized for $\beta > 0$ at $\beta= L_0$. So, we get that if
% \begin{equation}
%     \theta = \eta' = a\frac{\nu}{\sqrt{L_0}},\quad \gamma' = \frac{1}{\sqrt{L_0}}, \quad \eta = \beta \gamma' = \sqrt{L_0},
% \end{equation}
%  we have 
%  \begin{align}
%      E_t \le E_0,
%  \end{align}
%  or 
%  \begin{equation}
%      f(x_t)-f^\ast \le e^{-a\frac{\nu}{\sqrt{L_0}}t}E_0.
%  \end{equation}



%      Let be
% \begin{equation}
%     E_t = e^{\theta t}(f(x_t)-f^\ast) +\beta \frac{ e^{\theta t}}{2}\norm{z_t -  x^\ast}^2
% \end{equation}
% We have
% \begin{equation}
%     E_t = E_0 + \int_0^t \dot{E}_sds,
% \end{equation}
% where
% \begin{align}
%     \dot{E}_t &= \theta  e^{\theta t}(f(x_t)-f^\ast)  +\theta \beta \frac{ e^{\theta t}}{2}\norm{z_t -  x^\ast}^2\nonumber\\
%     &+ e^{\theta t}\dotprod{\nabla f(x_t),\dot{x}_t} + \beta  e^{\theta t}\dotprod{z_t -  x^\ast,\dot{z}_t}
%    \nonumber \\&= \theta  e^{\theta t}(f(x_t)-f^\ast)  +\theta \beta \frac{ e^{\theta t}}{2}\norm{z_t -  x^\ast}^2\nonumber\\
%     &+ e^{\theta t}\dotprod{\nabla f(x_t), \eta_t(z_t-x_t) - \gamma_t \nabla f(x_t)} +\beta   e^{\theta t}\dotprod{z_t - x^\ast,\eta'_t(x_t-z_t) - \gamma'_t \nabla f(x_t) }\nonumber\\
%     &=\theta  e^{\theta t}(f(x_t)-f^\ast)  +\theta \beta \frac{ e^{\theta t}}{2}\norm{z_t -  x^\ast}^2\nonumber\\
%     &+ e^{\theta t}\eta \dotprod{\nabla f(x_t), x^\ast - x_t} + e^{\theta t}(\eta -\beta \gamma')\dotprod{\nabla f(x_t),z_t - x^\ast}\nonumber\\
%     &+ \beta  \eta' e^{\theta t}\dotprod{z_t - x^\ast,x_t - z_t} - \gamma e^{\theta t}\norm{\nabla f(x_t)}^2\nonumber\\
%     &= \theta  e^{\theta t}(f(x_t)-f^\ast)  +\frac{ e^{\theta t}\beta }{2} \bpar{\theta  - \eta'}\norm{z_t -  x^\ast}^2\nonumber\\
%     &+ e^{\theta t}\eta \dotprod{\nabla f(x_t), x^\ast - x_t} + e^{\theta t}(\eta -\beta \gamma')\dotprod{\nabla f(x_t),z_t - x^\ast}\nonumber\\
%     &+ \frac{\beta  \eta' e^{\theta t}}{2}\norm{x^\ast-x_t}^2 - \gamma e^{\theta t}\norm{\nabla f(x_t)}^2.\label{eq:nest_ode_0}
% \end{align}
% We chose $\eta' = \theta$ to cancel the $\norm{z_t-x^\ast}^2$ term. We choose $\eta = \beta\gamma'$ to cancel the $\dotprod{\nabla f(x_t),z_t - x^\ast}$ term. Then,
% \begin{align}\label{eq:pl_corr_nest_1}
%     e^{-\theta t}\dot{E}_t = \theta(f(x_t)-f^\ast) +\eta\dotprod{\nabla f(x_t), x^\ast - x_t}+  \frac{\beta  \theta }{2}\norm{x^\ast-x_t}^2 - \gamma \norm{\nabla f(x_t)}^2.
% \end{align}
% We have
% \begin{equation}\label{eq:pl_corr_nest_2}
%    \dotprod{\nabla f(x_t), x^\ast - x_t} = \frac{\dotprod{\nabla f(x_t), x^\ast - x_t}}{ \norm{\nabla f(x_t)}\norm{x_t-x^\ast} } \frac{\norm{\nabla f(x_t)}}{\sqrt{f(x_t)-f^\ast}}\sqrt{f(x_t)-f^\ast} \norm{x_t-x^\ast}\le -a \sqrt{2\mu}\sqrt{\frac{2}{L_0}}(f(x_t)-f^\ast).
% \end{equation}
% In the same time, we have
% \begin{equation}\label{eq:pl_corr_nest_3}
%    \dotprod{\nabla f(x_t), x^\ast - x_t} = \frac{\dotprod{\nabla f(x_t), x^\ast - x_t}}{ \norm{\nabla f(x_t)}\norm{x_t-x^\ast} }\norm{x_t-x^\ast}\norm{\nabla f(x_t)}\le -a \mu\norm{x_t-x^\ast}^2.
% \end{equation}

% So, writing $\dotprod{\nabla f(x_t), x^\ast - x_t}$ in \eqref{eq:pl_corr_nest_1} as $\frac{1}{2}\dotprod{\nabla f(x_t), x^\ast - x_t} + \frac{1}{2}\dotprod{\nabla f(x_t), x^\ast - x_t}$, in which we inject \eqref{eq:pl_corr_nest_2}, \eqref{eq:pl_corr_nest_3}, and ignoring the $- \gamma \norm{\nabla f(x_t)}^2$ term, we deduce
% \begin{align}
%     e^{-\theta t}\dot{E}_t &\le  \theta(f(x_t)-f^\ast) -\beta \gamma' a \sqrt{\frac{\mu}{L_0}}(f(x_t)-f^\ast) - \frac{\eta a \mu}{2}\norm{x_t-x^\ast}^2+ \frac{\beta  \theta }{2}\norm{x^\ast-x_t}^2  \nonumber\\
%     &=\bpar{\theta - \eta a  \sqrt{\frac{\mu}{L_0}}}(f(x_t)-f^\ast) +\frac{1}{2} \bpar{\beta \theta - \eta a \mu}\norm{x_t-x^\ast}^2.\label{eq:pl_corr_nest_4}
% \end{align}
% We fix $\eta = \sqrt{L_0}$, such that $\theta = a\sqrt{\mu}$, which cancels the left term in \eqref{eq:pl_corr_nest_4}. The cancelation of the right term follows from the following elementary computations
% \begin{equation}
%     \beta \theta - \eta a \mu = \beta a \sqrt{\mu} - a \sqrt{L_0} \mu= 0 \Leftrightarrow \beta = \sqrt{\mu L_0}.
% \end{equation}

% From this, we deduce that $\eta' = \frac{1}{\sqrt{\mu}}$ and
% $$f(x_t)-f^\ast \le e^{-a\sqrt{\mu} t}E_0 \le e^{-a\sqrt{\mu}t}\frac{\sqrt{L_0}\sqrt{L_0-\mu}}{2}\norm{x_0-x^\ast}^2.$$

% If we add the constraint $\gamma' \eta = 1 \Rightarrow \gamma' = \frac{1}{\sqrt{L_0}}$, we then need to cancel
% \begin{align}
%     \bpar{\frac{\theta L_0}{2} + \frac{\beta \theta}{2} - a\nu\sqrt{\beta}} = 0 \Leftrightarrow \theta = 2\frac{a\nu \sqrt{\beta}}{L_0 + \beta}.
% \end{align}
% Basic derivative computations leads that $2\frac{a\nu \sqrt{\beta}}{L_0 + \beta}$ is maximized for $\beta > 0$ at $\beta= L_0$. So, we get that if
% \begin{equation}
%     \theta = \eta' = a\frac{\nu}{\sqrt{L_0}},\quad \gamma' = \frac{1}{\sqrt{L_0}}, \quad \eta = \beta \gamma' = \sqrt{L_0},
% \end{equation}
%  we have 
%  \begin{align}
%      E_t \le E_0,
%  \end{align}
%  or 
%  \begin{equation}
%      f(x_t)-f^\ast \le e^{-a\frac{\nu}{\sqrt{L_0}}t}E_0.
%  \end{equation}



  
% \subsection{Proof of Proposition~\ref{prop:exact}}
%     Set
%     $$E_t:= e^{2\int_{0}^t \mu(x_s)ds}(f(x_t)-f^\ast). $$
%     We have
%     \begin{equation}
%         \begin{aligned}\label{eq:pl_avg:1}
%                     \dot{E_t} &= 2\mu(x_t)e^{2\int_{0}^t \mu(x_s)ds}(f(x_t)-f^\ast) + e^{2\int_{0}^t \mu(x_s)ds}\dotprod{\nabla f(x_t),\dot x_t}\\
%         &\overset{\eqref{eq:gf}}{=} 2\mu(x_t)e^{2\int_{0}^t \mu(x_s)ds}(f(x_t)-f^\ast) - e^{2\int_{0}^t \mu(x_s)ds}\norm{\nabla f(x_t)}^2
%         \end{aligned}
%     \end{equation}

%     Recall $\mu(x) := 2\frac{\norm{\nabla f(x)}^2}{f(x)-f^\ast}$. Plugging it in \eqref{eq:pl_avg:1}, we deduce
%     \begin{align}
%         \dot{E_t} = 2\mu(x_t)e^{2\int_{0}^t \mu(x_s)ds}(f(x_t)-f^\ast) - 2\mu(x_t)e^{2\int_{0}^t \mu(x_s)ds}(f(x_t)-f^\ast)=0.
%     \end{align}
%     By integrating we deduce 
%     \begin{equation}
%         f(x_t)-f^\ast = e^{-2\int_{0}^t \mu(x_s)ds}(f(x_0)-f^\ast).
%     \end{equation}
  



 

%        \begin{assumption}[Strong growth condition]\label{ass:sgc}
%     We assume we have access to stochastic approximations $\nabla f(x,\xi)$ of the real gradient $\nabla f(x)$, where $\xi \sim \mathcal{P}(\Xi)$, such that
%     \begin{enumerate}[label=(\roman*)]
%         \item (Unbiased estimator) $\forall x \in \R^d, ~ \E{\nabla f(x,\xi)} = \nabla f(x)$. 
%         \item (Strong Growth Condition) $\exists \rho \ge 1, \forall x \in \R^d,~ \mathbb{E}_{\xi}\left[ \norm{\nabla f(x,\xi)}^2 \right] \le \rho \norm{\nabla f(x)}^2$. 
%     \end{enumerate} 
% \end{assumption}
% \begin{lemma}\label{lem:descent_sgc}
% Under Assumption~\ref{ass:l_smooth} and \ref{ass:sgc}, if $\gamma \le \frac{1}{L \rho}$, we have

% \begin{equation}
%      \mathbb{E}\left[ f(x-\gamma \nabla f(x,\xi)) - f(x) \right] \leq \gamma\left(\frac{L \rho}{2}\gamma - 1 \right) \lVert  \nabla f(x) \rVert^2.
% \end{equation}

% \end{lemma}
% \begin{proof}
%     By Assumption~\ref{ass:l_smooth}
%     \begin{eqnarray}
%         \begin{aligned}\label{eq:descent_lemma_0}
%                    f(x-\gamma \nabla f(x,\xi)) &\le f(x) + \dotprod{\nabla f(x),x-\gamma \nabla f(x,\xi) - x} + \frac{L}{2}\norm{ x-\gamma \nabla f(x,\xi) - x}^2\\
%         &= f(x) -\gamma \dotprod{\nabla f(x),\nabla f(x,\xi)} + \frac{\gamma^2L}{2}\norm{\nabla f(x,\xi)}^2.
%         \end{aligned}
%     \end{eqnarray}
% We take expectation on both sides of \eqref{eq:descent_lemma_0}
% \begin{equation}\label{eq:descent_lemma_1}
%     \E{f(x-\gamma \nabla f(x,\xi)) - f(x)} \le -\gamma \dotprod{\nabla f(x),\E{\nabla f(x,\xi)}} + \frac{\gamma^2L}{2}\E{\norm{\nabla f(x,\xi)}^2}.
% \end{equation}
% By Assumption~\ref{ass:sgc}, we have
% \begin{equation}\label{eq:descent_lemma_2}
%     \dotprod{\nabla f(x),\E{\nabla f(x,\xi)}}= \norm{\nabla f(x)}^2,\quad \E{\norm{\nabla f(x,\xi)}^2} \le \rho \norm{\nabla f(x)}^2.
% \end{equation}
% Injecting \eqref{eq:descent_lemma_2} in \eqref{eq:descent_lemma_1}, we have
% $$\E{f(x-\gamma \nabla f(x,\xi)) - f(x)} \le \gamma\left(\frac{L \rho}{2}\gamma - 1 \right) \lVert  \nabla f(x) \rVert^2.$$
% \end{proof}
% \subsection{Proof of Theorem~\ref{thm:gd_pl_&_ac}}
% \jh{SGC ??}

% As $E_n \ge f(\tilde x_k)-f^\ast$, and $E_0 \le \frac{\sqrt{L_0}(\sqrt{L_0}-\sqrt{\mu_0})}{2}\norm{x_0-x^\ast}^2$, see Lemma~\ref{lem:prod}, we conclude that
% $$ f(\tilde x_k)-f^\ast \le \bpar{ 1-  a\frac{\sqrt{\mu \mu_0 }}{L}}^n\frac{\sqrt{L_0}(\sqrt{L_0}-\sqrt{\mu_0})}{2}\norm{x_0-x^\ast}^2. $$
% \begin{equation}
%     f(\tilde x_k)-f^\ast \le \bpar{1-a\frac{\nu}{ L}}(f(x_0)-f^\ast)
% \end{equation}
% \subsection{Proof of Proposition~\ref{prop:gd_rsi}}
%     Define 
% \begin{equation}
%     E_n = f(\tilde x_k)-f^\ast 
%     + \frac{L_0}{2}\norm{\tilde x_k-x^\ast}^2.
% \end{equation}
% We have
% \begin{align}
%     E_{n+1} - E_n &\le -\frac{1}{2L}\norm{\nabla f(\tilde x_k)}^2 + \frac{L_0}{2L^2}\norm{\nabla f(\tilde x_k)}^2 - \frac{L_0}{L}\dotprod{\tilde x_k-x^\ast,\nabla f(\tilde x_k)} \\
%     &\le -\frac{L_0}{L}a\norm{\tilde x_k-x^\ast}\norm{\nabla f(\tilde x_k)} \\
%     &\le -\frac{L_0}{L}a\nu \norm{\tilde x_k-x^\ast}^2\\
%     &= -\frac{L_0}{2L}a\nu\norm{\tilde x_k-x^\ast}^2 -\frac{L_0}{2L}a\nu \norm{\tilde x_k-x^\ast}^2\\
%     &\le - a\frac{\nu}{L}(f\tilde x_k)-f^\ast) -\frac{L_0}{2L}a\nu \norm{\tilde x_k-x^\ast}^2\\
%     &= -a\frac{\nu}{L}E_n.
% \end{align}
% So, we have
% \begin{equation}
%     f(\tilde x_k)-f^\ast \le \bpar{1-a\frac{\nu}{ L}}(f(x_0)-f^\ast)
% \end{equation}


% \subsection{Proof of Theorem~\ref{thm:pl_accel}}\label{app:proof:pl_accel}
 

% \subsection{Proof of Proposition~\ref{prop:nest_rsi}}

%     Let $\overline{x}_t= (t,x_t,z_t)$. It satisfies $d\overline{x}_t = b(\overline{x}_t)dt + G(\overline{x}_t,\xi)dN(t,\xi)$, where
% \begin{equation}
%     b(\overline{x}_t) = \begin{pmatrix}
%     1 \\ \eta_t(z_t-x_t) \\ \eta_t'(x_t-z_t)
%     \end{pmatrix},\quad G(\overline{x}_t,\xi) = \begin{pmatrix}
%     0 \\ -\gamma_t \nabla f(x_t,\xi) \\ -\gamma_t' \nabla f(x_t,\xi).
%     \end{pmatrix}
% \end{equation}
% We apply Proposition~\ref{prop:sto_calc} to $\phi_t= \varphi(\overline{x}_t)$, where
% \begin{equation}
%     \varphi(t,x,z) =  A_t(f(x)-f^\ast) + \frac{B_t}{2}\norm{z - x^\ast}^2.
% \end{equation}
% so 
% \begin{equation}\label{eq:prop_calc_sto}
%     \varphi(\overline{x}_t) = \varphi(\overline{x}_0) + \int_{0}^t \dotprod{\nabla \varphi(\overline{x}_s),b(\overline{x}_s)}ds + \int_{[0,t]} \mathbb{E}_\xi( \varphi(\overline{x}_s + G(\overline{x}_s,\xi)) - \varphi(\overline{x}_s))ds + M_t,
% \end{equation}
% where $M_t$ is a martingale.
% We have 
% \begin{align}
%     \frac{\partial \varphi}{\partial s} &= \frac{dA_s}{ds}(f(x)-f^\ast + \frac{1}{2}\frac{dB_s}{ds}\norm{z - x^\ast}^2, \\
%     \frac{\partial \varphi}{\partial x} &= A_s\nabla f(x),\\
%     \quad  \frac{\partial \varphi}{\partial z} &= B_s(z - x^\ast).
% \end{align}
% So, 
% \begin{align}
%     \dotprod{\nabla \varphi(\overline{x}_s),b(\overline{x}_s)} &=\frac{dA_s}{ds}(f(x_s)-f^\ast) + \frac{1}{2}\frac{dB_s}{ds}\norm{z_s - x^\ast}^2 \nonumber \\
%     &+\dotprod{A_s \nabla f(x_s),\eta_s(z_s-x_s)} + \dotprod{B_s(z_s-x^\ast),\eta_s'(x_s - z_s)} \nonumber\\
%     &=\frac{dA_s}{ds}(f(x_s)-f^\ast) + \frac{1}{2}\frac{dB_s}{ds}\norm{z_s - x^\ast}^2\label{eq:nm_rsi_eq_1}\\
%     &+A_s \eta_s \dotprod{ \nabla f(x_s) ,z_s-x_s}+ B_s \eta_s' \dotprod{z_s-x^\ast,x_s - z_s}\label{eq:nm_rsi_eq_2}
% \end{align}
% Also, 
% \begin{align}
%    \mathbb{E}_{\xi} \varphi(\overline{x}_s+ G(\overline{x}_s)) - \varphi(\overline{x}_s) &= \mathbb{E}_{\xi}[A_s(f(x_s - \gamma_s \nabla f(x_s,\xi)) - f(x_s))] \nonumber\\
%    &+ \frac{B_s}{2}\mathbb{E}_{\xi}[\norm{ z_s - \gamma'_s\nabla f(x_s,\xi) - x^\ast}^2 - \norm{z_s -x^\ast}^2]\\
%     &= A_s\mathbb{E}_{\xi}[f(x_s - \gamma_s \nabla f(x_s,\xi)) - f(x_s)]  + \frac{B_s \gamma_s^{'2}}{2}\mathbb{E}_{\xi}[\norm{\nabla f(x_s,\xi)}^2] \label{eq:nm_rsi_eq_3}\\
%     &+ B_s\gamma'_s \dotprod{\nabla f(x_s),x^\ast - z_s} \label{eq:nm_rsi_eq_4}
% \end{align}
% where we used $\mathbb{E}_{\xi}[\nabla f(x_s,\xi)] = \nabla f(x_s)$.
% Using Assumptions \ref{ass:l_smooth}-\ref{ass:sgc} + Lemma~\ref{lem:descent_sgc}, we have
% \begin{align}
%     (\ref{eq:nm_rsi_eq_3}) \le \bpar{\rho B_s \gamma_s^{'2} + A_s \gamma_s( \rho\gamma_s L-2)}\frac{1}{2}\norm{\nabla f(x_s)}^2.
% \end{align}
% By Assumption~\ref{ass:eb}, we have
% \begin{equation}
%     \frac{dA_s}{ds}(f(x_s)-f^\ast)  \le \frac{L_0}{2}\frac{dA_s}{ds}\norm{x_s-x^\ast}^2.
% \end{equation}
% Also, we have
% \begin{equation}
%     \dotprod{z_s-x^\ast,x_s - z_s} = \frac{1}{2}\norm{x_s-x^\ast}^2 -\frac{1}{2}\norm{z_s-x^\ast}^2-\frac{1}{2}\norm{x_s-z_s}^2 \le\frac{1}{2}\norm{x_s-x^\ast}^2 -\frac{1}{2}\norm{z_s-x^\ast}^2.  
% \end{equation}
% Then,
% \begin{align}
%     A_s \eta_s \dotprod{ \nabla f(x_s) ,z_s-x_s} +  B_s\gamma'_s \dotprod{\nabla f(x_s),x^\ast - z_s} &= A_s \eta_s \dotprod{ \nabla f(x_s) ,x^\ast-x_s} + \bpar{A_s\eta_s - B_s \gamma'_s} \dotprod{ \nabla f(x_s) ,z_s-x_s}
%     \\
%     &\le- A_s \eta_s a \nu\norm{x_s-x^\ast}^2 + \bpar{A_s\eta_s - B_s \gamma'_s} \dotprod{ \nabla f(x_s) ,z_s-x_s},
% \end{align}
% where we used Assumptions~\ref{ass:uaac} and \ref{ass:eb}.
% So, we have
% \begin{align}
%     \dotprod{\nabla \varphi(x_s),b(x_s)} +  \mathbb{E}_{\xi} \varphi(\overline{x}_s+ G(\overline{x}_s)) - \varphi(\overline{x}_s) &\le \frac{1}{2}\bpar{ \frac{dB_s}{ds}-B_s \eta_s'}\norm{z_s-x^\ast}^2 +  \bpar{ \frac{L_0}{2}\frac{dA_s}{ds} + \frac{B_s \eta'_s}{2} - A_s \eta_s a \nu}\norm{x_s-x^\ast}^2\\
%     & + \bpar{A_s\eta_s - B_s \gamma'_s} \dotprod{ \nabla f(x_s) ,z_s-x_s}+\bpar{\rho B_s \gamma_s^{'2} + A_s \gamma_s( \rho\gamma_s L-2)}\frac{1}{2}\norm{\nabla f(x_s)}^2.
% \end{align}
% \item
% \paragraph{Parameter tuning}
% Inspired by the proof of Proposition~\ref{prop:nest_ode_triple}, we chose $A_s = e^{\theta s}$, $B_s = L_0e^{\theta_s}$, with $ \theta = \eta' = \frac{a}{\rho}\frac{\nu}{\sqrt{L_0}}\frac{1}{\sqrt{L}},\quad \gamma_s = \frac{1}{\rho L}\quad \gamma' = \frac{1}{\rho} \frac{1}{\sqrt{L_0 L}}, $ and $\eta = \frac{1}{\rho}\sqrt{\frac{L_0}{L}}$ \jh{donner minimum explication sur ça}. Such that
% \begin{align}
%     A_s\eta_s - B_s \gamma'_s = 0
% \end{align}
% \begin{align}
%     \rho B_s \gamma_s^{'2} + A_s \gamma_s( \rho\gamma_s L-2) = \frac{e^{-\theta_s}}{\rho L} - \frac{e^{-\theta_s}}{\rho L}=0, 
% \end{align}
% \begin{align}
%     \frac{dB_s}{ds}-B_s \eta_s' = \theta L_0 e^{-\theta_s} - L_0 \eta'e^{-\theta_s} = 0,
% \end{align}
% \begin{align}
%      \frac{L_0}{2}\frac{dA_s}{ds} + \frac{B_s \eta'_s}{2} - A_s \eta_s a \nu = \frac{\theta L_0}{2} + \frac{L_0\eta'}{2} - a\nu\sqrt{\frac{L_0}{L}} = 0.
% \end{align}
% This shows that
% \begin{equation}\label{eq:cont_supermart}
%     \phi(t) \le \phi(0) + M_t,
% \end{equation}
% and $\phi$ is a uniformly integrable supermartingale. We evaluate \eqref{eq:cont_supermart} at $t = T_k$, take expectation and use Theorem~\ref{thm:martingal_stopping} and Proposition~\ref{prop:discretization_constant_param}, to get
% \begin{equation}
% \E{e^{  \frac{a}{\rho}\frac{\nu}{\sqrt{L_0}}\frac{1}{\sqrt{L}}T_k}f(\tilde x_k)-f^\ast)} \le \phi(0).
% \end{equation}

% \section{Exact convergence rates}
%  \begin{proposition}\label{prop:exact}
%     Define the point-wise: Polyak-Lojasiewicz constant $\mu(x) :=  2\frac{\norm{\nabla f(x)}^2}{f(x)-f^\ast}$, aiming condition constant $a(x) := \frac{\dotprod{\nabla f(x),x-x^\ast}}{\norm{\nabla f(x)}\norm{x-x^\ast}}$ and quadratic growth constant $\mu_0(x) = 0.5\cdot\frac{f(x)-f^\ast}{\norm{x-x^\ast}^2}$.
    
%     (i) If $\xcont \sim \eqref{eq:gf}$, we have
%     $$f(x_t)-f^\ast = e^{-2\int_{0}^t \mu(x_s)ds}(f(x_0)-f^\ast).$$
%     (ii) If $\xcont \sim \eqref{eq:nmo}$ with $ \eta'_t = a(x_t)\sqrt{\mu(x_t)},\quad \eta_t = \frac{L_0}{\sqrt{\mu_0(x_t)}}, \quad \gamma'_t = \frac{1}{\sqrt{\mu_0(x_t)}},\quad \gamma_t \ge 0$, we have
%     $$f(x_t)-f^\ast \le e^{-\int_0^t a(x_s)\sqrt{\mu(x_s)}ds}L_0\norm{x_0-x^\ast}^2,$$
%     with $L_0 = \sup_{x\in \R^d} \mu_0(x)$. \jh{sup sur temps traj ?}
% \end{proposition}
% A straightforward consequence of Proposition~\ref{prop:exact} is that $f \in \pl \Rightarrow \int_{0}^t \mu(x_s)ds \ge \mu t$, and $f \in \pl \cap \ac \Rightarrow \int_0^t a(x_s)\sqrt{\mu(x_s)}ds \ge a\sqrt{\mu}t$, such that these rates are necessarily sharper.
% We observe that Proposition~\ref{prop:exact} \textit{(i)} provides an \textit{exact} decrease bound, fully determined by the initialization and the $\mu(\cdot)$ values along the trajectory. \jh{xp} Note that they are no other assumptions over $f$ except that of differentiability. In the case where the trajectory reach a local minima at time $t_0$, for which $\mu(x_{t_0}) = 0$, the bound indicates a linear decrease until this time, such that 
% $$f(x_t)-f^\ast = e^{-2\int_0^{\min(t,t_0)}\mu(x_s)ds}(f(x_0)-f^\ast).$$

%  %If it implies to recover the classic convergence rate in the first case, we deduce a different rate in the latter case.
%  Proposition~\ref{prop:exact} \textit{(ii)} indicates that the trajectory could cross regions such that $a(\cdot) \le 0$ (negative aiming condition), or $\mu(\cdot) = 0$ (saddle point of local minima), while still having $\int_0^t a(x_s)\sqrt{\mu(x_s)}ds \ge 0$. It may occur if the trajectory quickly fall in region with high aiming condition, or if momentum allows to quickly escape a saddle point/local minima. To sum up, all that matter is how the conditions are verified in average along the trajectory.
%  %However importantly, $ \int_0^t a(x_s)\sqrt{\mu(x_s)}ds$ could be significantly larger that $a\sqrt{\mu}$, the latter could even be reduced to zero.
% %Interestingly, bound (ii) allows the trajectory to cross a local minima $\hat{x}$, for which $\mu(\hat{x}) = 0$, while still having a linear convergence rate of convergence to the global minimum. All that matter is the average along the trajectory. 
% We further develop this analysis in Appendix~\ref{app:adptative_condition}, in the following directions:

% (1.) We extend Proposition~\ref{prop:exact} to the discrete setting. In the case of \eqref{eq:gf} and \eqref{eq:gd}, we also derive a characterization that makes use of the information of $a(x)$.

%  (2.) In the case of \eqref{eq:nmo}, the result of Proposition~\ref{prop:exact} is purely conceptual. Indeed, the tuning of the equations depends \textit{exactly} on untractable quantities. We thus propose a different viewpoint, with an assumption that holds on average along the trajectory. By instance in the continuous case, a relaxed aiming condition is that for some target convergence rate $\theta > 0$ and $a_{avg} > 0$, $\xcont$ is such that
%     \begin{align*}
%         \int_0^t e^{\theta s}\Delta_sds
%         \ge 0,
%     \end{align*}
%     $\Delta_s :=\dotprod{\nabla f(x_s),x_s-x^\ast}-a_{avg}\norm{\nabla f(x_s)}\norm{x_s-x^\ast}$.
%     Under these averaged assumptions, we recover the same rates as in Section~\ref{sec:cv_results}, with $a_{avg}$ replacing $a$.
% With this point of view as well, for $f\in \ac$ we may have $a_{avg} \gg a$.  These results could indicate why in practice we can select more aggressive parameters than what is allowed when evaluating geometrical constants. \jh{ça me semble intéressant de noter qu'on obtient des meilleurs résultats qu'en prenant l'inf des PL (par ex) le long de la trajectoire, ce qui est le truc déductible immédiatement}
% \subsubsection{Proof of Proposition~\ref{prop:exact}}\label{app:proof_exact}
% \textbf{(i) Case of gradient flow \eqref{eq:gf}. \:}
%     Set
%     $$E_t:= e^{2\int_{0}^t \mu(x_s)ds}(f(x_t)-f^\ast). $$
%     We have
%     \begin{equation}
%         \begin{aligned}\label{eq:pl_avg:1}
%                     \dot{E_t} &= 2\mu(x_t)e^{2\int_{0}^t \mu(x_s)ds}(f(x_t)-f^\ast) + e^{2\int_{0}^t \mu(x_s)ds}\dotprod{\nabla f(x_t),\dot x_t}\\
%         &\overset{\eqref{eq:gf}}{=} 2\mu(x_t)e^{2\int_{0}^t \mu(x_s)ds}(f(x_t)-f^\ast) - e^{2\int_{0}^t \mu(x_s)ds}\norm{\nabla f(x_t)}^2
%         \end{aligned}
%     \end{equation}

%     Recall $\mu(x) := 2\frac{\norm{\nabla f(x)}^2}{f(x)-f^\ast}$. Plugging it in \eqref{eq:pl_avg:1}, we deduce
%     \begin{align}
%         \dot{E_t} = 2\mu(x_t)e^{2\int_{0}^t \mu(x_s)ds}(f(x_t)-f^\ast) - 2\mu(x_t)e^{2\int_{0}^t \mu(x_s)ds}(f(x_t)-f^\ast)=0.
%     \end{align}
%     By integrating we deduce 
%     \begin{equation}
%         f(x_t)-f^\ast = e^{-2\int_{0}^t \mu(x_s)ds}(f(x_0)-f^\ast).
%     \end{equation}

% \noindent
% \textbf{(ii) Case of Nesterov Momentum ODE \eqref{eq:nmo}}
%      Let be
% \begin{equation}
%     E_t = e^{\int_0^t a(x_s)\sqrt{\mu(x_s)}ds}(f(x_t)-f^\ast) + \frac{ \beta}{2}e^{\int_0^t a(x_s)\sqrt{\mu(x_s)}ds}\norm{z_t -  x^\ast}^2.
% \end{equation}
% Denote $\Delta_t := e^{\int_0^t a(x_s)\sqrt{\mu(x_s)}ds}$ and $u(t) := a(x_t)\sqrt{\mu_t}$. We have

% \begin{align}
%     \Delta_t^{-1}\dot{E}_t &=  u_t(f(x_t)-f^\ast)  + \frac{ u_t \beta}{2}\norm{z_t -  x^\ast}^2\nonumber\\
%     &+ \dotprod{\nabla f(x_t),\dot{x}_t} + \beta \dotprod{z_t -  x^\ast,\dot{z}_t}
%    \nonumber \\&= u_t(f(x_t)-f^\ast)  + \frac{u_t \beta}{2}\norm{z_t -  x^\ast}^2\nonumber\\
%     &+ \dotprod{\nabla f(x_t), \eta_t(z_t-x_t) - \gamma_t \nabla f(x_t)} +\beta  \dotprod{z_t - x^\ast,\eta'_t(x_t-z_t) - \gamma'_t \nabla f(x_t) }\nonumber\\
%     &=u_t (f(x_t)-f^\ast)  +\frac{u_t \beta}{2}\norm{z_t -  x^\ast}^2\nonumber\\
%     &+ u_t \dotprod{\nabla f(x_t), x^\ast - x_t} + (\eta_t -\beta \gamma'_t)\dotprod{\nabla f(x_t),z_t - x^\ast}\nonumber\\
%     &+ \beta  \eta' \dotprod{z_t - x^\ast,x_t - z_t} - \gamma_t \norm{\nabla f(x_t)}^2\nonumber\\
%     &= u_t (f(x_t)-f^\ast)  +\frac{ \beta }{2} \bpar{u_t  - \eta'_t}\norm{z_t -  x^\ast}^2\nonumber\\
%     &+ \eta_t \dotprod{\nabla f(x_t), x^\ast - x_t} + (\eta_t -\beta_t \gamma'_t)\dotprod{\nabla f(x_t),z_t - x^\ast}\nonumber\\
%     &+ \frac{\beta  \eta'_t}{2}\norm{x^\ast-x_t}^2 - \gamma_t \norm{\nabla f(x_t)}^2.\label{eq:nest_ode_0}
% \end{align}
% We chose $\eta'_t = u_t$ to cancel the $\norm{z_t-x^\ast}^2$ term. We choose $\eta_t = \beta\gamma'_t$ to cancel the $\dotprod{\nabla f(x_t),z_t - x^\ast}$ term. Then,
% \begin{align}\label{eq:exact:nest_ode_1}
%    \Delta_t^{-1}\dot{E}_t = u_t(f(x_t)-f^\ast) +\beta \gamma'_t \dotprod{\nabla f(x_t), x^\ast - x_t}+  \frac{\beta  u_t }{2}\norm{x^\ast-x_t}^2 - \gamma_t \norm{\nabla f(x_t)}^2.
% \end{align}
% By Lemma~\ref{lem:tec:exact_scalar_prod_bound}, we have
% $$
%     \dotprod{\nabla f(x_t), x^\ast - x_t} =  -a(x_t) \sqrt{\frac{\mu(x_t)}{\mu_0(x_t)}}(f(x_t)-f^\ast) - a(x_t) \sqrt{\mu(x_t) \mu_0(x_t)}\frac{1}{2}\norm{x_t-x^\ast}^2,
% $$

% Such that fixing $\gamma'_t := \frac{1}{\sqrt{\mu_0(x_t)}}$ and ignoring the $- \gamma_t \norm{\nabla f(x_t)}^2$ term, \eqref{eq:exact:nest_ode_1} becomes

% \begin{equation}
% \begin{aligned}\label{eq:exact:nest_ode_2}
%     \Delta_t^{-1}\dot{E}_t &= u_t\bpar{1-\frac{\beta}{\mu_0(x_t)}}(f(x_t)-f^\ast).
% \end{aligned}
% \end{equation}
% It can be ensured choosing $\beta = \sup_{x\in \R^d} \mu_0(x) = L_0$. Then, under the set of parameters
% $$ \eta'_t = a(x_t)\sqrt{\mu(x_t)},\quad \eta_t = \frac{L_0}{\sqrt{\mu_0(x_t)}}, \quad \gamma'_t = \frac{1}{\sqrt{\mu_0(x_t)}},\quad \gamma_t \ge 0,$$
% we deduce $\dot E_t \le 0$ for all $t\in \R_+$. By integrating, we deduce
% \begin{equation}
%     E_t \le E_0 \Rightarrow f(x_t)-f^\ast \le e^{-\int_0^t a(x_s)\sqrt{\mu(x_s)}ds}E_0 \le e^{-\int_0^t a(x_s)\sqrt{\mu(x_s)}ds}L_0\norm{x_0-x^\ast}^2.
% \end{equation}
% The last equality uses $f(x_0)-f^\ast \le \frac{L_0}{2}\norm{x_0-x^\ast}^2$.
 
% \subsubsection{Proof of Proposition~\ref{prop:exact_pl_ac_gf}}

%   Set
%     $$E_t:= e^{\int_{0}^t a(x_s)\mu(x_s)ds}(f(x_t)-f^\ast) +e^{\int_{0}^t a(x_s)\mu(x_s)ds}(f(x_t)-f^\ast) \frac{\sqrt{\mu_0 L_0}}{2}\norm{x_t-x^\ast}^2. $$
% Noting $\delta_t := e^{\int_{0}^t a(x_s)\mu(x_s)ds}(f(x_t)-f^\ast)$, we have
% \begin{equation}
%     \begin{aligned}\label{eq:exact:pl_avg:1}
%       \Delta_t^{-1}  \dot E_t &:=  a(x_t)\mu(x_t) (f(x_t)-f^\ast) + \dotprod{\nabla f(x_t), \dot x_t} + \frac{a(x_t)\mu(x_t) \sqrt{\mu_0 L_0}}{2}\norm{x_t - x^\ast}^2 + \sqrt{\mu_0 L_0}\dotprod{x_t-x^\ast,\dot x_t}\\
%       &\overset{\eqref{eq:gf}}{=} a(x_t)\mu(x_t) (f(x_t)-f^\ast) + \frac{a(x_t)\mu(x_t) \sqrt{\mu_0 L_0}}{2}\norm{x_t - x^\ast}^2 - \norm{\nabla f(x_t)}^2 - \sqrt{\mu_0 L_0}\dotprod{x_t - x^\ast,\nabla f(x_t)}. 
%     \end{aligned}
% \end{equation}
%     Applying Lemma~\ref{lem:scalar_prod_bound}, namely
% $$ -\dotprod{x_t - x^\ast,\nabla f(x_t)} \le  -a \sqrt{\frac{\mu}{L_0}}(f(x_t)-f^\ast) - a \sqrt{\mu \mu_0}\frac{1}{2}\norm{x-x^\ast}^2,$$
% we then have that, ignoring the$ -\norm{\nabla f(x_t)}^2 $ term, \eqref{eq:exact:pl_avg:1} becomes
% $$    \Delta_t^{-1}  \dot E_t  \le \bpar{a(x_t)\mu(x_t) - }(f(x_t)-f^\ast)  + \bpar{}\frac{1}{2}\norm{x-x^\ast}^2 $$
    
%     \subsubsection{Proof of Proposition~\ref{prop:ode_nest_rsi_traj_depend}}
%          Let be
% \begin{equation}
%     E_t = e^{\int_0^t \frac{\nu(x_s)}{\sqrt{L_0(x_s)}}ds}(f(x_t)-f^\ast) + \frac{ \beta}{2}e^{\int_0^t \frac{\nu(x_s)}{\sqrt{L_0(x_s)}}ds}\norm{z_t -  x^\ast}^2.
% \end{equation}
% Denote $\Delta_t := e^{\int_0^t \frac{\nu(x_s)}{\sqrt{L_0(x_s)}}ds}$ and $u_t := \frac{\nu(x_t)}{\sqrt{L_0(x_t)}}$. We have

% \begin{align}
%     \Delta_t^{-1}\dot{E}_t &=  u_t(f(x_t)-f^\ast)  + \frac{ u_t \beta}{2}\norm{z_t -  x^\ast}^2\nonumber\\
%     &+ \dotprod{\nabla f(x_t),\dot{x}_t} + \beta \dotprod{z_t -  x^\ast,\dot{z}_t}
%    \nonumber \\&= u_t(f(x_t)-f^\ast)  + \frac{u_t \beta}{2}\norm{z_t -  x^\ast}^2\nonumber\\
%     &+ \dotprod{\nabla f(x_t), \eta_t(z_t-x_t) - \gamma_t \nabla f(x_t)} +\beta  \dotprod{z_t - x^\ast,\eta'_t(x_t-z_t) - \gamma'_t \nabla f(x_t) }\nonumber\\
%     &=u_t (f(x_t)-f^\ast)  +\frac{u_t \beta}{2}\norm{z_t -  x^\ast}^2\nonumber\\
%     &+ u_t \dotprod{\nabla f(x_t), x^\ast - x_t} + (\eta_t -\beta \gamma'_t)\dotprod{\nabla f(x_t),z_t - x^\ast}\nonumber\\
%     &+ \beta  \eta' \dotprod{z_t - x^\ast,x_t - z_t} - \gamma_t \norm{\nabla f(x_t)}^2\nonumber\\
%     &= u_t (f(x_t)-f^\ast)  +\frac{ \beta }{2} \bpar{u_t  - \eta'_t}\norm{z_t -  x^\ast}^2\nonumber\\
%     &+ \eta_t \dotprod{\nabla f(x_t), x^\ast - x_t} + (\eta_t -\beta_t \gamma'_t)\dotprod{\nabla f(x_t),z_t - x^\ast}\nonumber\\
%     &+ \frac{\beta  \eta'_t}{2}\norm{x^\ast-x_t}^2 - \gamma_t \norm{\nabla f(x_t)}^2.\label{eq:nest_ode_0}
% \end{align}
% We chose $\eta'_t = \theta_t$ to cancel the $\norm{z_t-x^\ast}^2$ term. We choose $\eta_t = \beta\gamma'_t$ to cancel the $\dotprod{\nabla f(x_t),z_t - x^\ast}$ term. Then,
% \begin{align}\label{eq:nest_ode_1}
%    \Delta_t^{-1}\dot{E}_t = u_t(f(x_t)-f^\ast) +\beta \gamma'_t \dotprod{\nabla f(x_t), x^\ast - x_t}+  \frac{\beta  u_t }{2}\norm{x^\ast-x_t}^2 - \gamma_t \norm{\nabla f(x_t)}^2.
% \end{align}


% Using $\nu(x):= \frac{\dotprod{\nabla f(x),x-x^\ast}}{\norm{x-x^\ast}^2}$, fixing $\gamma'_t := \frac{1}{\sqrt{L_0(x_t)}}$ and ignoring the $- \gamma_t \norm{\nabla f(x_t)}^2$ term, \eqref{eq:nest_ode_1} becomes

% \begin{equation}
% \begin{aligned}\label{eq:nest_ode_2}
%     \Delta_t^{-1}\dot{E}_t &= u_t(f(x_t)-f^\ast) -\beta \frac{\nu(x_t)}{\sqrt{L_0(x_t)}}\norm{x_t-x^\ast}^2+  \frac{\beta  u_t }{2}\norm{x^\ast-x_t}^2 \\
%     &=u_t(f(x_t)-f^\ast) -\beta u_t\norm{x_t-x^\ast}^2+  \frac{\beta  u_t }{2}\norm{x^\ast-x_t}^2 \\
%     &=u_t\bpar{f(x_t)-f^\ast - \frac{\beta}{2}\norm{x_t-x^\ast}^2}\\
%     &=u_t\bpar{L_0(x_t)-\beta}\norm{x_t-x^\ast}^2,
% \end{aligned}
% \end{equation}
% with the last line using $ L_0(x) := 2\frac{f(x)-f^\ast}{\norm{x-x^\ast}^2}$. Choosing $\beta = \sup_{x\in \R^d} L_0(x) := L_0$, we deduce $\dot E_t \le 0$ for all $t\in \R_+$. By integrating, we deduce
% \begin{equation}
%     E_t \le E_0 \Rightarrow f(x_t)-f^\ast \le e^{-\int_0^t \frac{\nu(x_s)}{\sqrt{L_0(x_s)}}ds}E_0 \le e^{-\int_0^t \frac{\nu(x_s)}{\sqrt{L_0(x_s)}}ds}\overline{L_0}\norm{x_0-x^\ast}^2.
% \end{equation}
% The last equality uses $f(x_0)-f^\ast \le \frac{L_0}{2}\norm{x_0-x^\ast}^2$.

\section{Proof of Proposition~\ref{lem:toy_example_ac}}
\textbf{Computation of $\nabla h(x)$.} We have
\begin{equation}
    \begin{aligned}
        \nabla h(x) &= \frac{x}{\norm{x}}f'(\norm{x})g\left(\frac{x}{\lVert x \rVert} \right) + \bpar{\frac{1}{\norm{x}}I_d - \frac{x x^T}{\norm{x}^3}}f(\lVert x \rVert)\nabla g\left(\frac{x}{\lVert x \rVert} \right)\\
        &=\frac{x}{\norm{x}}f'(\norm{x})g\left(\frac{x}{\lVert x \rVert} \right) + \frac{1}{\norm{x}}\bpar{I_d -\frac{x x^T}{\norm{x}^2}}f(\lVert x \rVert)\nabla g\left(\frac{x}{\lVert x \rVert} \right)
    \end{aligned}
\end{equation}
Noting $r := \norm{x}$ and $u :=x/\norm{x}$, this can be rewritten
\begin{equation}
    \begin{aligned}
        \nabla h(x) = \nabla  h(ru) = f'(r)g(u)u + \frac{1}{r}f(r) \nabla_{\sphere} g(u),
    \end{aligned}
\end{equation}
where $\nabla_{\sphere} g(u) = (I_d-uu^T)\nabla g(u)$ is the Riemannian gradient of $g$ on the sphere. $\nabla_{\sphere} g(u)$ belongs to the tangent space at $u$, namely $\nabla_{\sphere} g(u) \in  \mathcal{T}_u = \{ v, \dotprod{u,v} = 0 \}$. In particular, it implies that $u$ and $\nabla_{\sphere} g(u)$ are orthogonal, so
\[\norm{\nabla h(x)}^2 =f'(r)^ 2g(u)^2 + \frac{1}{r^2}f(r)^2 \norm{\nabla_{\sphere} g(u)}^2, \]
and 
\[\dotprod{u,\nabla_{\sphere} g(u)} = 0.\]
So, we have
\[\dotprod{\nabla h(x),x-x^\ast} = \dotprod{f'(r)g(u)u + \frac{1}{r}f(r) \nabla_{\sphere} g(u),ru} = rf'(r)g(u),\]
 and 
 \[\norm{\nabla h(x)}\norm{x-x^\ast} = \sqrt{f'(r)^ 2g(u)^2 + \frac{1}{r^2}f(r)^2 \norm{\nabla_{\sphere} g(u)}^2}r, \]
 such that
 \[\frac{\dotprod{\nabla h(x),x-x^\ast} }{\norm{\nabla h(x)}\norm{x-x^\ast}} = \frac{f'(r)g(u)}{\sqrt{f'(r)^ 2g(u)^2 + \frac{1}{r^2}f(r)^2 \norm{\nabla_{\sphere} g(u)}^2}}.\]
 So, $h \in \ac$ for $a$ such that
 \[a := \inf_{r>0,u \in \mathbb{S}^{d-1}} \frac{f'(r)g(u)}{\sqrt{f'(r)^ 2g(u)^2 + \frac{1}{r^2}f(r)^2 \norm{\nabla_{\sphere} g(u)}^2}} = \inf_{r>0,u \in \mathbb{S}^{d-1}} \frac{1}{\sqrt{1+ \bpar{\frac{f(r)}{rf'(r)}}^2 \bpar{\frac{\norm{\nabla_{\sphere} g(u)}}{g(u)}}^2}}  \]
 % If $f$ is star convex, we have
 % \[f'(r)r \ge f(r) \Rightarrow 1\ge \frac{f(r)}{f'(r)r},\]
 % which in turn implies 
 % \[ a \ge  \inf_{r>0,u \in \mathbb{S}^{d-1}} \frac{1}{\sqrt{1+ \bpar{\frac{\norm{\nabla_{\sphere} g(u)}}{g(u)}}^2}}.\]
 % If $g\ge 1$, we deduce
 % \[a \ge  \inf_{r>0,u \in \mathbb{S}^{d-1}} \frac{1}{\sqrt{1+ \norm{\nabla_{\sphere} g(u)}^2}}.\]
 % $a$ is smaller as the spherical variations become larger. In particular, if $g$ does not move spherically, \textit{e.g.} $g(x) = \norm{x}^2$ or $g$ is constant, this implies $\nabla_{\sphere} g(u)$. In this case, $a = 1$, and $h$ is a radial function.
 \subsection{Proof of Proposition~\ref{prop:converse_link}}\label{app:prop:converse}
 For the converse relation, we define a weaker relaxation of strong convexity than $\sqc$, defined by the Restricted Secant Inequality (RSI) \cite{zhang2013gradient}.
\begin{definition}
   For $\nu>0$, we call $\rsi$ the set of functions $f$ that satisfies the Restricted Secant Inequality, namely $\forall x \in \R^d,~ \dotprod{\nabla f(x),x-x^\ast} \ge \nu \norm{x-x^\ast}^2.$
\end{definition}
A function in $\sqc$ is also in $\rsif{\frac{2\mu}{2-\gamma}}$, such that the RSI condition is weaker \cite{hermant2024study}. 
\begin{proposition}
    Let $f\in\rsi$. If $f$ also satisfies
    \begin{equation}\label{eq:eb+}
         \forall x \in \R^d,~ \norm{\nabla f(x)}\le L\norm{x-x^\ast},
    \end{equation}

    then $f \in \acf{\frac{\nu}{L}}$.
\end{proposition}
\begin{proof}
    We have \[\dotprod{\nabla f(x),x-x^\ast} \ge \nu \norm{x-x^\ast}^2 \ge \frac{\nu}{L}\norm{x-x^\ast}\norm{\nabla f(x)}.\]
\end{proof}
The condition \eqref{eq:eb+} is sometimes named $\text{EB}^+$, as the upper version of the Error-Bound condition \cite{guille2021study}. It can also be interpreted as being Lipschitz with respect to the minimizers. The result follows from the previous proposition.
%%%%%%%%%%%%%%%%%%%%%%%%%%%%%%%%%%%%%%%%%%%%%%%%%%%%%%%%%%%%%%%%%%%%%%%%%%%%%%%
%%%%%%%%%%%%%%%%%%%%%%%%%%%%%%%%%%%%%%%%%%%%%%%%%%%%%%%%%%%%%%%%%%%%%%%%%%%%%%%


\end{document}


% This document was modified from the file originally made available by
% Pat Langley and Andrea Danyluk for ICML-2K. This version was created
% by Iain Murray in 2018, and modified by Alexandre Bouchard in
% 2019 and 2021 and by Csaba Szepesvari, Gang Niu and Sivan Sabato in 2022.
% Modified again in 2023 and 2024 by Sivan Sabato and Jonathan Scarlett.
% Previous contributors include Dan Roy, Lise Getoor and Tobias
% Scheffer, which was slightly modified from the 2010 version by
% Thorsten Joachims & Johannes Fuernkranz, slightly modified from the
% 2009 version by Kiri Wagstaff and Sam Roweis's 2008 version, which is
% slightly modified from Prasad Tadepalli's 2007 version which is a
% lightly changed version of the previous year's version by Andrew
% Moore, which was in turn edited from those of Kristian Kersting and
% Codrina Lauth. Alex Smola contributed to the algorithmic style files.
